% TOPLAS submission manuscript (acmart).  Converted on 2026-07-30 from the
% arXiv source (arXiv:2607.23500), which was generated from
% papers/POPL27/paper_draft.tex.  This file is edited by hand from now on.
%
% Per the TOPLAS author guidelines: acmsmall, no "review" option, no manual
% line numbers (the submission system adds them).
\documentclass[acmsmall,screen=true,acmthm=false]{acmart}
% acmthm=false keeps the paper's own amsthm+aliascnt theorem setup below, so
% that cleveref keeps reporting the right environment names (e.g., Lemma
% rather than Theorem) while all environments share one counter.

% Fresh-build workaround: latexmk seeds a first run with a dummy .aux that
% claims the previous run had one page (\gdef\@abspage@last{1}).  The kernel
% then fires its shipout/lastpage hook already on page 1, where hyperxmp
% (loaded by acmart) re-executes the author list to collect XMP metadata;
% acmart's affiliation city/country check misfires in that harvesting
% context and aborts with "No country present for an affiliation".
% Restoring the kernel default ("last page unknown") makes such a run behave
% like an ordinary first run; real page data from a genuine previous run is
% left untouched.
\makeatletter
\AtBeginDocument{\ifdefined\@abspage@last
  \ifnum\@abspage@last=1 \gdef\@abspage@last{2147483647}\fi\fi}
\makeatother

%% ---- ACM metadata (placeholders until acceptance) -----------------------
\setcopyright{none} % TODO(final): CC license under ACM Open Access
\acmJournal{TOPLAS}
\acmYear{2026}

%% amsmath, AMS symbols (via newtxmath), xcolor, and hyperref are loaded by
%% acmart itself; fullpage and authblk from the arXiv version do not apply.
\usepackage{mathtools}
\usepackage{amsthm}
\usepackage{stmaryrd}
\usepackage{listings}
\usepackage{booktabs}
\usepackage{microtype}
\usepackage{aliascnt}
\usepackage{tikz}

\newtheorem{theorem}{Theorem}[section]

\newaliascnt{lemma}{theorem}
\newtheorem{lemma}[lemma]{Lemma}
\aliascntresetthe{lemma}

\newaliascnt{corollary}{theorem}
\newtheorem{corollary}[corollary]{Corollary}
\aliascntresetthe{corollary}

\newaliascnt{definition}{theorem}
\newtheorem{definition}[definition]{Definition}
\aliascntresetthe{definition}

\newaliascnt{example}{theorem}
\newtheorem{example}[example]{Example}
\aliascntresetthe{example}

\newaliascnt{remark}{theorem}
\newtheorem{remark}[remark]{Remark}
\aliascntresetthe{remark}

\usepackage[nameinlink]{cleveref}
\crefname{theorem}{Theorem}{Theorems}
\Crefname{theorem}{Theorem}{Theorems}
\crefname{lemma}{Lemma}{Lemmas}
\Crefname{lemma}{Lemma}{Lemmas}
\crefname{corollary}{Corollary}{Corollaries}
\Crefname{corollary}{Corollary}{Corollaries}
\crefname{definition}{Definition}{Definitions}
\Crefname{definition}{Definition}{Definitions}
\crefname{example}{Example}{Examples}
\Crefname{example}{Example}{Examples}
\crefname{remark}{Remark}{Remarks}
\Crefname{remark}{Remark}{Remarks}

\newcommand{\R}{\mathbb{R}}
\newcommand{\Q}{\mathbb{Q}}
\newcommand{\N}{\mathbb{N}}
\newcommand{\Fin}[1]{\mathsf{Fin}(#1)}
\newcommand{\Flag}[1]{\mathcal{F}^{#1}}
\newcommand{\FlagAlg}[1]{\mathcal{A}^{#1}}
\newcommand{\poshom}[1]{\mathrm{Hom}^+(\mathcal{A}^{#1}, \mathbb{R})}
\newcommand{\den}[2]{p(#1;\,#2)}
\newcommand{\tdensity}[2]{\pi(#1;\,#2)}
% Retain acmart's semantic paragraph-heading structure and spacing, but make
% the run-in heading text more visible by setting it in upright bold.
\newcommand{\boldparagraph}[1]{\noindentparagraph{\textnormal{\textbf{#1}}}}
% Allow line breaks after underscores in long Lean identifiers; without this
% the narrower acmsmall measure produces overfull lines.  Robust so that
% \lean can appear inside captions and other moving arguments.
\DeclareRobustCommand{\lean}[1]{{\def\_{\textunderscore\allowbreak}\texttt{#1}}}
% Long unbreakable identifiers otherwise force overfull lines in the
% acmsmall measure; let TeX stretch interword spaces as a last resort.
\setlength{\emergencystretch}{3em}
\newcommand{\emptyt}{\emptyset_{\mathrm{t}}}
\newcommand{\leansmul}{\mathbin{\vcenter{\hbox{$\scriptstyle\bullet$}}}}
\newcommand{\Ext}{\operatorname{Ext}}

\tikzset{
  graph vertex/.style={circle, draw, fill=white, inner sep=0pt, minimum size=4.5pt},
  graph root/.style={circle, draw, fill=black, inner sep=0pt, minimum size=4.5pt},
  graph edge/.style={line width=0.45pt},
  graph nonedge/.style={line width=0.45pt, dashed},
  inline/.style={scale=0.20, baseline=-0.5ex},
  fv/.style={circle, fill=white, draw, inner sep=0pt, minimum size=3pt},
  fr/.style={circle, fill=black, inner sep=0pt, minimum size=3pt}
}

\newcommand{\fKtwobull}{%
  \begin{tikzpicture}[inline]
    \draw (-1,0) -- (1,0);
    \node[fr] at (-1,0) {};
    \node[fv] at (1,0) {};
  \end{tikzpicture}%
}
\newcommand{\fbarKtwobull}{%
  \begin{tikzpicture}[inline]
    \node[fr] at (-1,0) {};
    \node[fv] at (1,0) {};
  \end{tikzpicture}%
}
\newcommand{\fbarKthree}{%
  \begin{tikzpicture}[inline]
    \node[fv] at (90:1) {};
    \node[fv] at (210:1) {};
    \node[fv] at (330:1) {};
  \end{tikzpicture}%
}
\newcommand{\fbarPthree}{%
  \begin{tikzpicture}[inline]
    \draw (210:1) -- (330:1);
    \node[fv] at (90:1) {};
    \node[fv] at (210:1) {};
    \node[fv] at (330:1) {};
  \end{tikzpicture}%
}
\newcommand{\fPthree}{%
  \begin{tikzpicture}[inline]
    \draw (210:1) -- (90:1) -- (330:1);
    \node[fv] at (90:1) {};
    \node[fv] at (210:1) {};
    \node[fv] at (330:1) {};
  \end{tikzpicture}%
}
\newcommand{\fKthree}{%
  \begin{tikzpicture}[inline]
    \draw (210:1) -- (90:1) -- (330:1) -- (210:1);
    \node[fv] at (90:1) {};
    \node[fv] at (210:1) {};
    \node[fv] at (330:1) {};
  \end{tikzpicture}%
}
\newcommand{\fbarKthreebull}{%
  \begin{tikzpicture}[inline]
    \node[fr] at (90:1) {};
    \node[fv] at (210:1) {};
    \node[fv] at (330:1) {};
  \end{tikzpicture}%
}
\newcommand{\fEthreebull}{%
  \begin{tikzpicture}[inline]
    \draw (210:1) -- (330:1);
    \node[fr] at (90:1) {};
    \node[fv] at (210:1) {};
    \node[fv] at (330:1) {};
  \end{tikzpicture}%
}
\newcommand{\fbarPthreebull}{%
  \begin{tikzpicture}[inline]
    \draw (90:1) -- (330:1);
    \node[fr] at (90:1) {};
    \node[fv] at (210:1) {};
    \node[fv] at (330:1) {};
  \end{tikzpicture}%
}
\newcommand{\fPthreebull}{%
  \begin{tikzpicture}[inline]
    \draw (90:1) -- (330:1) -- (210:1);
    \node[fr] at (90:1) {};
    \node[fv] at (210:1) {};
    \node[fv] at (330:1) {};
  \end{tikzpicture}%
}
\newcommand{\fKonetwobull}{%
  \begin{tikzpicture}[inline]
    \draw (90:1) -- (210:1);
    \draw (90:1) -- (330:1);
    \node[fr] at (90:1) {};
    \node[fv] at (210:1) {};
    \node[fv] at (330:1) {};
  \end{tikzpicture}%
}
\newcommand{\fCfive}{%
  \begin{tikzpicture}[inline]
    \draw (90:1) -- (18:1) -- (-54:1) -- (-126:1) -- (162:1) -- cycle;
    \node[fv] at (90:1) {};
    \node[fv] at (18:1) {};
    \node[fv] at (-54:1) {};
    \node[fv] at (-126:1) {};
    \node[fv] at (162:1) {};
  \end{tikzpicture}%
}
\newcommand{\fCfivebull}{%
  \begin{tikzpicture}[inline]
    \draw (90:1) -- (18:1) -- (-54:1) -- (-126:1) -- (162:1) -- cycle;
    \node[fr] at (90:1) {};
    \node[fv] at (18:1) {};
    \node[fv] at (-54:1) {};
    \node[fv] at (-126:1) {};
    \node[fv] at (162:1) {};
  \end{tikzpicture}%
}

\lstdefinelanguage{Lean4}{
  keywords={def,theorem,lemma,instance,structure,class,import,open,namespace,
             end,where,fun,let,have,show,by,match,with,if,then,else,
             return,do,for,in,noncomputable,abbrev,variable,section,
             native_decide,decide,norm_num,simp,ring,linarith,omega,
             exact,apply,rw,calc,intro,constructor,use,refine,obtain,
             rcases,cases,induction,assumption,contradiction,trivial,rfl,
             macro,elab,notation,example,
             infix,infixl,infixr,prefix,postfix},
  keywordstyle=\color{blue!70!black}\bfseries,
  comment=[l]{--},
  morecomment=[l]{/--},
  commentstyle=\color{gray}\itshape,
  stringstyle=\color{orange!80!black},
  morestring=[b]",
  showstringspaces=false,
  sensitive=true,
  basicstyle=\ttfamily\small,
  breaklines=true,
  columns=fullflexible,
  escapeinside={(*@}{@*)},
  literate=
    {->}{{$\to$}}2
    {<-}{{$\leftarrow$}}2
    {<=}{{$\leq$}}1
    {>=}{{$\geq$}}1
    {alpha}{{$\alpha$}}1
    {forall}{{$\forall$}}1
    {exists}{{$\exists$}}1
}
\lstset{language=Lean4, frame=single, framesep=4pt,
        rulecolor=\color{black!20}, backgroundcolor=\color{gray!5}}

\title[A Certificate-to-Proof Compiler for Flag Algebras]{A Certificate-to-Proof Compiler for Flag Algebras in Lean~4}

%% TODO(submission): add \orcid{...} for every author.  ACM requires an ORCID
%% for the submitting author at submission time and for all authors before
%% production.

\author{Gyeongwon Jeong}
\authornote{These authors are joint first authors of this work.}
\affiliation{%
  \department{School of Computing}
  \institution{KAIST}
  \city{Daejeon}
  \country{Korea}}
\affiliation{%
  \department{School of Computational Sciences}
  \institution{Korea Institute for Advanced Study (KIAS)}
  \city{Seoul}
  \country{Korea}}
\email{jgyw0910@kaist.ac.kr}

\author{Seonghun Park}
\authornotemark[1]
\affiliation{%
  \department{School of Computing}
  \institution{KAIST}
  \city{Daejeon}
  \country{Korea}}
\affiliation{%
  \department{School of Computational Sciences}
  \institution{Korea Institute for Advanced Study (KIAS)}
  \city{Seoul}
  \country{Korea}}
\email{hun57@kaist.ac.kr}

\author{Jihoon Hyun}
\affiliation{%
  \department{School of Computing}
  \institution{KAIST}
  \city{Daejeon}
  \country{Korea}}
\email{qawbecrdtey@kaist.ac.kr}

\author{Sang-il Oum}
\affiliation{%
  \department{Discrete Mathematics Group}
  \institution{Institute for Basic Science (IBS)}
  \city{Daejeon}
  \country{Korea}}
\affiliation{%
  \department{Department of Mathematical Sciences}
  \institution{KAIST}
  \city{Daejeon}
  \country{Korea}}
\email{sagil@ibs.re.kr}

\author{Hongseok Yang}
\affiliation{%
  \department{School of Computational Sciences}
  \institution{Korea Institute for Advanced Study (KIAS)}
  \city{Seoul}
  \country{Korea}}
\email{hongseokyang@kias.re.kr}

%% Condensed contact block for the submission (one line per author).  ACM
%% requires authors' addresses in journal articles (acmart guide,
%% \authorsaddresses), so the block is condensed rather than suppressed;
%% TAPS regenerates the full block from metadata at production.
\authorsaddresses{%
  Authors' Contact Information:
  Gyeongwon Jeong, KAIST, Daejeon, Republic of Korea,
  jgyw0910@kaist.ac.kr;
  Seonghun Park, KAIST, Daejeon, Republic of Korea, hun57@kaist.ac.kr;
  Jihoon Hyun, KAIST, Daejeon, Republic of Korea,
  qawbecrdtey@kaist.ac.kr;
  Sang-il Oum, Institute for Basic Science (IBS), Daejeon, Republic of
  Korea, sagil@ibs.re.kr;
  Hongseok Yang, Korea Institute for Advanced Study (KIAS), Seoul,
  Republic of Korea, hongseokyang@kias.re.kr.}

% acmart sets the float separations at \begin{document}, so cap the stretch
% afterwards.  Keeping acmart's designed 15pt but shrinking "plus 8pt" stops a
% page that needs vertical filling from dumping all of its slack into the gap
% under a top float.
\AtBeginDocument{\setlength{\textfloatsep}{15pt plus 2pt minus 4pt}}

\begin{document}

\begin{abstract}
Razborov's flag algebra method proves asymptotic bounds in extremal graph
theory, such as upper bounds on how often a fixed graph pattern can occur in
large graphs that avoid other specified patterns.  It turns a claim about
infinitely many graphs into a finite algebraic verification.  In practice, a
semidefinite-programming solver searches for a certificate consisting of
matrices with rational entries and rational coefficients.  If these data
satisfy certain algebraic identities and positivity conditions, they establish
the desired bound.  The method has played a central role in proofs of major
results in extremal graph theory.  The certificates for some of these results,
however, contain thousands of exact rational values and are generated through
large toolchains that rely on unverified numerical computation.  They are
therefore too large to audit by hand and cannot serve as trusted proof
artifacts without independent checking.  We present a Lean~4 formalization of
flag algebras for finite simple graphs, together with a certificate-to-proof
compiler built on the formalization.  Given such a certificate, the compiler
treats it as untrusted input, independently verifies its algebraic identities
and positivity conditions against the formal definitions, and constructs a
complete, machine-checked Lean proof.

The development has three layers: a specification layer that formalizes the
mathematics of the method, a reflection layer that provides verified
computation and generates thousands of verified identities required by a
certificate, and an automation layer that drives the compilation.  The compiler
runs inside Lean during elaboration: the theorem statement is fixed in the
source file, the certificate serves only as untrusted input, and Lean recomputes
every fact claimed by the certificate against the mathematical definitions.
We evaluate the compiler on seven certificates for Tur\'an-type upper bounds
and explain how Lean checks each proof obligation in the generated proofs.  Two
of the compiled upper bounds, for Mantel's theorem and the Erd\H{o}s pentagon
theorem, are paired with matching lower-bound constructions proved directly in
Lean, thereby establishing the exact value in both cases.  In particular, this
gives, to our knowledge, the first machine-checked proof of the Erd\H{o}s
pentagon theorem, whose known proofs all use flag algebras.
\end{abstract}

\begin{CCSXML}
<ccs2012>
<concept>
<concept_id>10003752.10003790.10003792</concept_id>
<concept_desc>Theory of computation~Proof theory</concept_desc>
<concept_significance>500</concept_significance>
</concept>
<concept>
<concept_id>10003752.10003790.10002990</concept_id>
<concept_desc>Theory of computation~Logic and verification</concept_desc>
<concept_significance>300</concept_significance>
</concept>
<concept>
<concept_id>10003752.10003790.10011740</concept_id>
<concept_desc>Theory of computation~Type theory</concept_desc>
<concept_significance>300</concept_significance>
</concept>
<concept>
<concept_id>10002950.10003624.10003633.10003646</concept_id>
<concept_desc>Mathematics of computing~Extremal graph theory</concept_desc>
<concept_significance>300</concept_significance>
</concept>
</ccs2012>
\end{CCSXML}
%% All four concept_ids above were verified against the official CCS tool
%% at https://dl.acm.org/ccs on 2026-07-30.

\ccsdesc[500]{Theory of computation~Proof theory}
\ccsdesc[300]{Theory of computation~Logic and verification}
\ccsdesc[300]{Theory of computation~Type theory}
\ccsdesc[300]{Mathematics of computing~Extremal graph theory}

\keywords{flag algebras, semidefinite programming, formal verification,
  interactive theorem proving, tactic metaprogramming, Lean, proof by reflection, 
  extremal graph theory, Tur\'an density}

\maketitle
\renewcommand{\shortauthors}{Jeong, Park, Hyun, Oum, and Yang}

\section{Introduction}
\label{sec:intro}

Large external computations have played a central role in proofs of landmark
mathematical results.  The four-color theorem was settled by extensive machine
case analysis and later formalized and checked in a proof
assistant~\cite{gonthier2008fourcolour}.  The Kepler conjecture was resolved by
optimization-heavy computation, and a complete formal proof required more than
a decade of further work~\cite{hales2017kepler}.  The empty-hexagon number was
determined with a SAT solver and later verified by formalizing the geometric
reduction in Lean and checking the solver's clausal proof with a formally
verified checker~\cite{heule2024hexagon,subercaseaux2024hexagon}.  Despite their
mathematical differences, these proofs share a computational structure.  A
search tool produces candidate evidence such as a case analysis, an optimization
certificate, or an unsatisfiability proof, and trust in the final theorem rests
on the validation of that evidence and the correctness of the reduction
that connects the evidence to the mathematical statement.

The programming languages and formal verification communities have developed
methods for validating results obtained by untrusted computation.  SAT solvers
emit unsatisfiability certificates that are checked by formally verified
checkers~\cite{lammich2020unsat,tan2023cakelpr}; SMT results are validated either
by reconstructing proofs inside a proof assistant~\cite{bohme2010z3} or by
checking proof witnesses with a verified checker~\cite{armand2011modular}; and
numerical sums-of-squares certificates for polynomial inequalities have their
entries rounded to rational numbers and are then reverified
exactly~\cite{harrison2007verifying,besson2006fast,morrison2026sos}.
Across these examples, the search procedures remain unverified; assurance comes
from independently validating the evidence produced by each run, much as
translation validation checks compilation
results~\cite{pnueli1998translation}.

This paper applies the same validation principle to Razborov's flag algebra
method, a technique for proving asymptotic density bounds in extremal graph
theory~\cite{razborov2007flag}.  At a high level, it reduces the proof of such a
bound to finding matrices and numerical weights that satisfy a finite system of
algebraic identities and positivity constraints.  A
semidefinite-programming (SDP) solver searches numerically for
such data, which are then converted to exact form using rational numbers and
recorded in a certificate.  Checking the resulting exact certificate against
these conditions establishes the original bound.  This validation is especially
valuable because the flag algebra method has been central to proofs of major
results in extremal graph theory, including the Erd\H{o}s
pentagon theorem and the determination of the maximum induced $C_5$
density~\cite{hatami2012,grzesik2012,balogh2016induced5}.  Yet established
toolchains for classical flag algebras, such as
Flagmatic~\cite{vaughan2013flagmatic}, generate and check their certificates
outside proof assistants.  Formally verifying these large toolchains, including
their graph-generation, exact-counting, and numerical-optimization components,
would be a major undertaking, and the certificates considered here had not
been converted into machine-checked proofs.  Rather than verify an entire
toolchain, we build a reusable certificate-to-proof compiler that treats the
toolchain and its output as untrusted, validates each certificate, and
constructs a complete Lean proof.

\boldparagraph{Flag algebras and their certificates.}
To see what such a bound and its certificate look like, consider a classical
problem about finite simple graphs.\footnote{By a finite simple graph, we mean
an undirected graph with a finite vertex set, no self-loops, and at most one
edge between any two distinct vertices.  Formally, it is a pair $G=(V,E)$,
where $V$ is finite and $E\subseteq\binom{V}{2}$; here $\binom{V}{2}$ denotes
the set of two-element subsets of $V$.  We write $V(G)$ and $E(G)$ for the
vertex set $V$ and the edge set $E$, respectively.}  What is the largest
possible fraction of vertex pairs that can be edges if the graph contains no
triangle?  As the number of vertices grows, Mantel's theorem gives the exact
answer~\cite{mantel}:
\[
  \lim_{n\to\infty}
  \max_{\substack{|V(G)|=n\\G\text{ triangle-free}}}
  \frac{|E(G)|}{\binom{n}{2}}
  = \frac{1}{2}.
\]
The expression on the left is the Tur\'an density of edges in triangle-free
graphs.  The two directions of the equality have very different proof burdens.
For the lower bound, one explicit construction suffices: balanced complete
bipartite graphs contain no triangles and have edge densities approaching
$1/2$.  The details of the construction are unimportant here; its role is
simply to witness that $1/2$ is attainable in the limit.  The upper bound is
the universal part: it must show that no family of increasingly large
triangle-free graphs can asymptotically exceed $1/2$.  This universal claim is
the part established by the flag-algebra certificate.

Flag algebras make such a certificate possible by reducing the universal claim
to a finite algebraic calculation over small, partially labeled graphs, called
\emph{flags}, that encode the relevant local configurations.  A
\emph{candidate certificate} selects finitely many flags, supplies matrices
with rational entries, and gives rational coefficients that weight the
resulting algebraic terms.  It asserts that these data satisfy specified
algebraic identities and that the matrices are positive semidefinite.  Once
these finite claims are verified, flag-algebra theory lifts them to the desired
asymptotic upper bound over the graph class under consideration.

Finite, however, does not mean small or readily auditable.  A nontrivial
certificate combines many flags and depends on large tables of exact counting
results for finite graphs.  Each table entry is a rational number recording the
fraction of suitable vertex selections in one small, partially labeled graph
that induce another specified pattern.  For the Erd\H{o}s pentagon theorem,
these tables contain thousands of entries; the certificate also includes
matrices whose positive semidefiniteness must be established using exact
rational arithmetic.  This scale is beyond realistic human audit, and a formal
development cannot simply trust data produced by a large, unverified numerical
toolchain.  It must check these obligations automatically inside the proof
assistant, without trusting the raw data or requiring thousands of handwritten
proofs.

\boldparagraph{A certificate-to-proof compiler.}
To perform these checks at scale, we build a Lean~4~\cite{moura2021lean}
development that formalizes the flag algebra method for finite simple graphs
and implements a certificate-to-proof compiler.  The compiler turns candidate
certificates into complete, machine-checked proofs without trusting the
certificate data.  The development has three layers.

The \emph{specification layer} (\Cref{sec:abstract}) formalizes the
mathematics of the method: partially labeled graphs and their densities
in large graphs, the quotient algebra of density expressions, and its
semantics through positive homomorphisms.  It
follows the standard development closely, so flag-algebra arguments can
be stated in Lean much as they appear on paper.

The \emph{reflection layer} (\Cref{sec:reflection}) makes the
specification computable.  It supplies executable graph
representations, adequacy theorems proving that their computations
agree with the specification, and elaboration-time commands that
generate the verified identities a given certificate requires.

The \emph{automation layer} (\Cref{sec:flagmatic}) drives the
compilation.  It treats the output of the external SDP search as
candidate data only: Lean recomputes every claimed fact against the
mathematical definitions, verifies positive semidefiniteness exactly
over $\mathbb{Q}$ through $LDL^\top$ factorizations, and assembles the
final bound using the reflected identities and custom tactics.  The
compilation happens at elaboration time, inside Lean: the theorem
statement is fixed in the source file, and the certificate enters only
as untrusted hint data from which the proof is synthesized.

\boldparagraph{A formalization-driven semantics for forbidden subgraphs.}
Formalization did more than add assurance to the existing method: it suggested
a different organization of the mathematics.  All seven compiled bounds
concern graphs that avoid a specified subgraph $H$.  Razborov's standard
presentation builds this constraint into the flag algebra from the outset,
constructing a separate algebra for each forbidden family.  To make our Lean
development reusable across problems, we instead retain one ambient algebra of
all finite graphs and impose $H$-freeness only when interpreting an inequality;
we call this alternative the \emph{ensemble semantics}.  This design simplifies
the formalization by keeping its core definitions and lemmas independent of the
forbidden family.  Its use in the compiler is justified directly by flag-algebra
theory: we prove that the ensemble inequalities the compiler establishes yield
the corresponding asymptotic bounds in extremal graph theory.  The standard
semantics plays no role in this justification, so equivalence between the two
is not required for the soundness of the compiled bounds.

The alternative semantics nevertheless exposed a natural mathematical
question: when do the two semantics validate the same inequalities?  This
comparison is separate from the preceding justification, but it led to new
mathematics.
\Cref{sec:metatheory} proves that built-in validity always implies ensemble
validity and characterizes, through a structural condition that we call
\emph{root-plantability}, when the converse holds for all inequalities.  This
result illustrates how formalization can do more than check existing arguments:
it can also reveal new mathematical structures and questions.

\boldparagraph{Contributions.}
We make the following contributions.
\begin{itemize}
\item \emph{A Lean~4 formalization and certificate-to-proof compiler for flag
  algebras.}  Invoked as an elaboration-time tactic, the compiler turns
  externally generated, untrusted certificates into complete,
  machine-checked proofs of Tur\'an-type bounds.
\item \emph{Machine-checked extremal bounds}: seven compiled Tur\'an-type upper
  bounds.  For Mantel's theorem and the Erd\H{o}s pentagon theorem, we also
  prove matching lower-bound constructions in Lean, yielding exact theorems.
  To our knowledge, this is the first machine-checked proof of the Erd\H{o}s
  pentagon theorem; all known proofs use flag algebras.
\item \emph{A formalization-driven semantics for forbidden subgraphs.}  The
  ensemble semantics keeps the core formalization independent of the forbidden
  family, and we prove directly that it yields the intended asymptotic bounds.
  We also present selected results from the
  mathematical comparison it prompted, together with their autoformalization;
  a separate mathematics paper will develop the broader meta-theory and its
  applications.
\end{itemize}
The complete Lean development, including the certificate-to-proof compiler
and the formalized semantic results presented here, is  available at
\url{https://github.com/taeyool/lean-flag-algebras-release}.

\boldparagraph{Organization.}
\Cref{sec:background} reviews the mathematical background of flag algebras
and fixes the notation used throughout the paper.
\Cref{sec:abstract,sec:reflection,sec:flagmatic} present the specification,
reflection, and automation layers, respectively: the specification layer
includes the semantic treatment of forbidden-subgraph assumptions, and the
automation layer includes the evaluation of the compiler on seven certificates.
\Cref{sec:lowerbounds} illustrates direct uses of the formalization beyond
certificate compilation, and \Cref{sec:obstacles} analyzes the main engineering
obstacles.  \Cref{sec:metatheory} presents the selected results comparing the
ensemble and built-in semantics that are included in this paper.
\Cref{sec:unsure} records the remaining design questions,
\Cref{sec:related} discusses related work, and \Cref{sec:conclusion} concludes
the paper.

\section{Background: Flag Algebras}
\label{sec:background}

This section recalls the mathematical definitions underlying flag algebras,
following Razborov~\cite{razborov2007flag}, and establishes the notation used
later in the paper.  No background in flag algebras or extremal graph theory
is assumed.

\boldparagraph{Notation.}
For a non-negative integer $n$, we write $[n]$ for the set of positive integers
at most $n$.  We consider only undirected, finite, and simple graphs; for a
graph $G$, its vertex and edge sets are denoted by $V(G)$ and $E(G)$.  We use
the standard graph names: $K_n$ is the complete graph on $n$ vertices, $P_n$
the path and $C_n$ the cycle on $n$ vertices, $\overline{G}$ the complement of
$G$, and $K_{a,b}$ the complete bipartite graph with parts of sizes $a$
and~$b$.  We fix a (possibly empty) finite set $\mathcal{H}$ of forbidden
graphs; a graph is \emph{\(\mathcal{H}\)-free} if it contains no subgraph
isomorphic to any member of $\mathcal{H}$.  All flag types, flags, and flag
algebras defined below are relative to this fixed $\mathcal{H}$.

\subsection{Flag Types and Flags}
\label{sec:background-flags}

\boldparagraph{Flag Types.}
A flag type specifies the labeled part shared by all flags of that flag type.
Formally, a \emph{flag type} of size $k$ is an \(\mathcal{H}\)-free graph
\(\sigma\) with vertex set $[k]$.  The \emph{empty flag type} $\emptyset$, of
size zero, is the unique graph on the empty vertex set.
We use the term \emph{flag type} rather than Razborov's original \emph{type}
to avoid confusion with Lean's type system in later sections.

\boldparagraph{Flags.}
For a graph $G$ and a set $S\subseteq V(G)$, let $G[S]$ denote the
\emph{induced subgraph} of $G$ on $S$: its vertex set is $S$, and its edges are
exactly the edges of $G$ with both endpoints in $S$.
Fix a flag type $\sigma$ of size $k$.
A \emph{$\sigma$-flag} is a pair $F = (M, \theta_F)$ consisting of an
\(\mathcal{H}\)-free finite graph $M$ and an injective map
$\theta_F : [k] \hookrightarrow V(M)$ that induces a graph isomorphism from
$\sigma$ to $M[\operatorname{Im}(\theta_F)]$.  The \emph{size} of $F$ is
$|V(M)|$, and we write $V(F)$ for $V(M)$.
For each $i\in[k]$, the vertex $\theta_F(i)$ carries label $i$, so the vertices
in~$\operatorname{Im}(\theta_F)$ are the \emph{labeled vertices} of $F$, also
called its \emph{roots}; all remaining vertices are \emph{unlabeled}.
For $S \subseteq V(F)$ with $\operatorname{Im}(\theta_F) \subseteq S$,
we write $F[S]$ for the $\sigma$-flag $(M[S], \theta_F)$, the
\emph{$\sigma$-flag induced by $F$ on $S$}.
Flags over the empty flag type $\emptyset$ are ordinary
\(\mathcal{H}\)-free graphs with no labeled vertices.
\Cref{fig:flag-examples} illustrates several flags of different flag types.

Two $\sigma$-flags $(M_1,\theta_1)$ and $(M_2,\theta_2)$ are
\emph{isomorphic}, written $(M_1,\theta_1) \cong_\sigma (M_2,\theta_2)$, if
some graph isomorphism $\alpha: V(M_1)\to V(M_2)$ satisfies
$\alpha \circ \theta_1 = \theta_2$, an equation saying that $\alpha$
preserves every label.  We write $\mathcal{F}^\sigma_n$ for the set of
isomorphism classes of $\sigma$-flags of size $n$, and
$\mathcal{F}^\sigma = \bigcup_{n \geq k} \mathcal{F}^\sigma_n$.

\begin{figure}[t]
  \centering
  \setlength{\tabcolsep}{1.2em}
  \renewcommand{\arraystretch}{1.3}
  \begin{tabular}{cccccc}
    \begin{tikzpicture}[scale=0.8, baseline=-0.5ex]
      \node[graph vertex] (a) at (0,0) {};
      \node[graph vertex] (b) at (0.8,0) {};
      \draw[graph edge] (a) -- (b);
    \end{tikzpicture}
    &
    \begin{tikzpicture}[scale=0.8, baseline=-0.5ex]
      \node[graph vertex] (a) at (90:0.65) {};
      \node[graph vertex] (b) at (210:0.65) {};
      \node[graph vertex] (c) at (330:0.65) {};
    \end{tikzpicture}
    &
    \begin{tikzpicture}[scale=0.8, baseline=-0.5ex]
      \node[graph vertex] (a) at (90:0.65) {};
      \node[graph vertex] (b) at (210:0.65) {};
      \node[graph vertex] (c) at (330:0.65) {};
      \draw[graph edge] (b) -- (c);
    \end{tikzpicture}
    &
    \begin{tikzpicture}[scale=0.8, baseline=-0.5ex]
      \node[graph vertex] (a) at (90:0.65) {};
      \node[graph vertex] (b) at (210:0.65) {};
      \node[graph vertex] (c) at (330:0.65) {};
      \draw[graph edge] (a) -- (b);
      \draw[graph edge] (a) -- (c);
    \end{tikzpicture}
    &
    \begin{tikzpicture}[scale=0.8, baseline=-0.5ex]
      \node[graph vertex] (a) at (90:0.65) {};
      \node[graph vertex] (b) at (210:0.65) {};
      \node[graph vertex] (c) at (330:0.65) {};
      \draw[graph edge] (a) -- (b);
      \draw[graph edge] (a) -- (c);
      \draw[graph edge] (b) -- (c);
    \end{tikzpicture}
    &
    \begin{tikzpicture}[scale=0.7, baseline=-0.5ex]
      \node[graph vertex] (v1) at (90:0.65) {};
      \node[graph vertex] (v2) at (18:0.65) {};
      \node[graph vertex] (v3) at (-54:0.65) {};
      \node[graph vertex] (v4) at (-126:0.65) {};
      \node[graph vertex] (v5) at (162:0.65) {};
      \draw[graph edge] (v1) -- (v2) -- (v3) -- (v4) -- (v5) -- (v1);
    \end{tikzpicture}
    \\
    $K_2$ & $\overline{K_3}$ & $\overline{P_3}$ & $P_3$ & $K_3$ & $C_5$ \\[0.8em]
    \begin{tikzpicture}[scale=0.8, baseline=-0.5ex]
      \node[graph root] (a) at (0,0) {};
      \node[graph vertex] (b) at (0.8,0) {};
      \draw[graph edge] (a) -- (b);
    \end{tikzpicture}
    &
    \begin{tikzpicture}[scale=0.8, baseline=-0.5ex]
      \node[graph root] (a) at (0,0) {};
      \node[graph vertex] (b) at (0.8,0) {};
    \end{tikzpicture}
    &
    \begin{tikzpicture}[scale=0.8, baseline=-0.5ex]
      \node[graph vertex] (a) at (90:0.65) {};
      \node[graph root] (b) at (210:0.65) {};
      \node[graph vertex] (c) at (330:0.65) {};
      \draw[graph edge] (b) -- (c);
    \end{tikzpicture}
    &
    \begin{tikzpicture}[scale=0.8, baseline=-0.5ex]
      \node[graph vertex] (a) at (90:0.65) {};
      \node[graph root] (b) at (210:0.65) {};
      \node[graph vertex] (c) at (330:0.65) {};
      \draw[graph edge] (b) -- (a);
      \draw[graph edge] (a) -- (c);
    \end{tikzpicture}
    &
    \begin{tikzpicture}[scale=0.8, baseline=-0.5ex]
      \node[graph vertex] (a) at (90:0.65) {};
      \node[graph root] (b) at (210:0.65) {};
      \node[graph vertex] (c) at (330:0.65) {};
      \draw[graph edge] (b) -- (a);
      \draw[graph edge] (b) -- (c);
    \end{tikzpicture}
    &
    \begin{tikzpicture}[scale=0.7, baseline=-0.5ex]
      \node[graph root] (v1) at (90:0.65) {};
      \node[graph vertex] (v2) at (18:0.65) {};
      \node[graph vertex] (v3) at (-54:0.65) {};
      \node[graph vertex] (v4) at (-126:0.65) {};
      \node[graph vertex] (v5) at (162:0.65) {};
      \draw[graph edge] (v1) -- (v2) -- (v3) -- (v4) -- (v5) -- (v1);
    \end{tikzpicture}
    \\
    $K_2^\bullet$ & $\overline{K_2}^\bullet$ & $\overline{P_3}^\bullet$ & $P_3^\bullet$
    & ${P_3'}^\bullet$ & $C_5^\bullet$ \\
  \end{tabular}
  \Description{Two rows of small graphs drawn as dots and lines.  The top
  row shows ordinary unlabeled graphs on two, three, and five vertices; the
  bottom row shows partially labeled graphs in which one root vertex is
  drawn as a filled dot.}
  \caption{Examples of flags when no graphs are forbidden
  ($\mathcal{H}=\varnothing$).  \emph{Top row}: $\emptyset$-flags, equivalently
  ordinary graphs, of sizes 2, 3, and 5.  \emph{Bottom row}: $\sigma$-flags of the
  unique one-vertex flag type $\sigma$.  The filled vertex carries label~1,
  and the open vertices are unlabeled.  The $\sigma$-flags $P_3^\bullet$ and
  ${P_3'}^\bullet$ have the same underlying path; $P_3^\bullet$ labels an
  endpoint, whereas ${P_3'}^\bullet$ labels the middle vertex.}
  \label{fig:flag-examples}
\end{figure}

\subsection{Subflag Densities}
\label{sec:background-densities}

In a density calculation, flags play two roles.  The \emph{pattern flag} is the
small local configuration whose occurrence is being measured, while the
\emph{host flag} is the larger flag in which we sample vertices.  These are
roles in the calculation, not different kinds of flags.  We use \(F\), or
\(F_i\), for pattern flags, \(G\) for host flags, and \(G'\) for an intermediate
host flag when a third flag variable is needed.

Let $\sigma$ be a flag type of size $k$.  

\boldparagraph{Single-flag density.}
For a $\sigma$-flag $F$ of size $m$
(the pattern $\sigma$-flag) and a $\sigma$-flag $G = (M, \theta_G)$ of size $n \geq m$
(the host $\sigma$-flag), the \emph{density of $F$ in $G$}, written $\den{F}{G}$, is
the probability that a uniformly random $(m-k)$-element subset of
$V(G) \setminus \operatorname{Im}(\theta_G)$, together with the $k$ labeled
vertices of $G$, induces a $\sigma$-flag \(F'\) with \(F' \cong_\sigma F\).
Explicitly:
\[
  \den{F}{G}
  \;:=\;
  \frac{1}{\binom{n-k}{m-k}}
    \cdot {\Bigl|\Big\{S \subseteq V(G)\setminus \operatorname{Im}(\theta_G)
                  ~\Big|~ |S|=m-k,\;
                  G[S \cup \operatorname{Im}(\theta_G)]
                  \cong_\sigma F \Big\}\Bigr|}.
\]
This value lies in $[0,1]$ and is invariant under
isomorphism of both $F$ and $G$.
When $n < m$, we set $\den{F}{G} := 0$ by convention.

\begin{example}[Subflag densities in $C_5^\bullet$]
  Among the four unlabeled vertices of
  $\fCfivebull$, exactly two are adjacent to the labeled one.  Hence
  \[
    \den{\fKtwobull}{\fCfivebull} = \frac{2}{4} = \frac12,
    \qquad
    \den{\fbarKtwobull}{\fCfivebull} = \frac{2}{4} = \frac12.
  \]
\end{example}

\boldparagraph{Chain rule.}
The density $\den{F}{G}$ can be computed in two equivalent ways: by sampling
$m-k$ unlabeled vertices directly from $G$, or by first sampling an
intermediate $\sigma$-flag $G'$ of some size $n'$ with $m \leq n' \leq n$ inside $G$ and
then sampling~$F$ inside $G'$.
These procedures describe the same sampling experiment, yielding the
single-flag case of Razborov's chain rule~\cite[Lemma~2.2]{razborov2007flag}.

\begin{lemma}[Chain rule]\label{lem:chain-rule}
  Let $\sigma$ be a flag type of size $k$, let $m,n,n' \in \mathbb{N}$ satisfy
  $k \leq m \leq n' \leq n$, and let $F \in \mathcal{F}^\sigma_m$ and
  $G \in \mathcal{F}^\sigma_n$.  Then
  \[
    \den{F}{G}
    \;=\;
    \sum_{G' \in \mathcal{F}^\sigma_{n'}} \den{F}{G'}\,\den{G'}{G}.
  \]
\end{lemma}

\boldparagraph{Multi-flag density.}
For pattern $\sigma$-flags $F_1, \ldots, F_t$ of sizes $m_1, \ldots, m_t$ and
a host $\sigma$-flag $G = (M, \theta_G)$ of size $n$, we say $F_1, \ldots, F_t$
\emph{fit} in $G$ if
\[
  n - k \;\geq\; (m_1 - k) + \cdots + (m_t - k).
\]
When $F_1, \ldots, F_t$ fit in $G$, let $(S_1, \ldots, S_t)$ be a uniformly
random tuple of pairwise-disjoint subsets $S_i \subseteq V(G) \setminus
\operatorname{Im}(\theta_G)$ with $|S_i| = m_i - k$.
The \emph{density of $F_1, \ldots, F_t$ in $G$}, written $\den{F_1,\ldots,F_t}{G}$, is the
probability that $G[S_i \cup \operatorname{Im}(\theta_G)] \cong_\sigma F_i$
for every $i \in [t]$.
When $F_1, \ldots, F_t$ do not fit in $G$, we set
$\den{F_1,\ldots,F_t}{G} := 0$ by convention.

\subsection{The Flag Algebra}

Fix a flag type $\sigma$ of size $k$.
Informally, the flag algebra $\mathcal{A}^\sigma$ is an algebra\footnote{Here
``algebra'' means a vector space whose elements can also be multiplied.  In
this paper, the scalar coefficients are real numbers, multiplication is
bilinear, associative, and commutative, and there is a multiplicative
unit.} of density expressions generated
by $\sigma$-flags.  Each $\sigma$-flag is a symbolic building block: it stands for the
density of that local labeled pattern in a large host graph.  The quotient
below records the identity of Lemma~\ref{lem:chain-rule} among these building
blocks.
Formally, let
$\mathbb{R}[\mathcal{F}^\sigma]$ be the free real vector space on
$\mathcal{F}^\sigma$, that is, the vector space of finite formal real linear
combinations of $\sigma$-flags.  For a $\sigma$-flag $F$, write $\mathbf e_F$
for its corresponding basis vector.
Define the \emph{zero space} $\mathcal{Z}^\sigma$ as the subspace generated
by all elements of the form
\begin{equation}
  \label{eq:zero-element}
  \mathbf e_F - \sum_{G \in \mathcal{F}^\sigma_n} \den{F}{G} \cdot \mathbf e_G,
  \qquad \text{for}\ F \in \mathcal{F}^\sigma_m \ \text{and}\ n \geq m.
\end{equation}
The \emph{flag algebra} is the quotient vector space
\[
  \mathcal{A}^\sigma \;:=\; \mathbb{R}[\mathcal{F}^\sigma] \,/\, \mathcal{Z}^\sigma.
\]
We write $[F]$ for the class of $\mathbf e_F$ in $\mathcal{A}^\sigma$.
Throughout the paper, capital letters $F,G$ denote finite flags, while
lowercase letters $f,g,h$ denote arbitrary flag-algebra elements.
For familiar named or pictorial flags, such as $K_2$ or
$\fKtwobull$, we retain the standard shorthand of omitting the brackets in
algebraic formulas; there the surrounding algebraic context makes the
canonical image unambiguous.

We call $\sum_{G \in \mathcal{F}^\sigma_n} \den{F}{G} \cdot \mathbf e_G$ the
\emph{expansion of $F$ at size $n$}.
By Lemma~\ref{lem:chain-rule}, a $\sigma$-flag $F \in \mathcal{F}^\sigma_m$ and
its expansion at any size $n \geq m$ have the same density in every host
$\sigma$-flag of size $\geq n$.
The quotient by $\mathcal{Z}^\sigma$ records this expansion identity as the
equality
\begin{equation}
  \label{eq:expansion-identity}
  [F]
  =
  \sum_{G \in \mathcal{F}^\sigma_n} \den{F}{G}\,[G]
  \qquad\text{in }\mathcal{A}^\sigma .
\end{equation}
We may therefore freely replace $[F]$ by its expansion at any size $n \geq m$.

\begin{example}[$K_2$ expanded at size $3$ in the triangle-free algebra]
  Working in the triangle-free flag algebra (i.e., $\mathcal{H} = \{K_3\}$),
  the $\emptyset$-flags of size $3$ are the three triangle-free graphs on three vertices:
  $\fbarKthree$, $\fbarPthree$, and $\fPthree$.
  The density of an edge in each of them is
  \[
    \den{K_2}{\fbarKthree}=0,\quad
    \den{K_2}{\fbarPthree}=\frac13,\quad
    \den{K_2}{\fPthree}=\frac23,
  \]
  so the zero-space quotient identifies
  \begin{equation}
    K_2
    \;=\;
    \frac13\,\fbarPthree
    + \frac23\,\fPthree
    \qquad\text{in } \mathcal{A}^{\emptyset}.
    \label{eq:k2-expansion}
  \end{equation}
\end{example}

\boldparagraph{Multiplication.}
For $m_1,m_2 \in \mathbb{N}$ and $\sigma$-flags
$F_1 \in \mathcal{F}^\sigma_{m_1}$ and $F_2 \in \mathcal{F}^\sigma_{m_2}$,
choose an auxiliary size $\ell \in \mathbb{N}$ with
$\ell \geq m_1 + m_2 - k$ and define
\begin{equation}
  \label{eq:mul}
  [F_1] \cdot [F_2]
  \;:=\;
  \sum_{G \in \mathcal{F}^\sigma_\ell} \den{F_1,F_2}{G} \cdot [G]
  \qquad\text{in } \mathcal{A}^\sigma.
\end{equation}
At first glance, this definition appears to depend on both the choice of $\ell$
and the choice of representatives for~$[F_1]$ and $[F_2]$; the zero-space
quotient ensures that neither matters.  Extending bilinearly to all of
$\mathcal{A}^\sigma$ gives a commutative algebra over the real numbers.  Its
unit is the basis element represented by the $\sigma$-flag
$(\sigma, \operatorname{id}_{[k]})$ of size~$k$, whose underlying graph
is $\sigma$ itself with every vertex labeled; we write it as
$\mathbf{1}_\sigma$, and simply as $\mathbf{1}$ when
$\sigma=\emptyset$.

\begin{example}[Multiplication of $K_2^\bullet$ and $\overline{K_2}^\bullet$]
\label{ex:typed-edge-nonedge-product}
  Let $\sigma$ be the one-vertex flag type.  Taking $\ell = 3$, the density
  $\den{\fKtwobull,\,\fbarKtwobull}{G}$ equals $\frac{1}{2}$ for each size-3
  $\sigma$-flag $G$ whose labeled vertex has degree exactly $1$:
  \[
    \den{\fKtwobull,\,\fbarKtwobull}{\fbarPthreebull}
    \;=\;
    \den{\fKtwobull,\,\fbarKtwobull}{\fPthreebull}
    \;=\; \frac{1}{2},
  \]
  and vanishes on all other size-3 $\sigma$-flags:
  \[
    \den{\fKtwobull,\,\fbarKtwobull}{G} = 0
    \qquad\text{for all } G \in \mathcal{F}_3^\sigma \setminus \{\fbarPthreebull,\,\fPthreebull\}.
  \]
  Hence
  \[
    \fKtwobull \cdot \fbarKtwobull
    \;=\;
    \frac{1}{2}\,\fbarPthreebull + \frac{1}{2}\,\fPthreebull
    \qquad\text{in } \mathcal{A}^\sigma.
  \]
\end{example}

\subsection{Positive Homomorphisms and Semantic Order}
\label{sec:background-semantic-order}

Fix a flag type $\sigma$.  
The flag algebra $\mathcal{A}^\sigma$ provides algebraic syntax for expressing
densities of pattern $\sigma$-flags.  Flag algebra proofs, however, target
\emph{asymptotic} statements: rather than a single finite host $\sigma$-flag, we
consider increasing sequences~$\{G_n\}_{n\in\mathbb{N}}$ of host
$\sigma$-flags with $|V(G_n)|\to\infty$ and ask what happens to $\sigma$-flag
densities in the limit.  Positive homomorphisms can be thought of as the
limiting density profiles of convergent sequences of host $\sigma$-flags.
The semantic non-negativity cone
$\mathcal{C}^\sigma$ then consists of the flag-algebra elements
whose value under every positive homomorphism is non-negative, and establishing
membership in $\mathcal{C}^\sigma$ is precisely what flag algebra proofs aim
to do.

\boldparagraph{Convergent sequences.}
Call a sequence $\{G_n\}_{n\in\mathbb{N}}$ of host $\sigma$-flags \emph{increasing} if
$|V(G_1)| < |V(G_2)| < \cdots$.  Each host $\sigma$-flag $G$ defines a point
$p^G \in [0,1]^{\mathcal{F}^\sigma}$ by $p^G(F) := \den{F}{G}$.
An increasing sequence is \emph{convergent} if the sequence of points $p^{G_n}$
converges in $[0,1]^{\mathcal{F}^\sigma}$ endowed with the product topology;
equivalently, if $\lim_{n\to\infty}\den{F}{G_n}$ exists for every $\sigma$-flag $F$.

\boldparagraph{Positive homomorphisms.}
A \emph{positive homomorphism} on
$\mathcal{A}^\sigma$ is a map
$\phi:\mathcal{A}^\sigma \to \mathbb{R}$ that preserves addition,
multiplication, multiplication by real scalars, and the unit.  It must also
satisfy $\phi([F]) \geq 0$ for every $\sigma$-flag $F$.
Following standard flag-algebra terminology, we call such a homomorphism
\emph{positive}, although its defining inequalities are non-strict.
We write $\poshom{\sigma}$ for the set of all positive homomorphisms on
$\mathcal{A}^\sigma$.

In $\mathcal{A}^\sigma$, the following identity holds for every
$\ell \in \mathbb{N}$ with $\ell \geq |V(\sigma)|$:
\begin{equation}
  \sum_{F \in \mathcal{F}^\sigma_\ell} [F] \;=\; \mathbf{1}_\sigma.
  \label{eq:sum-to-one}
\end{equation}
Applying any $\phi \in \poshom{\sigma}$ to both sides gives
$\sum_{F \in \mathcal{F}^\sigma_\ell} \phi([F]) = 1$, and non-negativity then
forces $\phi([F]) \in [0,1]$ for every $\sigma$-flag $F$.  Each
$\phi \in \poshom{\sigma}$ therefore determines the point
$F \mapsto \phi([F])$ of $[0,1]^{\mathcal{F}^\sigma}$, identifying
$\poshom{\sigma}$ with a subset of that space.  %

The following theorem makes precise the sense in which positive homomorphisms
are exactly the limiting density profiles of convergent sequences of
$\sigma$-flags~\cite{razborov2007flag}.

\begin{theorem}[Razborov]\label{thm:convergent-hom}
  Let $\sigma$ be a flag type.
  \begin{enumerate}
  \item[\textup{(a)}] If $\{G_n\}_{n\in\mathbb{N}}$ is a convergent sequence of
    $\sigma$-flags, then
    \[
      \lim_{n\to\infty} p^{G_n} \in \poshom{\sigma}.
    \]
  \item[\textup{(b)}] Conversely, every $\phi \in \poshom{\sigma}$ satisfies
    $\phi = \lim_{n\to\infty} p^{G_n}$ for some convergent sequence
    $\{G_n\}_{n\in\mathbb{N}}$ of $\sigma$-flags.
  \end{enumerate}
\end{theorem}

\boldparagraph{Semantic non-negativity cone and order.}
For a fixed flag type $\sigma$, define the \emph{semantic non-negativity cone}
\[
  \mathcal{C}^\sigma
  :=
  \{f \in \mathcal{A}^\sigma
    \mid \phi(f) \geq 0
    \text{ for every } \phi \in \poshom{\sigma}\}.
\]
We write
\[
  f \leq_{\mathcal{H},\sigma} g
  \qquad\text{when}\qquad
  g-f \in \mathcal{C}^\sigma,
\]
where the subscript $\mathcal{H}$ records that the comparison is made in the
\(\mathcal{H}\)-free world: both $\mathcal{A}^\sigma$ and $\mathcal{C}^\sigma$ are
built from \(\mathcal{H}\)-free graphs.
Equivalently, $f\leq_{\mathcal{H},\sigma} g$ means that
$\phi(f)\leq\phi(g)$ for every $\phi \in \poshom{\sigma}$.  When $\sigma=\emptyset$, we write $f \leq_{\mathcal{H}} g$ in place of $f \leq_{\mathcal{H},\emptyset} g$, and $\mathcal{C}$ in place of $\mathcal{C}^\emptyset$.
Note that $f \in \mathcal{C}^\sigma$ is the same as saying $0 \leq_{\mathcal{H},\sigma} f$.

Three facts about $\mathcal{C}^\sigma$ are immediate from the definition.
Every $\sigma$-flag $F$ gives an element
$[F]\in\mathcal{C}^\sigma$, because positive homomorphisms send every flag
basis element to a non-negative real by definition.
For any $f \in \mathcal{A}^\sigma$, the square $f^2$ satisfies $0 \leq_{\mathcal{H},\sigma} f^2$,
because $\phi(f^2) = \phi(f)^2 \geq 0$ for every~$\phi \in \poshom{\sigma}$.
And since every $\phi \in \poshom{\sigma}$ is linear, $\mathcal{C}^\sigma$ is
closed under addition and under multiplication by non-negative reals.

\boldparagraph{From semantic non-negativity to Tur\'an density.}
For a finite simple graph $F$ and a finite collection $\mathcal{H}$ of finite
simple graphs, write $\mathrm{ex}(n,\,F;\,\mathcal{H})$ for the maximum number
of induced copies of $F$ in an $n$-vertex graph that contains no member of
$\mathcal{H}$ as a (not necessarily induced) subgraph.  Thus an induced copy
of $F$ preserves both edges and nonedges, whereas a subgraph copy of a member
of $\mathcal{H}$ need preserve only its edges; additional edges between the
corresponding vertices are allowed.  The associated
\emph{Tur\'an density} is
\[
  \tdensity{F}{\mathcal{H}}
  :=
  \lim_{n\to\infty}
  \frac{\mathrm{ex}(n,\,F;\,\mathcal{H})}{\binom{n}{|V(F)|}}.
\]
It records the asymptotic proportion of $|V(F)|$-vertex subsets that induce a
copy of $F$.  Whenever $\mathcal{H}=\{H\}$ is a singleton, we write $H$ in
place of $\mathcal{H}$.

The problem of bounding the Tur\'an density $\tdensity{F}{\mathcal{H}}$ from above
reduces to establishing membership in $\mathcal{C}^\emptyset$, as the following
lemma shows.

\begin{lemma}\label{lem:semantic-bound}
  Let $F \in \mathcal{F}^\emptyset$ and $c \in \mathbb{R}$.  If
  $[F] \leq_{\mathcal{H}} c \cdot \mathbf{1}$, then
  $\tdensity{F}{\mathcal{H}} \leq c$.
\end{lemma}
\begin{proof}
  Suppose for contradiction that $\tdensity{F}{\mathcal{H}} > c$.  Then some
  $\varepsilon > 0$ and some increasing sequence of \(\mathcal{H}\)-free graphs
  $\{G_k\}_{k\in\mathbb{N}}$ satisfy $\den{F}{G_k} > c + \varepsilon$ for all
  $k$.  The points $p^{G_k}$ lie in the product space
  $[0,1]^{\mathcal{F}^\emptyset}$, which is compact and metrizable because
  $\mathcal{F}^\emptyset$ is countable.  Hence some subsequence
  $\{p^{G_{k_j}}\}$ converges, and by \Cref{thm:convergent-hom}(a) its limit
  is some $\phi \in \poshom{\emptyset}$.
  Since $\phi$ is the pointwise limit of $\{p^{G_{k_j}}\}$, it satisfies
  $\phi([F]) = \lim_j \den{F}{G_{k_j}} \geq c + \varepsilon > c$.  But
  $[F] \leq_{\mathcal{H}} c \cdot \mathbf{1}$ forces $\phi([F]) \leq c$ for
  every~$\phi \in \poshom{\emptyset}$, a contradiction.
\end{proof}

A matching lower bound typically comes from an explicit \(\mathcal{H}\)-free construction
whose $F$-density approaches~$c$, pinning down $\tdensity{F}{\mathcal{H}}$ exactly.
Two classical examples illustrate this.

\begin{example}[Mantel's theorem~\cite{mantel}]\label{ex:mantel}
  Mantel's theorem states $\tdensity{K_2}{K_3} = \tfrac{1}{2}$, meaning
  asymptotically at most half of all vertex pairs in a triangle-free graph
  are edges.  The upper bound translates into the flag algebra inequality
  $[K_2] \leq_{K_3} \tfrac{1}{2}\cdot\mathbf{1}$; the balanced complete bipartite
  graphs $K_{\lfloor n/2\rfloor,\lceil n/2\rceil}$ witness the lower bound.
\end{example}

\begin{example}[Erd\H{o}s Pentagon Theorem~{\cite{grzesik2012,hatami2012}}]\label{ex:pentagon}
  The Erd\H{o}s pentagon problem asks for the maximum density of $C_5$ in
  triangle-free graphs; the asymptotic answer is $\tdensity{C_5}{K_3} = \tfrac{24}{625}$.  The upper bound corresponds to the flag algebra inequality
  $[C_5] \leq_{K_3} \tfrac{24}{625}\cdot\mathbf{1}$; an explicit \(K_3\)-free construction
  (the balanced blow-up of $C_5$) witnesses the lower bound.
\end{example}

\subsection{The Downward Operator}
\label{sec:background-downward}
Labeled flag algebras (i.e., $\mathcal{A}^\sigma$ with $\sigma \neq \emptyset$) are often
richer than the unlabeled algebra $\mathcal{A}^\emptyset$: establishing membership
in $\mathcal{C}^\sigma$ can be easier than working directly in $\mathcal{C}$.
The final extremal statement, however, is always about unlabeled graphs.
The downward operator bridges this gap by mapping $\mathcal{A}^\sigma$ to
$\mathcal{A}^\emptyset$ in a way that preserves semantic non-negativity
(\Cref{thm:downward-nonneg}), so valid inequalities proved in a labeled algebra
can be transferred to the unlabeled one.

For a $\sigma$-flag $F = (M, \theta_F)$, let $F|_\emptyset$ denote the
$\emptyset$-flag obtained from $F$ by forgetting the embedding $\theta_F$,
i.e., viewing $M$ as a $\emptyset$-flag.
For a $\sigma$-flag $F$ of size $m$ with $|V(\sigma)|=k$, let
\[
  q_\sigma(F)
  \;:=\;
  \frac{
    \bigl|\{\theta' : [k]\hookrightarrow V(M)
       \mid \theta' \text{ embeds }\sigma\text{ into }M
       \text{ and } (M,\theta')\cong_\sigma F\}\bigr|}
       {m(m-1)\cdots(m-k+1)}
\]
be the probability that a uniformly random injection $[k]\hookrightarrow V(M)$
recovers the same $\sigma$-flag as $F$.
The \emph{downward operator}
$\llbracket\cdot\rrbracket_\sigma:\mathcal{A}^\sigma\to\mathcal{A}^\emptyset$
is then defined on flag basis elements by
\[
  \llbracket [F] \rrbracket_\sigma
  \;:=\;
  q_\sigma(F)\cdot [F|_\emptyset],
\]
extended linearly to all of $\mathcal{A}^\sigma$.
For $f\in\mathcal{A}^\sigma$, we call
$\llbracket f\rrbracket_\sigma$ the \emph{downward image} of $f$.

We write $\langle\sigma\rangle_0\in\mathcal{A}^\emptyset$ for the empty-type
flag-algebra basis element represented by the underlying graph of the flag type
$\sigma$, with its vertex labels forgotten.  Applying the downward operator to
the typed unit therefore gives
\[
  \llbracket \mathbf{1}_\sigma \rrbracket_\sigma
  =
  q_\sigma(\mathbf{1}_\sigma)\,\langle\sigma\rangle_0.
\]

\boldparagraph{Random homomorphisms.}
A positive homomorphism $\phi_0 \in \poshom{\emptyset}$ records the limiting
densities of unlabeled graphs along a convergent sequence of finite graphs.
For a nonempty flag type $\sigma$, its random extension can be understood
through any such sequence realizing $\phi_0$: in each finite graph, choose
uniformly at random an embedding of $\sigma$ and use the embedded vertices as
labels.  Razborov's random-extension theorem describes the limiting
distribution of the resulting $\sigma$-flag densities.
Whenever $\phi_0(\langle\sigma\rangle_0) > 0$, meaning that the underlying
unlabeled graph of $\sigma$ has positive limiting density, the following
holds~\cite{razborov2007flag}:
there exists a unique probability measure on $\poshom{\sigma}$,
which we denote $\Ext_\sigma(\phi_0)$,
such that for every~$f \in \mathcal{A}^\sigma$, the random positive
homomorphism $\phi^\sigma \sim \Ext_\sigma(\phi_0)$ satisfies
\begin{equation}\label{eq:ensemble}
  \mathbb{E}\bigl[\phi^\sigma(f)\bigr]
  \;=\;
  \frac{\phi_0\bigl(\llbracket f \rrbracket_\sigma\bigr)}
       {\phi_0\bigl(\llbracket \mathbf{1}_\sigma \rrbracket_\sigma\bigr)}.
\end{equation}
Note that the denominator is positive by the assumption
$\phi_0(\langle\sigma\rangle_0)>0$, since
$\phi_0\bigl(\llbracket \mathbf{1}_\sigma \rrbracket_\sigma\bigr)
=q_\sigma(\mathbf{1}_\sigma)\cdot
\phi_0(\langle\sigma\rangle_0)$ and $q_\sigma(\mathbf{1}_\sigma)>0$ always.

Equation~\eqref{eq:ensemble} implies that the downward operator $\llbracket \cdot \rrbracket_\sigma$ maps $\mathcal{C}^\sigma$ to $\mathcal{C}$:

\begin{theorem}[Razborov~{\cite{razborov2007flag}}]\label{thm:downward-nonneg}
  For any flag type $\sigma$ and any $f \in \mathcal{A}^\sigma$ with $0 \leq_{\mathcal{H},\sigma} f$,
  \[
    0 \;\leq_{\mathcal{H}}\; \llbracket f \rrbracket_\sigma.
  \]
\end{theorem}
\begin{proof}[Proof sketch]
  Let $\phi_0 \in \poshom{\emptyset}$; we show $\phi_0(\llbracket f \rrbracket_\sigma) \geq 0$.
  If $\phi_0(\langle\sigma\rangle_0) = 0$, then $\phi_0$ assigns density zero
  to the underlying unlabeled graph of $\sigma$.  The underlying graph of
  every $\sigma$-flag contains an induced copy of $\sigma$, so the downward
  image of every $\sigma$-flag basis element also evaluates to zero under
  $\phi_0$.  By linearity, $\phi_0(\llbracket f \rrbracket_\sigma) = 0$.
  Otherwise, let $\phi^\sigma \sim \Ext_\sigma(\phi_0)$.
  Since $0 \leq_{\mathcal{H},\sigma} f$, every positive homomorphism on
  $\mathcal{A}^\sigma$ evaluates $f$ non-negatively, so
  $\mathbb{E}[\phi^\sigma(f)] \geq 0$.
  Rearranging Equation~\eqref{eq:ensemble}:
  \[
    \phi_0\bigl(\llbracket f \rrbracket_\sigma\bigr)
    \;=\;
    \phi_0\bigl(\llbracket \mathbf{1}_\sigma \rrbracket_\sigma\bigr)
    \cdot
    \mathbb{E}\bigl[\phi^\sigma(f)\bigr]
    \;\geq\; 0.
  \]
  Since $\phi_0$ was arbitrary, we have $0 \leq_{\mathcal{H}} \llbracket f \rrbracket_\sigma$.
\end{proof}

A flag algebra proof of $[F] \leq_{\mathcal{H}} c\cdot\mathbf{1}$ typically proceeds as follows:
first, one exhibits $f \in \mathcal{A}^\sigma$ with $0 \leq_{\mathcal{H},\sigma} f$ for some
nonempty flag type $\sigma$ (e.g., a square $f = g^2$) and applies
\Cref{thm:downward-nonneg} to obtain $0 \leq_{\mathcal{H}} \llbracket f \rrbracket_\sigma$;
then, one adds $\llbracket f \rrbracket_\sigma$ to $[F]$ or an expansion of $[F]$
at an appropriate size, and uses non-negativity of flag basis elements and
Equation~\eqref{eq:sum-to-one} to conclude $[F] \leq_{\mathcal{H}} c\cdot\mathbf{1}$.
This is the shape of proof our certificate-to-proof compiler produces, and
\Cref{sec:flagmatic-background} follows it through in full for Mantel's
theorem.

\section{Formalizing Flag Algebras in Lean~4}
\label{sec:abstract}

Informal flag algebra proofs are compact partly because their notation is
deliberately overloaded.  The same symbol may denote a concrete labeled graph,
its isomorphism class, a basis vector in a free vector space, an element of the
quotient algebra, or the real number obtained by applying a positive
homomorphism to that element.  On paper, these changes of viewpoint are usually
harmless and are left implicit.  In Lean, however, they create an immediate
formalization challenge: each role must be represented by a separate type, and
each passage between roles must be justified.  In particular, constructions
on concrete graphs must be invariant under isomorphism, expansion identities
must hold in the quotient, multiplication must be well defined on the quotient
algebra, and positive homomorphisms must preserve the algebraic operations.
Discharging these obligations, and proving in Lean the flag-algebra theorems
of \Cref{sec:background} that rest on them, is the work of the
\emph{specification layer}, presented in the following subsections.

One design choice pervades the development: how forbidden-subgraph assumptions
are represented.  Where
\Cref{sec:background} fixes \(\mathcal{H}\) first and builds flags and flag
algebras inside the class of \(\mathcal{H}\)-free graphs, our formalization
constructs the ambient flag algebras with no forbidden-family parameter and
imposes \(\mathcal{H}\)-freeness only where the semantic order is tested, by
restricting the positive homomorphisms it quantifies over.  Carrying a
forbidden family through the definitions would burden even flags, densities,
and multiplication with an extra parameter and side conditions; postponing it
lets us formalize the \(\mathcal{H}\)-independent theory once and reuse it for
every forbidden graph.  \Cref{sec:forbidden} defines the resulting ensemble
semantic order and its transfer theorem, and \Cref{sec:metatheory} compares it
with the order obtained by building the constraint in from the outset.

The Lean snippets below are lightly adapted for exposition.  We spell out
parameters that Lean would normally infer, so that a definition's types and
dependencies are visible without the surrounding source context; we gather
those shared by a group of definitions into a \lean{variable} block, which
applies to the definitions after it in the same subsection, leaving
each definition to show only its own arguments; and we occasionally replace
implementation-oriented names with clearer paper-facing names.  Following the
development itself, we write the empty flag type as \(\emptyt\), rather than
the \(\emptyset\) of \Cref{sec:background-flags}, to distinguish it from ordinary
empty sets.

\subsection{Flag Types and Flags}
\label{sec:flagtypes}

Before giving the Lean definitions of flags, we recall the small amount of
Mathlib notation used below.  In \Cref{sec:background}, we described a finite
graph as having a vertex set $V(G)$.  In Lean, the vertex set is represented
by a type: a term of type \lean{SimpleGraph V} is a simple graph whose
vertices are the elements of the type \lean{V}.  For example,
\lean{SimpleGraph (Fin n)} is a graph on the type \lean{Fin n}, whose elements
are the natural numbers \(0,\ldots,n-1\).  Finiteness of the vertex type is
not part of \lean{SimpleGraph} itself; it is supplied separately, by
assumptions such as \lean{Fintype V}.

Mathlib provides ${G\hookrightarrow_g H}$ for the type of graph embeddings
from $G$ to $H$.  Such an embedding is an injective vertex map that preserves
adjacency and non-adjacency, so its image is an induced copy of $G$ in $H$.  Mathlib also provides $G\simeq_g H$ for graph isomorphisms.
With this notation in place, the concrete pair $F=(M,\theta_F)$ from
\Cref{sec:background-flags} is represented by \lean{LabeledGraph}, while its
isomorphism class is represented by \lean{Flag}.

\begin{lstlisting}
abbrev FlagType (T : Type) := SimpleGraph T

structure LabeledGraph {T : Type} [Fintype T] ((*@$\sigma$@*) : FlagType T) (V : Type) where
  graph      : SimpleGraph V
  type_embed : (*@$\sigma$@*) (*@$\hookrightarrow$@*)g graph
\end{lstlisting}
Here \lean{FlagType T} is simply a name for \lean{SimpleGraph T}: a flag type
is represented by a simple graph on the label type \lean{T}.  The structure
\lean{LabeledGraph} packages the two pieces of a concrete $\sigma$-flag: the
field \lean{graph} stores the underlying graph $M$, and the field
\lean{type\_embed} stores the embedding
$\theta_F:\sigma\hookrightarrow_g M$.  In the signature, braces mark implicit
parameters that Lean normally infers from context, whereas square brackets
mark typeclass assumptions such as \lean{[Fintype T]}, which records that the
label type \lean{T} is finite.

We define \lean{LabeledGraphIso} as a structure representing isomorphisms
between two \lean{LabeledGraph} terms, and define \lean{Flag} by quotienting
\lean{LabeledGraph} by this isomorphism relation:
\begin{lstlisting}
structure LabeledGraphIso {T V V' : Type} [Fintype T]
    {(*@$\sigma$@*) : FlagType T} (G : LabeledGraph (*@$\sigma$@*) V) (G' : LabeledGraph (*@$\sigma$@*) V') where
  graph_iso     : G.graph (*@$\simeq$@*)g G'.graph
  type_preserve : graph_iso (*@$\circ$@*) G.type_embed = G'.type_embed

infixl:50 " (*@\textcolor{orange!80!black}{$\simeq$}@*)f " => LabeledGraphIso

instance labeledGraphSetoid {T : Type} [Fintype T] ((*@$\sigma$@*) : FlagType T) (V : Type)
    : Setoid (LabeledGraph (*@$\sigma$@*) V) where
  r     := fun (G G' : LabeledGraph (*@$\sigma$@*) V) => Nonempty (G (*@$\simeq$@*)f G')
  iseqv := ... -- proof that r defines an equivalence relation

def Flag {T : Type} [Fintype T] ((*@$\sigma$@*) : FlagType T) (V : Type) : Type :=
  Quotient (labeledGraphSetoid (*@$\sigma$@*) V)
\end{lstlisting}
The structure \lean{LabeledGraphIso} represents an isomorphism of concrete
$\sigma$-flags.
It consists of a graph isomorphism \lean{graph\_iso} between the underlying
graphs, together with a compatibility proof \lean{type\_preserve} that this
isomorphism sends the labeled copy of $\sigma$ in $G$ to the labeled copy in
$G'$, that is, $\alpha\circ\theta_G=\theta_{G'}$.  The notation
$G\simeq_f G'$ denotes this type of isomorphism data.

Lean's quotient type does not accept an arbitrary relation directly; it
requires a \emph{setoid}, which packages a relation with proofs that it is an
equivalence relation.  The instance \lean{labeledGraphSetoid} provides this
package for \lean{LabeledGraph}~$\sigma$~\lean{V}: its field \lean{r}
identifies two terms when a labeled-graph isomorphism between them exists, and
\lean{iseqv} holds the corresponding proofs.  Then
\lean{Flag}~$\sigma$~\lean{V} is the quotient of
\lean{LabeledGraph}~$\sigma$~\lean{V} by that setoid, so its elements are
precisely isomorphism classes of concrete $\sigma$-flags on the vertex type
\lean{V}.  Lean writes such classes with double brackets:
\(\llbracket G\rrbracket\) is the $\sigma$-flag represented by the concrete
labeled graph \(G\).

We can now represent $\mathcal{F}^\sigma_n$ in Lean by
\lean{Flag}~$\sigma$~\lean{(Fin}~$n$\lean{)}.  The type
\lean{FinFlag}~$\sigma$ collects $\sigma$-flags of all finite sizes and
represents $\mathcal{F}^\sigma = \bigcup_n \mathcal{F}^\sigma_n$.  It is
defined as a dependent sum type: each element is a dependent pair
$\hat F = \langle n, F\rangle$, where $n$ records the size and $F$ is the
$\sigma$-flag of that size.  In Lean, the two components are accessed as
$\hat F.1$ and $\hat F.2$.
\begin{lstlisting}
def FinFlag {k : (*@$\mathbb{N}$@*)} ((*@$\sigma$@*) : FlagType (Fin k)) : Type :=
  (*@$\Sigma$@*) (n : (*@$\mathbb{N}$@*)), Flag (*@$\sigma$@*) (Fin n)
\end{lstlisting}

Unlike the definitions in \Cref{sec:background-flags}, none of the definitions above
carries the \(\mathcal{H}\)-free condition: the underlying graphs are
arbitrary.  This is a deliberate interface choice.  The core quotient and
density machinery is developed once for all finite graphs, while
\(\mathcal{H}\)-freeness is imposed later by restricting the positive
homomorphisms used in the ensemble semantic order
(\Cref{sec:forbidden}).

\subsection{Subflag Densities}
\label{sec:formal-subflag-densities}

Formalizing the density function is not straightforward.  The density
$\den{F_1,\ldots,F_t}{G}$ from \Cref{sec:background-densities} should be a
function taking $t$ pattern $\sigma$-flags and a host $\sigma$-flag to a value in $[0,1]$, but
\lean{Flag} is a quotient type of \lean{LabeledGraph}.  The standard way to
define a function from a quotient type in Lean is to first define it on the
underlying type, prove that it respects the equivalence relation, and lift it
via \lean{Quotient.lift}.  That is, one must first define the density on
labeled graphs, and then prove that replacing any input by an isomorphic one
gives the same result.

We first introduce \lean{LabeledGraphList} as an abbreviation for a
(\lean{Fin t})-indexed tuple of labeled graphs.  Given such a tuple of
patterns and a host, \lean{labeledGraphListDensity} computes the density
directly as a ratio: the number of ways to choose pairwise-disjoint sets of
unlabeled host vertices, one per pattern, that induce the given patterns, over
the number of ways to choose such sets at all.  This is easier to work with in
Lean than the sampling formulation of \Cref{sec:background-densities}.
\begin{lstlisting}
abbrev LabeledGraphList {T : Type} [Fintype T]
    ((*@$\sigma$@*) : FlagType T) (t : (*@$\mathbb{N}$@*)) (Vl : Fin t (*@$\to$@*) Type) : Type :=
  (i : Fin t) (*@$\to$@*) LabeledGraph (*@$\sigma$@*) (Vl i)

noncomputable def labeledGraphListDensity
    {T W : Type} [Fintype T] [Fintype W] [DecidableEq W]
    {t : (*@$\mathbb{N}$@*)} {Vl : Fin t (*@$\to$@*) Type} [FintypeList Vl] [DecidableEqList Vl]
    {(*@$\sigma$@*) : FlagType T} (Hl : LabeledGraphList (*@$\sigma$@*) t Vl) (G : LabeledGraph (*@$\sigma$@*) W) : (*@$\mathbb{Q}$@*) :=
  let r_list (i : Fin t) := (Hl i).size - (*@$\sigma$@*).size
  labeledGraphListCount Hl G / multinomialCoefficient r_list (G.size - (*@$\sigma$@*).size)
\end{lstlisting}
Here a \lean{LabeledGraphList} is a function assigning to each index
\lean{i : Fin t} a labeled graph on its own vertex type \lean{Vl i}, so the
$t$ patterns need not share one ambient vertex type.  In the second
definition, \lean{Hl} is the tuple of patterns, \lean{G} the host, and the
local function \lean{r\_list} records how many unlabeled vertices each pattern
contributes: its value at index $i$ is \(|V(F_i)| - |V(\sigma)|\).  In the
ratio above, the numerator \lean{labeledGraphListCount} counts the tuples of
pairwise-disjoint unlabeled vertex sets of \lean{G} that induce \lean{Hl}, and
the denominator, which \lean{multinomialCoefficient} computes from
\lean{r\_list} and \(|V(G)| - |V(\sigma)|\), counts all such tuples.

The square-bracketed parameters record what finite data is available: for a
type \lean{X}, \lean{[Fintype X]} exhibits a finite list of all elements of
\lean{X}, and \lean{[DecidableEq X]} a procedure deciding equality on
\lean{X}.  Since the pattern tuple may use a different vertex type
\lean{Vl i} at each index, \lean{[FintypeList Vl]} and
\lean{[DecidableEqList Vl]} package the same data for every \lean{Vl i}.
Lean normally supplies all of these by typeclass inference.

The \lean{noncomputable} keyword signals that Lean cannot synthesize an
algorithm to evaluate this definition directly.  This poses no obstacle to
formalizing flag algebra itself, but becomes relevant when applying the algebra
to concrete extremal problems, a point we return to in \Cref{sec:reflection}.

To make \lean{labeledGraphListDensity} a function of $\sigma$-flags, we lift
its two inputs separately.  A \lean{Quotient.lift} on the host argument gives
\(\lean{LabeledGraphList}\to\lean{Flag}\to\mathbb{Q}\); a second one on the
pattern tuple gives
\(\lean{QuotLabeledGraphList}\to\lean{Flag}\to\mathbb{Q}\), where
\lean{QuotLabeledGraphList} is the quotient of concrete tuples by entrywise
isomorphism; and repackaging that input as a tuple of $\sigma$-flags gives the
final \(\lean{FlagList}\to\lean{Flag}\to\mathbb{Q}\).  The congruence proofs
for the two lifts, omitted below, show that the density is unchanged when the
host $\sigma$-flag or any pattern $\sigma$-flag is replaced by an isomorphic
representative.
\begin{lstlisting}
variable {T W : Type} [Fintype T] [Fintype W] [DecidableEq W]
         {t : (*@$\mathbb{N}$@*)} {Vl : Fin t (*@$\to$@*) Type} [FintypeList Vl] [DecidableEqList Vl]
         {(*@$\sigma$@*) : FlagType T}

def labeledGraphListEqv (Gl Gl' : LabeledGraphList (*@$\sigma$@*) t Vl) : Prop :=
  (*@$\forall$@*) (i : Fin t), Nonempty (Gl i (*@$\simeq$@*)f Gl' i)

def QuotLabeledGraphList ((*@$\sigma$@*) : FlagType T) (t : (*@$\mathbb{N}$@*)) (Vl : Fin t (*@$\to$@*) Type) : Type :=
  Quotient (labeledGraphListSetoid (*@$\sigma$@*) t Vl) -- the setoid of labeledGraphListEqv

abbrev FlagList ((*@$\sigma$@*) : FlagType T) (t : (*@$\mathbb{N}$@*)) (Vl : Fin t (*@$\to$@*) Type) : Type :=
  (*@$\forall$@*) (i : Fin t), Flag (*@$\sigma$@*) (Vl i)

noncomputable def labeledGraphListDensityLifted (Hl : LabeledGraphList (*@$\sigma$@*) t Vl)
    : Flag (*@$\sigma$@*) W (*@$\to$@*) (*@$\mathbb{Q}$@*) :=
  Quotient.lift (fun G => labeledGraphListDensity Hl G)
    (...) -- proof that host-graph equivalence preserves density

noncomputable def quotLabeledGraphListDensity
    : QuotLabeledGraphList (*@$\sigma$@*) t Vl (*@$\to$@*) Flag (*@$\sigma$@*) W (*@$\to$@*) (*@$\mathbb{Q}$@*) :=
  Quotient.lift labeledGraphListDensityLifted
    (...) -- proof that pattern-list equivalence preserves density

noncomputable def FlagList.coe (Fl : FlagList (*@$\sigma$@*) t Vl)
    : QuotLabeledGraphList (*@$\sigma$@*) t Vl :=
  (*@$\llbracket$@*)fun i => (Fl i).out(*@$\rrbracket$@*)

noncomputable def flagListDensity
    : FlagList (*@$\sigma$@*) t Vl (*@$\to$@*) Flag (*@$\sigma$@*) W (*@$\to$@*) (*@$\mathbb{Q}$@*) :=
  fun Fl => quotLabeledGraphListDensity (FlagList.coe Fl)
\end{lstlisting}
Here \lean{labeledGraphListEqv} is the entrywise isomorphism relation whose
setoid, formed as in \Cref{sec:flagtypes}, defines
\lean{QuotLabeledGraphList}.  The three density definitions carry out the
three steps above, the last one through \lean{FlagList.coe}, which chooses a
representative \lean{(Fl i).out} for each entry of \lean{Fl} and forms the
quotient class of the resulting concrete tuple.

The following \lean{flagDensity}$_1$ and \lean{flagDensity}$_2$ are
specializations of \lean{flagListDensity} to singleton and pair lists:
\begin{lstlisting}
variable {U U(*@$_1$@*) U(*@$_2$@*) : Type} [Fintype U] [DecidableEq U]
         [Fintype U(*@$_1$@*)] [DecidableEq U(*@$_1$@*)] [Fintype U(*@$_2$@*)] [DecidableEq U(*@$_2$@*)]

noncomputable def flagDensity(*@$_1$@*) (F : Flag (*@$\sigma$@*) U) (G : Flag (*@$\sigma$@*) W) : (*@$\mathbb{Q}$@*) :=
  flagListDensity [F](*@$^\mathtt{f}$@*) G

noncomputable def flagDensity(*@$_2$@*) (F(*@$_1$@*) : Flag (*@$\sigma$@*) U(*@$_1$@*)) (F(*@$_2$@*) : Flag (*@$\sigma$@*) U(*@$_2$@*)) (G : Flag (*@$\sigma$@*) W) : (*@$\mathbb{Q}$@*) :=
  flagListDensity [F(*@$_1$@*), F(*@$_2$@*)](*@$^\mathtt{f}$@*) G
\end{lstlisting}
Here the superscript \(\mathtt{f}\) marks a tuple constructor for
\lean{FlagList}, distinguishing it from ordinary list notation and from
quotient brackets.

Building on the definitions above, we prove the following chain rule
(\Cref{lem:chain-rule}) and several of its variants.
The statement is an elementary counting identity, but its formal proof is not:
the counts are defined on concrete labeled graphs, and the identity between
them comes from an explicit bijection between the choices the two sampling
procedures make (\Cref{sec:proof-style}).
\begin{lstlisting}
theorem flagDensity_eq_sum_density_prods
    {(*@$\ell_0$@*) (*@$\ell_1$@*) (*@$\ell$@*) (*@$\ell'$@*) : (*@$\mathbb{N}$@*)} {(*@$\sigma$@*) : FlagType (Fin (*@$\ell_0$@*))} (F(*@$_1$@*) : Flag (*@$\sigma$@*) (Fin (*@$\ell_1$@*))) (G : Flag (*@$\sigma$@*) (Fin (*@$\ell$@*)))
    (h(*@$\ell_1$@*) : (*@$\ell_0$@*) (*@$\leq$@*) (*@$\ell_1$@*)) (h(*@$\ell'$@*) : (*@$\ell_1$@*) (*@$\leq$@*) (*@$\ell'$@*)) (h(*@$\ell$@*) : (*@$\ell'$@*) (*@$\leq$@*) (*@$\ell$@*))
    : flagDensity(*@$_1$@*) F(*@$_1$@*) G = (*@$\sum$@*) G' : Flag (*@$\sigma$@*) (Fin (*@$\ell'$@*)), flagDensity(*@$_1$@*) F(*@$_1$@*) G' * flagDensity(*@$_1$@*) G' G := ...
\end{lstlisting}

\subsection{The Flag Algebra}
\label{sec:formal-flag-algebra}

To define the flag algebra, we first construct the free real vector space on
$\mathcal{F}^\sigma$.  Each element of this vector space is represented as a
finitely supported function from $\mathcal{F}^\sigma$ to $\mathbb{R}$, using
Mathlib's \lean{Finsupp} type.  We call the resulting type
\lean{FlagVector}~$\sigma$; the constructor \lean{basisVector} produces the
basis vector for each $\sigma$-flag:
\begin{lstlisting}
variable {k : (*@$\mathbb{N}$@*)} {(*@$\sigma$@*) : FlagType (Fin k)}

abbrev FlagVector ((*@$\sigma$@*) : FlagType (Fin k)) : Type :=
  FinFlag (*@$\sigma$@*) (*@$\to_0$@*) (*@$\mathbb{R}$@*)

def basisVector (F : FinFlag (*@$\sigma$@*)) : FlagVector (*@$\sigma$@*) :=
  Finsupp.single F 1
\end{lstlisting}
Here $\to_0$ denotes the type of finitely supported functions, and \lean{Finsupp.single F 1} constructs the basis vector $\mathbf e_F$ for $F \in \mathcal{F}^\sigma$: the function taking value $1$ at $F$ and $0$ elsewhere.

Next, we define the zero space $\mathcal{Z}^\sigma$.  The auxiliary function
\lean{flagExpansion F}~$\ell$ is the Lean version of
$\sum_{G \in \mathcal{F}^\sigma_\ell} \den{F}{G}\cdot \mathbf e_G$, the expansion of
$F$ at size $\ell$, with the black dot below for the scalar multiplication
written $\cdot$ here.  We encode the generating element of
Equation~\eqref{eq:zero-element} as \lean{zeroElement}, \lean{zeroSet} $\sigma$
collects all elements of this form, and \lean{ZeroSpace} $\sigma$ is the
linear subspace they span, represented in Lean by a \lean{Submodule}:
\begin{lstlisting}
noncomputable def flagExpansion (F : FinFlag (*@$\sigma$@*)) ((*@$\ell$@*) : (*@$\mathbb{N}$@*)) : FlagVector (*@$\sigma$@*) :=
  (*@$\sum$@*) G : Flag (*@$\sigma$@*) (Fin (*@$\ell$@*)), flagDensity(*@$_1$@*) F.2 G (*@$\leansmul$@*) basisVector (*@$\langle$@*)(*@$\ell$@*), G(*@$\rangle$@*)

noncomputable def zeroElement (F : FinFlag (*@$\sigma$@*)) ((*@$\ell$@*) : (*@$\mathbb{N}$@*)) : FlagVector (*@$\sigma$@*) :=
  basisVector F - flagExpansion F (*@$\ell$@*)

noncomputable def zeroSet ((*@$\sigma$@*) : FlagType (Fin k)) : Set (FlagVector (*@$\sigma$@*)) :=
  {z | (*@$\exists$@*) (F : FinFlag (*@$\sigma$@*)) ((*@$\ell$@*) : (*@$\mathbb{N}$@*)), F.1 (*@$\leq$@*) (*@$\ell$@*) (*@$\wedge$@*) z = zeroElement F (*@$\ell$@*)}

noncomputable def ZeroSpace ((*@$\sigma$@*) : FlagType (Fin k)) : Submodule (*@$\mathbb{R}$@*) (FlagVector (*@$\sigma$@*)) :=
  Submodule.span (*@$\mathbb{R}$@*) (zeroSet (*@$\sigma$@*))
\end{lstlisting}
Following the same pattern as the definition of \lean{Flag} from
\lean{LabeledGraph} in \Cref{sec:flagtypes}, we now define the flag algebra
$\mathcal{A}^\sigma$ as a quotient of \lean{FlagVector} $\sigma$, by the
relation \lean{flagVectorEqv} that two flag vectors differ by an element of
\lean{ZeroSpace} $\sigma$.
\begin{lstlisting}
def flagVectorEqv (f g : FlagVector (*@$\sigma$@*)) : Prop :=
  f - g (*@$\in$@*) ZeroSpace (*@$\sigma$@*)

abbrev FlagAlgebra ((*@$\sigma$@*) : FlagType (Fin k)) : Type :=
  Quotient (flagVectorSetoid (*@$\sigma$@*)) -- the setoid of flagVectorEqv
\end{lstlisting}
Because \lean{ZeroSpace} $\sigma$ is spanned by the differences
\lean{zeroElement F}~$\ell$, this quotient imposes a family of expansion
identities: for every $\sigma$-flag \lean{F : FinFlag}~$\sigma$ and every size
$\ell \geq F.1$, the equivalence class of \lean{basisVector F} is identified
with that of its expansion \lean{flagExpansion F}~$\ell$.  As
\lean{FlagAlgebra}~$\sigma$ is again a quotient type, if $f$ is a
\lean{FlagVector}~$\sigma$ then $\llbracket f\rrbracket$ is its equivalence
class in this quotient.

Formalizing multiplication in $\mathcal{A}^\sigma$ also requires several
steps.  The term \lean{flagMulWithSize} defines
the product of two $\sigma$-flags at a chosen output size $\ell$ as in
Equation~\eqref{eq:mul}, and \lean{flagMul} specializes it to the minimal
admissible size $\ell = F.1 + F'.1 - k$.  Then \lean{bilinearExtension}
extends \lean{flagMul} linearly in each argument to give a product on
\lean{FlagVector} $\sigma$, an operation on formal sums for which none of the
algebra laws yet hold, and \lean{Quotient.map$_2$} lifts that product to
\lean{FlagAlgebra} $\sigma$.
\begin{lstlisting}
noncomputable def flagMulWithSize (F F' : FinFlag (*@$\sigma$@*)) ((*@$\ell$@*) : (*@$\mathbb{N}$@*)) : FlagVector (*@$\sigma$@*) :=
  (*@$\sum$@*) G : Flag (*@$\sigma$@*) (Fin (*@$\ell$@*)), flagDensity(*@$_2$@*) F.2 F'.2 G (*@$\leansmul$@*) basisVector (*@$\langle$@*)(*@$\ell$@*), G(*@$\rangle$@*)

noncomputable def flagMul (F F' : FinFlag (*@$\sigma$@*)) : FlagVector (*@$\sigma$@*) :=
  flagMulWithSize F F' (F.1 + F'.1 - k)

noncomputable instance ((*@$\sigma$@*) : FlagType (Fin k)) : Mul (FlagVector (*@$\sigma$@*)) where
  mul := bilinearExtension flagMul

noncomputable instance ((*@$\sigma$@*) : FlagType (Fin k)) : Mul (FlagAlgebra (*@$\sigma$@*)) where
  mul := Quotient.map(*@$_2$@*) ((*@$\cdot$@*) * (*@$\cdot$@*))
    (...) -- proof that the product respects flagVectorEqv in each argument
\end{lstlisting}
Here \lean{Quotient.map$_2$} lifts a binary operation to a quotient type given
a proof that the operation respects the equivalence relation.  The expression
\lean{($\cdot$ * $\cdot$)} is not a new operation: its two arguments are flag
vectors, so Lean resolves \lean{*} using the immediately preceding
\lean{Mul (FlagVector}~$\sigma$\lean{)} instance, whose multiplication is
\lean{bilinearExtension flagMul}.  Its proof obligation is that replacing
either factor by an equivalent flag vector gives an equivalent product, which
settles the representative half of the independence claimed for
Equation~\eqref{eq:mul}.  A separate theorem settles the other half, that any
two admissible output sizes give equivalent flag vectors.
\begin{lstlisting}
theorem flagMulWithSize_indep_on_size {F(*@$_1$@*) F(*@$_2$@*) : FinFlag (*@$\sigma$@*)} {(*@$\ell_1$@*) (*@$\ell_2$@*) : (*@$\mathbb{N}$@*)}
    (h(*@$\ell_1$@*) : F(*@$_1$@*).1 + F(*@$_2$@*).1 (*@$\leq$@*) (*@$\ell_1$@*) + k) (h(*@$\ell_2$@*) : F(*@$_1$@*).1 + F(*@$_2$@*).1 (*@$\leq$@*) (*@$\ell_2$@*) + k)
    : flagVectorEqv (flagMulWithSize F(*@$_1$@*) F(*@$_2$@*) (*@$\ell_1$@*)) (flagMulWithSize F(*@$_1$@*) F(*@$_2$@*) (*@$\ell_2$@*)) := ...
\end{lstlisting}

With these definitions in place, we prove that \lean{FlagAlgebra} $\sigma$ is
a commutative algebra over the real numbers, completing the formalization of
the algebraic structure of $\mathcal{A}^\sigma$.

\subsection{Positive Homomorphisms and Semantic Order}

This subsection describes how we encode the semantic side of flag algebras in
Lean: sequences of finite $\sigma$-flags, their limiting density functions, positive
homomorphisms, and the order induced by testing against all positive
homomorphisms.  The convergence theorem from
Theorem~\ref{thm:convergent-hom} is the main mathematical bridge between finite
$\sigma$-flag sequences and positive homomorphisms, but the formalization first has to
make each of these surrounding notions explicit.

A sequence of $\sigma$-flags indexed by $\mathbb{N}$ is represented by
\lean{FlagSeq}~$\sigma$, and \lean{ConvergesTo s a} says that such a sequence
$s$ is convergent in the sense of \Cref{sec:background-semantic-order}, with
limit $a$.
\begin{lstlisting}
variable {k : (*@$\mathbb{N}$@*)} {(*@$\sigma$@*) : FlagType (Fin k)}

abbrev FlagSeq ((*@$\sigma$@*) : FlagType (Fin k)) :=
  (*@$\mathbb{N}$@*) (*@$\to$@*) FinFlag (*@$\sigma$@*)

def ConvergesTo (s : FlagSeq (*@$\sigma$@*)) (a : FinFlag (*@$\sigma$@*) (*@$\to$@*) (*@$\mathbb{R}$@*)) : Prop :=
  Increases s (*@$\wedge$@*) Tendsto (flagDensitySeq s) atTop ((*@$\mathcal{N}$@*) a)
\end{lstlisting}
Here the candidate limit \(a\) is a real-valued function on finite
\(\sigma\)-flags.  The first conjunct, \lean{Increases s}, asserts that the
underlying sizes of the \(\sigma\)-flags \(s(0),s(1),\ldots\) strictly increase.  The
second conjunct uses Lean's filter notation for convergence:
\lean{atTop} is the limit \(n\to\infty\) on \(\mathbb{N}\), and
\lean{$\mathcal{N}$\,\,a} is the neighborhood filter of the function~\(a\).  The term
\lean{flagDensitySeq s} is the sequence of density functions associated to
\(s\); its \(n\)-th value sends a finite \(\sigma\)-flag \(F\) to
\(\den{F}{s(n)}\).  Thus, the \lean{Tendsto} line is the Lean statement that,
for every finite \(\sigma\)-flag~\(F\),
\[
  \lim_{n\to\infty}\den{F}{s(n)}=a(F).
\]

The Lean definition of a positive homomorphism follows the mathematical one of
\Cref{sec:background-semantic-order} almost literally: a map
\(\phi:\mathcal{A}^\sigma\to\mathbb{R}\) preserving the algebraic structure and
non-negative on every flag basis element.
\begin{lstlisting}
abbrev Hom ((*@$\sigma$@*) : FlagType (Fin k)) :=
  FlagAlgebra (*@$\sigma$@*) (*@$\to_a$@*)[(*@$\mathbb{R}$@*)] (*@$\mathbb{R}$@*)

def PositiveHom ((*@$\sigma$@*) : FlagType (Fin k)) : Type :=
  { (*@$\phi$@*) : Hom (*@$\sigma$@*) // (*@$\forall$@*) F : FinFlag (*@$\sigma$@*), (*@$\phi$@*) (*@$\llbracket$@*)basisVector F(*@$\rrbracket$@*) (*@$\geq$@*) 0 }
\end{lstlisting}
The arrow \lean{$\to_a$[$\mathbb{R}$]} is Lean's notation for the type of
functions preserving addition, multiplication, multiplication by real scalars,
and the unit.  The second definition is Lean's subtype construction, whose
separator \lean{//} precedes the defining predicate; mathematically we read
\lean{PositiveHom}~$\sigma$ as the set
\[
  \{\phi\in \operatorname{Hom}(\mathcal{A}^\sigma,\mathbb{R})
    \mid \phi([F])\ge 0
    \text{ for every }F\in\mathcal{F}^{\sigma}\}.
\]

Using these definitions, we formalize Theorem~\ref{thm:convergent-hom} as two
separate theorems.\footnote{For readability, the displayed statements give the
limit profile the ambient function type
\lean{FinFlag $\sigma$ $\to$ $\mathbb{R}$}.  The Lean implementation instead
packages such profiles in the auxiliary subtype
\lean{FlagDensitySpace}~\(\sigma\), which records that every value lies in
\([0,1]\), and packages the profile induced by a positive homomorphism as
\lean{$\phi$.coe}.  Accordingly, the implemented conclusions are written more
compactly as \lean{$\phi$.coe = a} and
\lean{ConvergesTo s $\phi$.coe}.  The formulations are equivalent: every flag
density lies in \([0,1]\), so the pointwise limit of a convergent flag sequence
does as well.}
\begin{lstlisting}
(*@\textcolor{gray}{\itshape -- Theorem~\ref*{thm:convergent-hom} (a)}@*)
theorem flagSeq_limit_mem_positiveHom (s : FlagSeq (*@$\sigma$@*))
    {a : FinFlag (*@$\sigma$@*) (*@$\to$@*) (*@$\mathbb{R}$@*)} (hs_conv : ConvergesTo s a)
    : (*@$\exists$@*) ((*@$\phi$@*) : PositiveHom (*@$\sigma$@*)), (*@$\forall$@*) (F : FinFlag (*@$\sigma$@*)), (*@$\phi$@*) (*@$\llbracket$@*)basisVector F(*@$\rrbracket$@*) = a F := ...

(*@\textcolor{gray}{\itshape -- Theorem~\ref*{thm:convergent-hom} (b)}@*)
theorem positiveHom_as_flagSeq_limit ((*@$\phi$@*) : PositiveHom (*@$\sigma$@*))
    : (*@$\exists$@*) (s : FlagSeq (*@$\sigma$@*)), ConvergesTo s (fun F => (*@$\phi$@*) (*@$\llbracket$@*)basisVector F(*@$\rrbracket$@*)) := ...
\end{lstlisting}
Part~(a) assumes \lean{ConvergesTo\,\,\,s\,\,\,a} and produces a positive
homomorphism \(\phi\) satisfying
\(\phi(\llbracket\lean{basisVector F}\rrbracket)=a(F)\) for every finite
\(\sigma\)-flag \(F\).  Thus, the limiting density profile of every convergent
\(\sigma\)-flag sequence is induced by a positive homomorphism.  Conversely,
part~(b) starts with a positive homomorphism \(\phi\) and produces a
\(\sigma\)-flag sequence whose limiting density profile is the one \(\phi\)
induces.
Both directions require nontrivial formalization work.  In part (a), one must
show that the limit \(a\) of a \(\sigma\)-flag sequence \(s\) satisfies all algebraic
properties of a positive homomorphism, including details that informal proofs
leave implicit.  Part~(b) goes the other way, and its proof is probabilistic:
a sequence of \(\sigma\)-flags sampled using the values of \(\phi\) as weights
almost surely has \(\phi\) as its limiting density profile.  Formalizing it
means constructing the product measure on flag sequences, proving that the
coordinatewise convergence event has probability one, and selecting a
realization from that event.

We also define the semantic non-negativity cone $\mathcal{C}^\sigma$
and the order it induces on $\mathcal{A}^\sigma$.
Here, 
\lean{semanticCone}~$\sigma$ is the set of elements in
\lean{FlagAlgebra}~$\sigma$ that every positive homomorphism
evaluates non-negatively.  The induced order is installed as a Lean
\lean{LE} instance, with $f \leq g$ defined to mean
$g - f \in \mathcal{C}^\sigma$.
\begin{lstlisting}
def semanticCone ((*@$\sigma$@*) : FlagType (Fin k)) : Set (FlagAlgebra (*@$\sigma$@*)) :=
  { f | (*@$\forall$@*) (*@$\phi$@*) : PositiveHom (*@$\sigma$@*), (*@$\phi$@*) f (*@$\geq$@*) 0 }

instance ((*@$\sigma$@*) : FlagType (Fin k)) : LE (FlagAlgebra (*@$\sigma$@*)) where
  le := fun f g => g - f (*@$\in$@*) semanticCone (*@$\sigma$@*)
\end{lstlisting}

\subsection{The Downward Operator}
\label{sec:formal-downward}

The downward operator from \Cref{sec:background-downward} forgets the
labels of a typed flag-algebra expression while keeping the correct
normalizing factor.  On a single $\sigma$-flag $F$, the intended formula is
\[
  \llbracket [F] \rrbracket_\sigma = q_\sigma(F)\cdot [F|_\emptyset].
\]
The displayed formula specifies the downward operator on flag basis elements.
To implement it, we first define its two ingredients, the normalizing factor
$q_\sigma(F)$ and the unlabeled flag $F|_\emptyset$, on concrete
$\sigma$-typed labeled graphs and lift both to flags.  This gives the downward
image of an individual $\sigma$-flag.  We then extend that operation linearly
to flag vectors and lift it once more, to $\mathcal{A}^\sigma$, giving a
well-defined operator
\(\mathcal{A}^\sigma\longrightarrow\mathcal{A}^\emptyset\).

For a concrete labeled graph $G$ on $n$ vertices, the normalizing factor is a
ratio: the numerator \lean{isomorphismCount G} counts the injections
$[k]\hookrightarrow V(G)$ that produce a typed labeled graph isomorphic
to~$G$, and the denominator $n!/(n-k)!$ counts all possible injections.  The
other ingredient, $F|_\emptyset$, is computed by
\lean{unlabel\_labeledGraph}, which forgets the type embedding.  Lifting these
two operations to flags takes the \lean{Quotient.lift} pattern of
\Cref{sec:formal-subflag-densities}, so each needs a proof of invariance under
isomorphism: isomorphic representatives admit the same label
placements, and their unlabelings are isomorphic as $\emptyt$-flags.  The two
then combine in \lean{downwardFlag}, the downward image of an individual
$\sigma$-flag.
\begin{lstlisting}
variable {k n : (*@$\mathbb{N}$@*)} {V : Type} {(*@$\sigma$@*) : FlagType (Fin k)}

noncomputable def downwardNormalizingFactor : Flag (*@$\sigma$@*) (Fin n) (*@$\to$@*) (*@$\mathbb{Q}$@*) :=
  Quotient.lift (fun G => isomorphismCount G / (n.factorial / (n - k).factorial)) (...)

noncomputable def unlabel : Flag (*@$\sigma$@*) V (*@$\to$@*) Flag (*@$\emptyt$@*) V :=
  Quotient.lift (fun G => (*@$\llbracket$@*)unlabel_labeledGraph G(*@$\rrbracket$@*)) (...)

noncomputable def downwardFlag (F : Flag (*@$\sigma$@*) (Fin n)) : FlagVector (*@$\emptyt$@*) :=
  downwardNormalizingFactor F (*@$\leansmul$@*) basisVector (*@$\langle$@*)n, unlabel F(*@$\rangle$@*)
\end{lstlisting}

We next lift this construction to the algebra.  The map
\lean{downwardFlagVector} extends \lean{downwardFlag} linearly to
$\sigma$-typed flag vectors, and \lean{downwardFlagVectorQuot} takes the
zero-space quotient of its output, so its values lie in
\lean{FlagAlgebra}~$\emptyt$ while its input is still a flag vector.
Quotienting the input as well requires the respect lemma below: two
$\sigma$-typed flag vectors differing by a zero-space element have downward
images representing the same element of the empty-type algebra.  This is the
compatibility condition needed for the final quotient lift.
\begin{lstlisting}
noncomputable def downwardFlagVector : FlagVector (*@$\sigma$@*) (*@$\to$@*) FlagVector (*@$\emptyt$@*) :=
  linearExtension (fun F : FinFlag (*@$\sigma$@*) => downwardFlag F.2)

noncomputable def downwardFlagVectorQuot : FlagVector (*@$\sigma$@*) (*@$\to$@*) FlagAlgebra (*@$\emptyt$@*) :=
  fun f : FlagVector (*@$\sigma$@*) => (*@$\llbracket$@*)downwardFlagVector f(*@$\rrbracket$@*)

lemma downwardFlagVectorQuot_respect_eqv {f f' : FlagVector (*@$\sigma$@*)}
    (h : f - f' (*@$\in$@*) ZeroSpace (*@$\sigma$@*))
    : downwardFlagVectorQuot f = downwardFlagVectorQuot f' := ...

noncomputable def downward : FlagAlgebra (*@$\sigma$@*) (*@$\to$@*) FlagAlgebra (*@$\emptyt$@*) :=
  Quotient.lift (fun f : FlagVector (*@$\sigma$@*) => downwardFlagVectorQuot f) (...)

notation "(*@\textcolor{orange!80!black}{$\llbracket$}@*)" f "(*@\textcolor{orange!80!black}{$\rrbracket_0$}@*)" => (downward f)
\end{lstlisting}
The last definition is now a genuine map
\(\lean{FlagAlgebra}~\sigma\to\lean{FlagAlgebra}~\emptyt\), independent of all
representative choices.  The final line introduces the Lean notation
\(\llbracket f\rrbracket_0\) for \lean{downward f}; the subscript $0$
emphasizes that the result lies in the empty-type algebra.  In mathematical
prose we use \(\llbracket f\rrbracket_\sigma\), following
\Cref{sec:background-downward}, where the subscript instead records the type
being forgotten.

Having constructed the algebraic map, we turn to its semantic meaning: a
formal counterpart of the random-extension identity of
Equation~\eqref{eq:ensemble},
\[
  \mathbb{E}\bigl[\phi^\sigma(f)\bigr]
  \;=\;
  \frac{\phi_0\bigl(\llbracket f \rrbracket_\sigma\bigr)}
       {\phi_0\bigl(\llbracket \mathbf{1}_\sigma \rrbracket_\sigma\bigr)}.
\]
Stating it in Lean requires a measurable space for the measure to live on,
and the implementation obtains one by identifying a positive homomorphism
with its density profile: the sample space \lean{PositiveHomSpace}~$\sigma$ is
the measurable subspace of \([0,1]^{\mathcal{F}^\sigma}\) consisting of the
profiles induced by positive homomorphisms.  The formal counterpart is then
the following existence statement.
\begin{lstlisting}
theorem (*@\texttt{exists\_probMeasure\_extend\_emptyType\_positiveHom}@*)
    {(*@$\phi_0$@*) : PositiveHom (*@$\emptyt$@*)} (h(*@$\sigma$@*) : (*@$\phi_0$@*) ((*@$\langle\sigma\rangle_0$@*)) > 0)
    : (*@$\exists$@*) (*@$\mathbb{P}$@*) : ProbabilityMeasure (PositiveHomSpace (*@$\sigma$@*)), (*@$\forall$@*) f : FlagAlgebra (*@$\sigma$@*),
        (*@$\int$@*) ((*@$\phi$@*) : PositiveHomSpace (*@$\sigma$@*)), PositiveHomSpace.toPosHom (*@$\phi$@*) f (*@$\partial$@*)(*@$\mathbb{P}$@*)
        = (*@$\phi_0$@*) (*@$\llbracket$@*)f(*@$\rrbracket_0$@*) / (*@$\phi_0$@*) (*@$\llbracket$@*)(1 : FlagAlgebra (*@$\sigma$@*))(*@$\rrbracket_0$@*)
  := ...
\end{lstlisting}
Mathlib writes \lean{$\int$ x, g x $\partial\mu$} for the integral of
$x\mapsto g(x)$ with respect to $\mu$; here $\partial\mu$ specifies the
measure and is not a derivative.  In the displayed theorem, $\phi$ ranges over
\lean{PositiveHomSpace}~$\sigma$, the integrand
\lean{PositiveHomSpace.toPosHom $\phi$ f} evaluates at $f$ the homomorphism
that $\phi$ induces, and the measure is $\mathbb{P}$.  Because
$\mathbb{P}$ is a probability measure, the left-hand side is the expectation
$\mathbb{E}_{\phi\sim \mathbb{P}}[\phi(f)]$.  The theorem therefore says
exactly that $\mathbb{P}$ satisfies the identity above, so $\mathbb{P}$
realizes the random-extension measure $\Ext_\sigma(\phi_0)$
of \Cref{sec:background-downward}.

The Lean counterpart of \Cref{thm:downward-nonneg} now follows from the
random-extension identity:
\begin{lstlisting}
theorem downward_preserve_semanticCone (f : FlagAlgebra (*@$\sigma$@*))
    (hf : f (*@$\in$@*) semanticCone (*@$\sigma$@*)) : (*@$\llbracket$@*)f(*@$\rrbracket_0$@*) (*@$\in$@*) semanticCone (*@$\emptyt$@*) := ...
\end{lstlisting}
Every typed inequality used by the certificate proofs of
\Cref{sec:flagmatic} reaches the unlabeled algebra through this theorem.

\subsection{Forbidden Graphs and the Ensemble Semantic Order}
\label{sec:forbidden}

In \Cref{sec:background}, a finite forbidden family \(\mathcal{H}\) is fixed
first, and the flag types, flags, algebras, and positive homomorphisms are all
built inside the \(\mathcal{H}\)-free world.  In that setup, an inequality
\(f\leq_{\mathcal{H},\sigma} g\) is already an inequality in the constrained
\(\sigma\)-typed algebra: the forbidden configurations have been removed
before the semantic order is tested.

Our Lean formalization deliberately separates building the algebra from
imposing the constraint.  The core development constructs the ambient flag
algebra \(\mathcal{A}^\sigma\) of all finite simple graphs, with no forbidden
family as a parameter.  A forbidden graph enters only when we compare two
ambient algebra elements \(f,g\in\mathcal{A}^\sigma\), and even then it defines
no separate flag-algebra universe: the forbidden graph restricts the
empty-type positive homomorphisms from which the random extensions of
\Cref{sec:background-downward} are formed.  The whole construction of the
preceding subsections is therefore reused unchanged across extremal problems.

The order used for this a posteriori condition is the
\emph{ensemble semantic order}.  We state it for one forbidden graph \(H\); a
family would impose the zero-density condition below for each member.

\begin{definition}[Ensemble semantic order]\label{def:ensemble-semantic-order}
Fix a forbidden graph \(H\), a flag type \(\sigma\), and elements
\(f,g\in\mathcal{A}^\sigma\) in the ambient flag algebra.  We write
\(f\leq^{\mathrm{ens}}_{H,\sigma} g\) if, for every empty-type positive
homomorphism \(\phi_0\) such that, viewing graphs as
empty-type flag-algebra elements,
\[
  \phi_0([F])=0
  \quad\text{for every graph \(F\) containing \(H\) as a subgraph}
  \qquad\text{and}\qquad
  \phi_0(\langle\sigma\rangle_0)>0,
\]
the random extension \(\phi\sim\Ext_\sigma(\phi_0)\) satisfies
\[
  \mathbb{P}\bigl[\phi(f)\leq\phi(g)\bigr]=1.
\]
\end{definition}

The first condition restricts \(\phi_0\) to the \(H\)-free unlabeled limits: it
assigns density zero to every graph containing \(H\).  The second condition,
\(\phi_0(\langle\sigma\rangle_0)>0\), says that the underlying unlabeled graph of
\(\sigma\) has positive limiting density along any sequence of finite
graphs converging to \(\phi_0\).  This is precisely the hypothesis under which the random-extension
measure \(\Ext_\sigma(\phi_0)\) is defined.  The final display asks that
\(\phi(f)\leq\phi(g)\) hold with probability one under this measure.

\begin{remark}[Relation to built-in \(H\)-free semantics]\label{rem:builtin-hfree}
The built-in approach of \Cref{sec:background} and the ensemble approach agree
at the empty type, where no labels are chosen, so the distinction does not
affect the final unlabeled Tur\'an-density bounds.  For nonempty types,
however, they can differ: the built-in approach may fix the labels at an
exceptional vertex configuration, which uniform random sampling makes
invisible to the ensemble comparison.
\Cref{sec:metatheory} gives a precise criterion for when the two approaches
agree and an example showing that they need not.
\end{remark}

The Lean predicate \lean{forbidLE} is the formalization of
\Cref{def:ensemble-semantic-order}:
\begin{lstlisting}
variable {k m : (*@$\mathbb{N}$@*)} {(*@$\sigma$@*) : FlagType (Fin k)}

def forbidLEWith (C : ForbidCondition) (f g : FlagAlgebra (*@$\sigma$@*)) : Prop :=
  (*@$\forall\;$@*)((*@$\phi_0$@*) : PositiveHom (*@$\emptyt$@*)),
    (*@$\phi_0$@*) (*@$\langle\sigma\rangle_0$@*) > 0
    -> (*@$\,$@*) C (*@$\phi_0$@*)
    -> (*@$\,$@*) (*@$\mathbb{P}$@*)[(*@$\phi_0$@*)] {(*@$\phi$@*) : PositiveHomSpace (*@$\sigma$@*) | (*@$\phi$@*) f (*@$\leq\;$@*)(*@$\phi$@*) g} = 1

def forbidLE (H : SimpleGraph (Fin m)) (f g : FlagAlgebra (*@$\sigma$@*)) : Prop :=
  forbidLEWith (forbiddenCondition H) f g
\end{lstlisting}
Read \lean{forbidLE} first.  Its explicit argument \lean{H} is the forbidden
graph, and \lean{f} and \lean{g} are the two ambient \(\sigma\)-typed
flag-algebra elements being compared.  It delegates to \lean{forbidLEWith}, the
condition-generic version parameterized by a \lean{ForbidCondition}, a
predicate selecting which empty-type positive homomorphisms are allowed.  Here
that predicate is \lean{forbiddenCondition H}, the \(H\)-freeness condition of
\Cref{def:ensemble-semantic-order}.

In the body of \lean{forbidLEWith}, the hypothesis \lean{C $\phi_0$} imposes
that condition and the conclusion is the probability-one requirement of the
same definition.  Two points are specific to the formalization: the event is
carved out of \lean{PositiveHomSpace $\sigma$}, the measurable space of density
profiles of \Cref{sec:formal-downward}, each of which determines a positive
homomorphism on \(\mathcal{A}^\sigma\); and the measure
\lean{$\mathbb{P}$[$\phi_0$]} is the random-extension measure
\(\Ext_\sigma(\phi_0)\) on that space.

What remains is to connect this semantic predicate to the corresponding
asymptotic bound on induced densities, the Lean counterpart of
\Cref{lem:semantic-bound} with its hypothesis stated in the ensemble order
rather than the built-in order of \Cref{sec:background-semantic-order}.  In the
theorem below, \(H\) is the forbidden graph, \(F\) is the target graph, and
\lean{F.toFlagAlgebra} converts \(F\) into the element of
\lean{FlagAlgebra $\emptyt$} whose evaluation is the induced \(F\)-density.
\begin{lstlisting}
theorem generalizedTuranDensity_le_of_forbidLE
    {n m : (*@$\mathbb{N}$@*)} {H : SimpleGraph (Fin n)}
    {F : SimpleGraph (Fin m)} {c : (*@$\mathbb{R}$@*)} (hc : 0 (*@$\leq$@*) c)
    (h : forbidLE H F.toFlagAlgebra (c (*@$\leansmul$@*) (1 : FlagAlgebra (*@$\emptyt$@*))))
    : generalizedTuranDensity H F (*@$\leq$@*) c := ...
\end{lstlisting}
The hypothesis \lean{h} is the empty-type ensemble inequality
\([F]\leq^{\mathrm{ens}}_{H,\emptyset}c\cdot\mathbf{1}\), and the conclusion is
the Tur\'an-density bound \(\tdensity{F}{H}\leq c\).  Each later application
proves the hypothesis and calls this theorem for the final conversion.

The Tur\'an density \(\tdensity{F}{H}\) is represented in Lean by
\lean{generalizedTuranDensity H F}, whose definition applies Mathlib's
\lean{limUnder} to the normalized extremal numbers
\(\mathrm{ex}(n,\,F;\,H)/\binom{n}{|V(F)|}\) of
\Cref{sec:background-semantic-order}.  That operation returns a real number
even before the sequence has been shown to converge, in which case the value
would have no asymptotic meaning.  We therefore prove
\lean{tendsto\_\allowbreak generalizedTuranDensity}: once \(n\geq |V(F)|\) the
sequence is non-increasing and bounded below by zero, so it converges, and its
limit is \lean{generalizedTuranDensity H F}.

\section{The Reflection Layer}
\label{sec:reflection}

The finite calculations a formal proof depends on can be discharged in two
ways.  One is to proceed as a
mathematician usually would: when a finite claim is needed, write a proof of
that claim directly.  In the flag-algebra setting, this would mean proving, one
by one, that two particular flags are isomorphic, that a certain density is a
given rational number, or that a product of flag-algebra basis elements expands into a specified
linear combination.  This is possible in principle, but it is the wrong scale:
a single nontrivial application contains many such finite checks, and hand-proving each
one would bury the mathematical argument under routine case analysis.

The second approach is computation-oriented.  Instead of proving each finite
calculation from scratch, we prove once that an algorithm correctly performs the
kind of calculation we need, and then we let Lean run that algorithm inside later
proofs.  This is the idea of \emph{reflection}: for each such operation, a
program that carries out the calculation, paired with a theorem saying that the
program's output agrees with the abstract definition.  Later proofs run the
program and invoke the theorem, rather than reconstructing the finite
verification case by case.  Flag algebras are especially well suited to this
style, because the candidate bounds and their supporting numerical data are
found by an SDP solver, so validating the resulting proof already amounts to
checking finite data.  Such data should
neither be taken on trust nor replayed as a hand proof, and reflection is the
middle path: Lean recomputes the entries itself.

\Cref{sec:abstract} gave the specification layer: definitions close to the
mathematics, where flags are identified up to isomorphism and quotient classes
hide arbitrary choices of representatives.  Those definitions are the right ones
for stating theorems, but many are marked \lean{noncomputable} because they are
specifications rather than efficient programs.  The reflection layer of our Lean
formalization is their computable counterpart.

\subsection{Concrete Graph Representations and Finite Search}
\label{sec:reflection-concrete}

The finite coefficients in a flag-algebra proof come from concrete counting
problems.  Computing a density \(\den{F}{G}\) means carrying out the sampling
experiment of \Cref{sec:background-densities}: enumerate the candidate vertex
subsets of the host \(G\), test which of them induce a copy of the pattern
\(F\), and divide by the number of choices.  The downward normalizing factor
\(q_\sigma(F)\) is a similar finite count, but of label placements rather than
subflags.

An executable flag-algebra development therefore needs two basic ingredients: a
way to enumerate the finite objects being counted, such as candidate induced
subflags, and a way to decide yes/no questions, such as whether two flags are
isomorphic.  In Lean these are the \lean{Fintype} and \lean{Decidable}
typeclasses, and together they make a definition that counts the objects
satisfying a predicate executable.

The specification layer of \Cref{sec:abstract}, however, is not designed
around producing such executable instances.  Its starting point is Mathlib's graph type
\lean{SimpleGraph V}, whose adjacency relation returns a proposition:
\begin{lstlisting}
structure SimpleGraph (V : Type*) where
  Adj      : V -> (*@$\,$@*) V -> (*@$\,$@*) Prop
  symm     : Symmetric Adj
  loopless : Irreflexive Adj
\end{lstlisting}
This works well for the mathematical specification, but it is too general to be
the data structure of the evaluator: the vertex type \lean{V} need not come with
an algorithm for enumerating its vertices, and \lean{Adj} need not come with a
Boolean decision procedure.  \lean{LabeledGraph}s and \lean{Flag}s, built on top
of this graph type, inherit the same generality; knowing that only finitely many
flags of each size exist up to isomorphism supplies neither ingredient.

The reflection layer therefore does not use \lean{SimpleGraph} directly as the
evaluator's representation.  Instead, it introduces a family of more rigid
edge-set representations: an \(n\)-vertex graph has vertex type \lean{Fin n},
whose vertices \(0,1,\ldots,n-1\) Lean can enumerate, and its edge set is stored
as \lean{Finset (Sym2 (Fin n))}, an explicit finite set of unordered pairs
against which adjacency is a membership test.  This representation is packaged
as \lean{Sym2Graph} and reused to define the concrete flag type
\lean{Sym2FlagType} and labeled graph \lean{Sym2LabeledGraph}.
\begin{lstlisting}
structure Sym2Graph (n : (*@$\mathbb{N}$@*)) where
  edges       : Finset (Sym2 (Fin n))
  edges_valid : forall e in edges, (*@$\neg$@*) e.IsDiag

abbrev Sym2FlagType (k : (*@$\mathbb{N}$@*)) := Sym2Graph k

structure Sym2LabeledGraph {k : (*@$\mathbb{N}$@*)} ((*@$\sigma$@*) : Sym2FlagType k) (n : (*@$\mathbb{N}$@*)) where
  edges       : Finset (Sym2 (Fin n))
  edges_valid : forall e in edges, (*@$\neg$@*) e.IsDiag
  type_embed  : fromEdgeSet (SetLike.coe (*@$\sigma$@*).edges) (*@$\hookrightarrow$@*)g fromEdgeSet (SetLike.coe edges)
\end{lstlisting}
Mathlib's constructor \lean{fromEdgeSet} builds an abstract \lean{SimpleGraph}
from a set of unordered pairs, and \lean{SetLike.coe} presents the edge
\lean{Finset} as such a set, so \lean{fromEdgeSet (SetLike.coe edges)} is the abstract graph decoded
from the concrete edge data.  Between the two decoded graphs, the field
\lean{type\_embed} carries the same embedding \(\theta_F\) as the
\lean{LabeledGraph} of \Cref{sec:flagtypes}.

These basic choices then scale to the finite searches used by the evaluator.
To compute a density \(\den{F}{G}\) it must range over the placements of the
pattern \(F\) inside the host \(G\), each of which chooses a vertex set of
\(G\) containing all labeled vertices and thereby determines an induced subflag
to be tested against \(F\).  The following \lean{Fintype} instance enumerates
those placements:
\begin{lstlisting}
variable {k m n : (*@$\mathbb{N}$@*)} {(*@$\sigma$@*) : Sym2FlagType k}

structure Sym2InducedLabeledSubgraph (G : Sym2LabeledGraph (*@$\sigma$@*) n) where
  verts : Finset (Fin n)
  verts_subset : G.type_verts (*@$\subseteq$@*) verts

def Sym2InducedLabeledSubgraph.edges {G : Sym2LabeledGraph (*@$\sigma$@*) n}
    (H : Sym2InducedLabeledSubgraph G) : Finset (Sym2 (Fin n)) :=
  G.edges.filter (fun (e : Sym2 (Fin n)) =>
    forall (v : Fin n), v (*@$\in$@*) e -> (*@$\,$@*) v (*@$\in$@*) H.verts)

instance (G : Sym2LabeledGraph (*@$\sigma$@*) n) : Fintype (Sym2InducedLabeledSubgraph G) where
  elems := (Finset.univ : Finset (Finset (Fin n))).filterMap
    (fun (V : Finset (Fin n)) =>
      if hV : G.type_verts (*@$\subseteq$@*) V
      then some ({ verts := V, verts_subset := hV } : Sym2InducedLabeledSubgraph G)
      else none)
    (by ...)
  complete := fun (H : Sym2InducedLabeledSubgraph G) => ...
\end{lstlisting}
The type \lean{Sym2InducedLabeledSubgraph G} represents induced subflags
of \(G\): it stores the chosen vertex set together with a proof that this set
includes \lean{G.type\_verts}, the image of the type embedding.  Its edge set
is then computed by filtering
\(G\)'s explicit edge set to the edges whose endpoints both lie in the
chosen vertex set.  The \lean{Fintype (Sym2InducedLabeledSubgraph G)} instance
provides an enumeration of all induced subflags of \(G\).  Concretely,
the instance contains an \lean{elems} field listing all elements of the type \lean{Sym2InducedLabeledSubgraph~G},
together with a \lean{complete} proof that this enumeration misses none of
them.  Here the
enumeration is deliberately simple: since the vertices of \(G\) are
\lean{Fin n}, Lean enumerates all subsets of \(\{0,\ldots,n-1\}\), keeps
exactly those containing \lean{G.type\_verts}, and stores the retained subset
with the proof of containment.
What remains is the second basic ingredient: a decidable predicate that, given
a candidate \(H\), decides whether its labeled graph is isomorphic to the
pattern \(F\) with the type labels respected.

The definitions below isolate this test at a single pattern.  They are written
for exposition: our Lean development states the same condition for tuples of
patterns, as \lean{predIsoSym2LabeledHl}, and decides it the same way.  As
stated, \lean{candidateMatches F H} is a proposition, not a program; the
reflection step is the instance that shows Lean how to decide it from the
concrete edge data.
\begin{lstlisting}
def candidateMatches (F : Sym2LabeledGraph (*@$\sigma$@*) m) {G : Sym2LabeledGraph (*@$\sigma$@*) n}
    (H : Sym2InducedLabeledSubgraph G) : Prop :=
  Nonempty (H.toLabeledSubgraph.coe (*@$\simeq$@*)f F.toLabeledGraph)

instance candidateMatchesDecidable
    (F : Sym2LabeledGraph (*@$\sigma$@*) m) (G : Sym2LabeledGraph (*@$\sigma$@*) n)
    : DecidablePred (fun H : Sym2InducedLabeledSubgraph G => candidateMatches F H) :=
  fun H => by
    have : Fintype H.toLabeledSubgraph.subgraph.verts := by ...
    have : DecidableRel H.toLabeledSubgraph.coe.graph.Adj := by ...
    have : DecidableRel F.toLabeledGraph.graph.Adj := by ...
    exact if hsize : H.verts.card = m then
      (by
        dsimp [candidateMatches] -- unfold the definition of candidateMatches
        infer_instance)
    else
      isFalse (fun hIso => hsize (verts_card_of_coe_iso H F hIso))

def sym2LabeledGraphDensity(*@$_1$@*) (F : Sym2LabeledGraph (*@$\sigma$@*) m) (G : Sym2LabeledGraph (*@$\sigma$@*) n)
    [DecidablePred (fun H : Sym2InducedLabeledSubgraph G => candidateMatches F H)]
    : (*@$\mathbb{Q}$@*) :=
  ((Finset.univ : Finset (Sym2InducedLabeledSubgraph G)).filter
    (fun H => candidateMatches F H)).card / Nat.choose (n - k) (m - k)
\end{lstlisting}
The body of \lean{candidateMatches} compares two abstract \lean{LabeledGraph}s,
each decoded from its concrete representation.
\lean{F.toLabeledGraph} decodes the pattern as above, and the candidate is
decoded in two steps: \lean{H.toLabeledSubgraph} is the labeled subgraph of
\lean{G.toLabeledGraph} on the vertices \lean{H.verts}, and \lean{.coe} views
that subgraph as a labeled graph in its own right, forgetting that it sits
inside \(G\).  The two are then compared with the isomorphism type
\(\simeq_f\) of \Cref{sec:flagtypes}, and \lean{Nonempty} asserts that such an
isomorphism exists.

The proof of the instance is the computational heart of the example.  After the
local \lean{Fintype} and \lean{DecidableRel} instances have been assembled, the
\lean{hsize} branch reaches \lean{infer\_instance}, which asks typeclass search
to decide that \lean{Nonempty} proposition; the generic instance does so by
enumerating finitely many graph isomorphisms and checking the type-preservation
condition.  The predicate has therefore been reduced to a finite search.  The
\lean{hsize} branch itself is not needed for that, but it reproduces the
implementation's prefilter: a candidate whose vertex set has the wrong
cardinality cannot be isomorphic to \(F\), so it is rejected before the
isomorphism search runs.

Once the matching predicate is decidable, the density is ordinary finite
evaluation: \lean{sym2LabeledGraphDensity}$_1$ filters the candidates by it,
counts the survivors, and divides by the number of ways to choose the non-type
vertices.  The multi-flag density \(\den{F_1,\ldots,F_t}{G}\) of
\Cref{sec:formal-subflag-densities} follows the same pattern, with the
tuple-level predicate additionally requiring the chosen non-type vertices to be
pairwise disjoint.  In particular, the computable counterparts of
\(\mathtt{flagDensity}_1\) and \(\mathtt{flagDensity}_2\) are
\lean{sym2FlagDensity}$_1$ and \lean{sym2FlagDensity}$_2$, whose arguments have
type \lean{Sym2Flag}, the reflection-layer counterpart of \lean{Flag}
constructed in \Cref{sec:reflection-adequacy}.  Unlike the specification-level
density functions, these definitions do not carry the \lean{noncomputable}
keyword: Lean has an executable algorithm for evaluating them.  But
computability alone does not say that they implement the density functions from
the specification layer.  That correspondence is established by the adequacy
theorems of the next subsection.

\subsection{Adequacy Theorems for Reflection}
\label{sec:reflection-adequacy}

\Cref{sec:flagtypes} defined flags by quotienting labeled graphs by
isomorphism.
The reflection layer repeats the same construction for the concrete graph
representation from \Cref{sec:reflection-concrete}:
\begin{lstlisting}
variable {k m m(*@$_0$@*) m(*@$_1$@*) n : (*@$\mathbb{N}$@*)} {(*@$\sigma$@*) : Sym2FlagType k}

instance sym2LabeledGraphSetoid ((*@$\sigma$@*) : Sym2FlagType k) (n : (*@$\mathbb{N}$@*))
    : Setoid (Sym2LabeledGraph (*@$\sigma$@*) n) where
  r     := sym2LabeledGraphEqv -- related by a label-preserving isomorphism
  iseqv := ... -- reflexivity, symmetry, and transitivity of r

def Sym2Flag ((*@$\sigma$@*) : Sym2FlagType k) (n : (*@$\mathbb{N}$@*)) : Type :=
  Quotient (sym2LabeledGraphSetoid (*@$\sigma$@*) n)
\end{lstlisting}
We now have two representations of flags: \lean{Flag} in the specification
layer and \lean{Sym2Flag} in the reflection layer.  To use a reflection-layer
computation in a specification-layer theorem, we need a map from
\lean{Sym2Flag} to \lean{Flag}.  That map is \lean{Sym2Flag.toFlag}:
\begin{lstlisting}
def Sym2Flag.toFlag (G : Sym2Flag (*@$\sigma$@*) n)
    : Flag (fromEdgeSet (SetLike.coe (*@$\sigma$@*).edges)) (Fin n) :=
  Quotient.lift
    (fun G : Sym2LabeledGraph (*@$\sigma$@*) n => (*@$\llbracket$@*)G.toLabeledGraph(*@$\rrbracket$@*))
    ... -- proof that the lift respects the equivalence relation of the quotient
    G
\end{lstlisting}

With \lean{Sym2Flag.toFlag} in hand, we can carry out the reflection described
in \Cref{sec:reflection}.  The preceding subsection supplied the programs; the
\emph{adequacy theorems} below are the theorems that certify them, saying that
\lean{sym2FlagDensity$_1$} and \lean{sym2FlagDensity$_2$} agree with
\lean{flagDensity$_1$} and \lean{flagDensity$_2$} once the \lean{Sym2Flag}
arguments are decoded:
\begin{lstlisting}
theorem flagDensity(*@$_1$@*)_eq_sym2FlagDensity(*@$_1$@*) (F : Sym2Flag (*@$\sigma$@*) m) (G : Sym2Flag (*@$\sigma$@*) n)
    : flagDensity(*@$_1$@*) F.toFlag G.toFlag = sym2FlagDensity(*@$_1$@*) F G := ...

theorem flagDensity(*@$_2$@*)_eq_sym2FlagDensity(*@$_2$@*)
    (F(*@$_0$@*) : Sym2Flag (*@$\sigma$@*) m(*@$_0$@*)) (F(*@$_1$@*) : Sym2Flag (*@$\sigma$@*) m(*@$_1$@*)) (G : Sym2Flag (*@$\sigma$@*) n)
    : flagDensity(*@$_2$@*) F(*@$_0$@*).toFlag F(*@$_1$@*).toFlag G.toFlag = sym2FlagDensity(*@$_2$@*) F(*@$_0$@*) F(*@$_1$@*) G := ...
\end{lstlisting}
A computable definition with its adequacy theorem is what we call a
\emph{reflected} computation.  A specification-level claim can then be
rewritten into the reflection layer and evaluated there.  For example,
the density \(\den{\fKtwobull,\,\fbarKtwobull}{\fbarPthreebull} = 1/2\) of
\Cref{ex:typed-edge-nonedge-product} is proved as follows:
\begin{lstlisting}
example : flagDensity(*@$_2$@*)
             Sym2Flag_2_1_0_0.toFlag
             Sym2Flag_2_1_0_1.toFlag
             Sym2Flag_3_1_0_1.toFlag = 1 / 2
  := by
  rw [flagDensity(*@$_2$@*)_eq_sym2FlagDensity(*@$_2$@*)]
  decide +kernel
\end{lstlisting}
The three constants are concrete \lean{Sym2Flag}s: \lean{Sym2Flag\_2\_1\_0\_0}
is \(\fbarKtwobull\), \lean{Sym2Flag\_2\_1\_0\_1} is \(\fKtwobull\), and
\lean{Sym2Flag\_3\_1\_0\_1} is \(\fbarPthreebull\).  Applying \lean{toFlag} to
each makes the statement a specification-level one, which \lean{rw} turns by
the adequacy theorem into the closed reflection-layer computation
\begin{lstlisting}
  sym2FlagDensity(*@$_2$@*)  Sym2Flag_2_1_0_0  Sym2Flag_2_1_0_1  Sym2Flag_3_1_0_1  =  1 / 2
\end{lstlisting}
The tactic \lean{decide +kernel} then closes this decidable goal by evaluating
that computation inside the Lean kernel.

The downward operator from \Cref{sec:formal-downward} follows the same
reflection pattern.  The specification-level definition uses the abstract
downward normalizing factor \(q_\sigma(F)\), implemented as
\lean{downwardNormalizingFactor}.  Its executable counterpart is
\lean{downwardNormalizingFactor\_Sym2Flag}.  On a concrete
\lean{Sym2LabeledGraph} over a flag type with \(k\) labels, on \(n\) vertices,
Lean enumerates the finite set of type embeddings that
produce an equivalent flag, takes its cardinality, and divides by the total
number \(n!/(n-k)!\) of injections of the \(k\) labels into the \(n\) vertices.
Because the computation is invariant under flag equivalence, it lifts through
the \lean{Sym2Flag} quotient.
\begin{lstlisting}
def isoEmbeddingCount_sym2LabeledGraph (G : Sym2LabeledGraph (*@$\sigma$@*) n) : (*@$\mathbb{N}$@*) :=
  (isoSym2TypeEmbeddingSetWithSameGraph G).card

def downwardNormalizingFactor_sym2LabeledGraph (G : Sym2LabeledGraph (*@$\sigma$@*) n) : (*@$\mathbb{Q}$@*) :=
  let num_of_all_injections := n.factorial / (n - k).factorial
  isoEmbeddingCount_sym2LabeledGraph G / num_of_all_injections

def downwardNormalizingFactor_Sym2Flag (F : Sym2Flag (*@$\sigma$@*) n) : (*@$\mathbb{Q}$@*) :=
  Quotient.lift
    (fun G => downwardNormalizingFactor_sym2LabeledGraph G)
    ... -- proof that downwardNormalizingFactor_sym2LabeledGraph respects the equivalence relation
    F
\end{lstlisting}
These are executable \lean{def}s rather than specification-level
\lean{noncomputable} ones.  We then prove that the specification-level and
executable implementations of the downward normalizing factor agree.
\begin{lstlisting}
theorem downwardNormalizingFactor_eq (F : Sym2Flag (*@$\sigma$@*) n)
    : downwardNormalizingFactor F.toFlag = downwardNormalizingFactor_Sym2Flag F := ...
\end{lstlisting}
Rewriting with it turns a goal about the specification-level coefficient into a
closed reflection-layer computation, which \lean{decide +kernel} evaluates as
before.

The reflected normalizing factor is what makes the downward operator executable
on concrete $\sigma$-flags.  Unlabeling such a flag yields the corresponding
$\emptyset$-flag by construction, so the coefficient is the only part that has
to be computed, and the adequacy theorem certifies
it.  By linearity, a finite linear combination of flag basis
elements is handled term by term, allowing later downward identities to be
established computationally rather than by separate counting arguments.

\subsection{Elaboration-Time Generation of Flags and Theorems}
\label{sec:reflection-generation}

The preceding subsection shows the basic unit of reflected computation: rewrite
a specification-level goal with an adequacy theorem, then evaluate the
resulting closed computation.  The adequacy theorems themselves are
uniform: they prove the reflected evaluators correct for every possible input.
Only their specialization to each concrete $\sigma$-flag and density instance
used by the later algebraic proof is repeated.

The generated flags are named by a fixed scheme, already visible in the
constants of \Cref{sec:reflection-adequacy}.  A name such as
\lean{Flag\_n\_k\_m\_i} or \lean{FlagAlgebra\_n\_k\_m\_i} carries four
indices: \(n\) is the size of the flag, \(k\) is the size of its flag type,
\(m\) is the position of the flag type's graph in the enumeration of the
\(k\)-vertex graphs, and \(i\) is the position of the flag itself in the
enumeration of the \(n\)-vertex flags of that flag type.  Both enumerations are
in a fixed canonical order, with positions counted from
zero.  Graphs are ordered first by their number of edges
and then by their canonical edge lists.  The canonical edge list of a graph is
the lexicographically smallest sorted edge list obtained by relabeling its
vertices in all possible ways.  The flags of a fixed type are grouped by their
underlying graph, in that same order, and within a group are ordered
lexicographically by the vertices the labels occupy.  The
reflection-layer constants \lean{Sym2Flag\_n\_k\_m\_i} follow the same
scheme, and the empty flag type is written \(k=m=0\).  Thus,
\lean{Sym2Flag\_3\_1\_0\_1} from \Cref{sec:reflection-adequacy} is flag
number \(1\) among the three-vertex flags of the unique one-vertex type,
namely \(\fbarPthreebull\), and \lean{FlagAlgebra\_3\_1\_0\_1} is its
flag-algebra basis element.

We automate this repetition with custom command elaborators written in Lean's
metaprogramming framework~\cite{ebner2017structured}.  During elaboration,
these commands run the finite enumerators and generate proofs of the
specialized facts.  A generation block has the following form:
\begin{lstlisting}[emph={generate_empty_typed_flags,generate_flags,generate_flag_pair_density_theorems_no_forbid,generate_mul_theorems},emphstyle=\color{blue!70!black}\bfseries]
-- Generate empty-typed flags at pattern size 2 and host size 3.
generate_empty_typed_flags 2
generate_empty_typed_flags 3

-- Use the unique flag type on one labeled vertex (k = 1, m = 0),
-- and generate its flags at the sizes 2 and 3.
generate_flags 2 1 0
generate_flags 3 1 0

-- Prove all pair densities, then all product expansions.
generate_flag_pair_density_theorems_no_forbid 2 3 1 0
generate_mul_theorems 2 3 1 0
\end{lstlisting}

The numerical arguments are the \(n\), \(k\), \(m\) of the naming scheme; the
density and multiplication commands take a pattern size \(p\) and a host size
\(h\) before the \(k,m\) of the flag type, and the suffix \lean{no\_forbid}
says that no forbidden-graph filter is being applied.
The flag-generating commands also emit a completeness lemma: every
\(n\)-vertex \(\sigma\)-flag, up to isomorphism, appears among the generated
names.  That lemma specializes a theorem proved once and for all, by induction
on the size: our enumeration algorithm produces a representative of every
isomorphism class.
\Cref{tab:generation-commands} collects the command family, including the
forbid-free variants below.

\begin{table}
  \centering
  \caption{The generation commands.  \(n\) is a flag size and \(k\), \(m\)
  select the flag type, as in the naming scheme; the density and
  multiplication commands take a pattern size \(p\) and a host size \(h\).
  The forbid-free variants take a further argument \lean{H}, a
  \lean{Sym2Graph} term for the forbidden graph.}
  \label{tab:generation-commands}
  \small
  \begin{tabular}{@{}lp{0.32\textwidth}@{}}
    \toprule
    Command & Generates \\
    \midrule
    \lean{generate\_empty\_typed\_flags n} &
      the empty-typed flags on \(n\) vertices, with the completeness lemma \\
    \lean{generate\_flags n k m} &
      the \(n\)-vertex flags of flag type \((k,m)\), with unlabelings, downward
      coefficients, and the completeness lemma \\
    \lean{generate\_flag\_pair\_density\_theorems\_no\_forbid p h k m} &
      the pair-density theorems for size-\(p\) patterns in size-\(h\) hosts \\
    \lean{generate\_mul\_theorems p h k m} &
      the product expansion~\eqref{eq:mul} for pairs of size-\(p\)
      patterns in size-\(h\) hosts \\
    \midrule
    \lean{generate\_forbid\_free\_empty\_typed\_flags n H} &
      the \(H\)-free empty-typed flags, with the filtered completeness lemma \\
    \lean{generate\_forbid\_free\_flags n k m H} &
      the \(H\)-free typed flags, with unlabelings, downward coefficients,
      and the filtered completeness lemma \\
    \lean{generate\_forbid\_free\_flag\_pair\_density\_theorems p h k m H} &
      the pair-density theorems over the \(H\)-free flags \\
    \lean{generate\_forbid\_free\_mul\_theorems p h k m H} &
      the product expansions, modulo the \(H\)-free condition \\
    \lean{generate\_forbid\_free\_flag\_density\_theorems p i h H} &
      the single-flag densities of \lean{Flag\_p\_0\_0\_i} in every
      \(H\)-free size-\(h\) host \\
    \bottomrule
  \end{tabular}
\end{table}

Generating the typed flags also proves a downward identity for each of them,
such as \(q_\sigma(\fbarPthreebull)=2/3\), and the pair-density command's
family includes the identity
\(\den{\fKtwobull,\,\fbarKtwobull}{\fbarPthreebull}=1/2\) proved earlier.  The
multiplication command combines those coefficients with the completeness lemma
for the hosts to prove the product expansion
\([F_1] \cdot [F_2]=\sum_G \den{F_1,F_2}{G}\,[G]\) of~\eqref{eq:mul} for every
pattern pair; completeness is what justifies summing over exactly the named
hosts.

Every bound compiled in \Cref{sec:flagmatic} is a bound on \(H\)-free graphs,
and each of the commands above has a forbid-free variant for that setting.  The
\lean{generate\_forbid\_free\_*} commands take one additional argument, a
concrete \lean{Sym2Graph} term for the forbidden graph \(H\); our Mantel
example, for instance, defines \lean{K3} as the complete graph on three
vertices and runs \lean{generate\_forbid\_free\_flags 3 1 0 K3}.  These
variants emit constants and theorems only for flags whose underlying graphs
avoid \(H\) as a subgraph, and their filtered completeness lemmas state that
the named flags exhaust precisely this set.

\medskip
\noindent\textbf{Optimizations.}
The following optimizations reduce the computational cost of elaboration.
They do not eliminate the inherent combinatorial growth of flag enumeration
and density computation, but make this paper's nontrivial examples feasible.

\begin{itemize}
\item \textbf{Size-inductive enumeration.}
The enumeration of flags is inductive in the size: the underlying
\(n\)-vertex graph representatives are obtained by extending the
\((n-1)\)-vertex representatives by one vertex in all possible ways, and the
flag type's labels are then placed on them in every valid
way.  The working lists therefore stay
close to one representative per isomorphism class, rather than covering all
\(2^{\binom{n}{2}}\) graphs on a fixed vertex set.

\item \textbf{Deduplication by cheap invariants.}
Both steps of that induction produce isomorphic duplicates: different vertex
extensions, and different label placements, can yield the same flag.  The
enumeration therefore deduplicates as it goes, keeping a candidate only when
it is isomorphic to no kept representative.  These isomorphism tests are the
expensive part, so each candidate carries a cheap isomorphism-invariant key,
derived from its edge count and degree sequence, and the expensive test runs
only when two keys collide.

\item \textbf{Pruning under a forbidden graph.}
For a complete forbidden graph, the size induction itself prunes: only
\(H\)-free representatives are extended, only \(H\)-free extensions kept,
and a cheap clique test decides containment, so flags containing \(H\) are
never generated.  Pruning misses no \(H\)-free isomorphism
class, since \(H\)-freeness is isomorphism-invariant and preserved by vertex
deletion, so the filtered completeness lemma again follows by
induction on size.

\end{itemize}

Metaprogramming automates theorem construction, not verification: Lean still
checks every generated proof, and a wrong coefficient makes that proof fail.
Once accepted, these theorems let later arguments simplify downward images,
densities, and products by rewriting rather than reopening the finite searches.
Thus, the story of this section has three steps: concrete representations make
the finite operations executable, adequacy theorems connect those operations
to the abstract definitions, and elaboration-time generation turns their
individual results into tables of proved facts.  The next section
uses those tables to automate complete flag-algebra arguments.

\section{Automating Flag-Algebra Proofs}
\label{sec:flagmatic}

Fix a forbidden family \(\mathcal H\), a target graph \(F\), and a proposed
bound \(c\).  The task addressed in this section is to turn finite certificate
data into a Lean proof of
\[
  [F]\leq_{\mathcal H}c\cdot\mathbf{1},
\]
which expresses the desired asymptotic upper bound on the induced density of
\(F\) in \(\mathcal H\)-free graphs.  A \emph{flag-algebra certificate} is a
finite witness for this inequality: it specifies flags, rational matrices, and
coefficients that purport to express the gap
\(c\cdot\mathbf{1}-[F]\) as a sum of terms that are non-negative under the
\(\mathcal H\)-free condition.  Flagmatic~\cite{vaughan2013flagmatic} searches
for such witnesses by semidefinite programming.

This section concerns certificate verification rather than certificate
search.  Our compiler accepts a certificate produced by Flagmatic as untrusted
input and runs inside Lean's elaborator: the desired theorem is stated in the
source file, and the \lean{flag\_certificate} tactic reads the certificate
file at elaboration time and synthesizes a proof of exactly that
statement.  Automating this translation is important because nontrivial
certificates contain extensive finite bookkeeping that would be tedious and
error-prone to reproduce manually.  The \emph{automation layer}, presented in
the following subsections, drives this compilation.

We first explain how positive semidefinite matrices justify the inequalities
encoded by a certificate.  We then identify the four obligations common to our
examples and show how the compiler discharges them.  Finally, we state the
trust model and evaluate the compiler on seven certificates.

\subsection{How a Certificate Proves a Bound}
\label{sec:flagmatic-background}

At a high level, a flag-algebra certificate proves
$[F] \leq_{\mathcal H} c\cdot\mathbf{1}$ by constructing two kinds of
non-negative correction term.  One kind consists of terms $Q_t$, each obtained
by applying the downward operator to a non-negative element of some
$\mathcal{A}^\sigma$, which a certificate supplies as a positive semidefinite
quadratic form in flags.  The other is a remainder $r$, a linear combination of
$\mathcal H$-free flags with non-negative coefficients.  Under the
$\mathcal H$-free condition, these data yield the chain
\begin{equation}
  \label{eq:cert-chain}
  [F]
  \leq_{\mathcal H} [F]+\sum_t Q_t
  = c\cdot\mathbf{1}-r
  \leq_{\mathcal H} c\cdot\mathbf{1}.
\end{equation}
The first inequality uses the non-negativity of the $Q_t$, the equality is the
finite calculation recorded by the certificate, and the final inequality uses
the non-negativity of $r$.

Mantel's theorem gives a small certificate of this shape.
Here the target is $F=K_2$, whose density is the edge density, and we
abbreviate its basis element $[K_2]$ by
$K_2$.
We work with the one-vertex type $\sigma$ and its two size-$2$ flags
$\fKtwobull$ and $\fbarKtwobull$, using these pictures directly in the
algebraic calculation below and omitting their brackets as usual.  Set
\[
  f := \tfrac12((\fKtwobull)-(\fbarKtwobull))^2,
  \qquad
  Q := \llbracket f \rrbracket_\sigma.
\]
We then show that
\[
  K_2 + Q = \tfrac12\cdot\mathbf{1} - \tfrac13\cdot\fbarPthree.
\]

For every positive homomorphism $\phi\in\poshom{\sigma}$,
\begin{align*}
  \phi(f)
  = \tfrac12\,\phi\bigl(((\fKtwobull)-(\fbarKtwobull))^2\bigr)
  = \tfrac12\,\bigl(\phi((\fKtwobull)-(\fbarKtwobull))\bigr)^2
  \;\geq\; 0.
\end{align*}
The first equality uses preservation of
scalar multiplication; the second uses preservation of multiplication, which
rewrites the flag-algebra square as the square of the real number
\(\phi((\fKtwobull)-(\fbarKtwobull))\).  Thus, $f$ is non-negative in
$\mathcal A^\sigma$, and \Cref{thm:downward-nonneg} gives $0\leq_{K_3}Q$.

To calculate $Q$ in the common size-three basis, we use the triangle-free
flags $\fbarKthree$, $\fbarPthree$, and $\fPthree$.  The size-three expansion
of $K_2$ from Equation~\eqref{eq:k2-expansion} is
\[
  K_2 = \tfrac13\cdot\fbarPthree + \tfrac23\cdot\fPthree.
\]
Applying Equation~\eqref{eq:sum-to-one} to the triangle-free empty-type flags of size
three gives
\[
  \mathbf{1} = \fbarKthree + \fbarPthree + \fPthree.
\]
A size-three product expansion followed by the downward operator gives
\begin{align*}
  Q
  &=
  \left\llbracket
    \tfrac12\left((\fKtwobull)-(\fbarKtwobull)\right)^2
  \right\rrbracket_\sigma\\
  &=
  \left\llbracket
    \tfrac12\left(
       \fbarKthreebull+\fEthreebull
       -\fbarPthreebull-\fPthreebull
       +\fKonetwobull
     \right)
  \right\rrbracket_\sigma\\
  &= \tfrac12\left(
       1\cdot\fbarKthree
       + \tfrac13\cdot\fbarPthree
       - \tfrac23\cdot\fbarPthree
       - \tfrac23\cdot\fPthree
       + \tfrac13\cdot\fPthree
     \right)\\
  &= \tfrac12\cdot\fbarKthree
     -\tfrac16\cdot\fbarPthree
     -\tfrac16\cdot\fPthree.
\end{align*}
The coefficients produced by the downward operator are its normalizing factors:
$1/3$, $2/3$, or $1$ according as one, two, or all three choices of the
labeled vertex recover the displayed $\sigma$-flag.

Combining these identities now gives a derivation of the desired bound:
\begin{align*}
  K_2
  &\leq_{K_3} K_2+Q \\
  &= \left(\tfrac13\cdot\fbarPthree + \tfrac23\cdot\fPthree\right)
     + \left(\tfrac12\cdot\fbarKthree
     -\tfrac16\cdot\fbarPthree
     -\tfrac16\cdot\fPthree\right)\\
  &= \tfrac12\cdot\fbarKthree
     + \tfrac16\cdot\fbarPthree
     + \tfrac12\cdot\fPthree\\
  &= \tfrac12\cdot
     \left(\fbarKthree+\fbarPthree+\fPthree\right)
     - \tfrac13\cdot\fbarPthree\\
  &= \tfrac12\cdot\mathbf{1} - \tfrac13\cdot\fbarPthree\\
  &\leq_{K_3} \tfrac12\cdot\mathbf{1}.
\end{align*}
The final inequality follows because the remainder
$r=\tfrac13\cdot\fbarPthree$ is non-negative.

The same idea scales when the single square is replaced by a positive
semidefinite quadratic form.  Choose $\sigma$-flags $E_1,\ldots,E_r$, let
$e_i:=[E_i]\in\mathcal A^\sigma$ be their basis elements, and collect these
algebra elements into a vector $v=(e_1,\ldots,e_r)^\top$.  A symmetric matrix
$M\in\mathbb{R}^{r\times r}$ is \emph{positive semidefinite}
(PSD) if $x^\top Mx\geq 0$ for every $x\in\mathbb{R}^r$.  Consequently, when $M=(M_{ij})$
is PSD, the flag-algebra quadratic form
$v^\top Mv=\sum_{i,j}M_{ij} e_i e_j$ is non-negative under every
positive homomorphism: evaluating the algebra elements turns it into the ordinary real
quadratic form $x^\top Mx$.  The downward operator then transfers this
non-negativity from the labeled algebra to the unlabeled one.

Flagmatic finds the matrices by turning the search for them into a
\emph{semidefinite program} (SDP).  It first fixes a host size $N$ and chooses
the flag types and flags that may occur.  It expands the target, the unit, and
the downward quadratic forms in the common basis of allowed $N$-vertex graphs,
namely those that avoid $\mathcal H$, with the matrix entries as
unknowns.  Requiring the coefficient left over for every host graph to be
non-negative gives linear constraints on those entries, and requiring each
matrix to be PSD gives the semidefinite constraints.  An SDP solver finds
numerical matrices satisfying these conditions, after which Flagmatic converts
them into exact rational matrices.

A full certificate uses several flag types at once, so the single matrix $M$
above becomes a family: the index $t$ ranges over the chosen flag types
$\sigma_t$, each contributing its own flag vector $v_t$ and matrix $M_t$.  We
call one such triple an \emph{SDP block};
the name reflects that the matrices $M_t$ are the diagonal blocks of the
semidefinite program's matrix variable.  The downward quadratic form
block~$t$ produces is
\[
  Q_t := \bigl\llbracket
    v_t^\top M_t\, v_t
  \bigr\rrbracket_{\sigma_t}.
\]
The resulting certificate specifies an identity, modulo the $\mathcal H$-free
condition, of the form
\[
  [F] + \sum_t Q_t
  =
  c\cdot\mathbf{1} - r,
\]
where every $M_t$ is PSD and $r$ is a linear combination of $\mathcal H$-free
flags with non-negative coefficients, so these data justify the
chain~\eqref{eq:cert-chain}.  Its equality is now the certificate identity
checked after all terms have been expanded in the common basis at host
size~$N$.  The certificate records the bound $c$, the chosen
flags, and the rational matrices $M_t$.  Our compiler takes none of these
numbers on trust: it recomputes the
densities and product expansions in Lean, proves the
matrices are PSD, and checks the calculation.

\subsection{Compiling a Certificate}
\label{sec:flagmatic-shape}
\label{sec:flagmatic-pipeline}
\label{sec:flagmatic-gen}\label{sec:flagmatic-psd}\label{sec:flagmatic-expand}%
\label{sec:flagmatic-close}\label{sec:flagmatic-script}%

Turning a certificate into a verified proof raises four obligations, referred
to below as Parts~1--4:

\begin{enumerate}
  \item \textbf{Generating the flags and their identities.}  The proof must name
    the flags and corresponding flag-algebra basis elements it uses,
    restricted to those free of the chosen forbidden graph, and must prove the
    density and product identities it needs.
  \item \textbf{Certifying the matrices.}  For each SDP block, the proof must
    show that the certificate's matrix $M_t$ is positive semidefinite.
  \item \textbf{Expanding the target (when needed).}  If the target graph $F$
    has fewer vertices than the host size $N$ chosen to bound the certificate
    search in \Cref{sec:flagmatic-background}, the proof must expand $[F]$ in
    the size-$N$ basis.
  \item \textbf{Assembling the main theorem.}  The proof must assemble the
    chain~\eqref{eq:cert-chain} from the results of the preceding parts.
\end{enumerate}

Reusable proof-producing components handle each part, so a new theorem
primarily requires new certificate data rather than a new proof architecture.

For the certificate shapes supported by the current implementation, this
process is exposed as the \lean{flag\_certificate} tactic.  A source file
states the desired bound and proves it by naming the certificate file and the
forbidden graph:
\begin{lstlisting}[emph={flag_certificate},emphstyle=\color{blue!70!black}\bfseries]
theorem Mantel_flagAlgebra
    : FlagAlgebra_2_0_0_1 (*@$\leq$@*)[completeGraph (Fin 3)] (1 / 2 : (*@$\mathbb{R}$@*)) (*@$\leansmul$@*) (1 : FlagAlgebra (*@$\emptyt$@*))
  := by
  flag_certificate "<mantel-cert>.json" K3
\end{lstlisting}
The variant \lean{flag\_certificate?} additionally offers
the synthesized script as a standard ``Try this:'' suggestion, so a single
editor click replaces the tactic call with the explicit proof (inline
matrices, PSD checks, and the closing normalization), which can then be
audited, archived, or checked independently of the certificate file.

A complete source file needs a few pieces around that theorem: the forbidden
graph as a \lean{Sym2Graph} term, the \lean{generate\_forbid\_free\_*} commands
of \Cref{sec:reflection-generation} that supply Part~1, option settings
choosing kernel evaluation and lifting Lean's heartbeat and recursion limits,
and, after the theorem, a restatement of the bound as a generalized Tur\'an
density (\Cref{sec:results}).  They are written once per certificate, by hand
or all at once by the \lean{flagmatic\_to\_lean.py} script:
\begin{lstlisting}[language={}]
python flagmatic_to_lean.py gen-skeleton <cert>.json <target>.lean
\end{lstlisting}
Passing \lean{--materialize} instead emits the fully
expanded legacy form, in which the matrices of Part~2, the target-expansion
lemma of Part~3, and the complete tactic proof of Part~4 appear in the
source file itself and the certificate is not read at build time.  We now
explain the synthesized proof one part at a time.

\boldparagraph{Generating the flags and their identities (Part~1).}
This is where the reflection layer of \Cref{sec:reflection} enters the
compiler, making the finite bookkeeping over flags, densities, and products
executable and checkable.  Part~1 is the scaffolding that lives in the source
file, and the prerequisite the \lean{flag\_certificate} tactic checks for,
failing with the exact block of generation commands to add when any are
missing.  It
defines the certificate's forbidden graph as a \lean{Sym2Graph} term and
invokes the \lean{generate\_forbid\_free\_*} commands of
\Cref{sec:reflection-generation} on it.  The tactic matches each of the certificate's
compact graph strings against the resulting enumeration, selecting the
corresponding
\lean{FlagAlgebra\_n\_k\_m\_i} constant, so Flagmatic's ordering need not
agree with Lean's.

The same commands prove the identities Parts~3 and~4 consume: a downward
normalizing factor $q_\sigma(F)$ for each emitted $\sigma$-flag, the product
expansions of the flag pairs the certificate uses, and, when the target must be
expanded in Part~3, the required single-flag densities.  The adequacy theorems
of \Cref{sec:reflection-adequacy} state all of them at the specification level.

\boldparagraph{Certifying the matrices (Part~2).}
Part~2 discharges the
positive-semidefiniteness obligations.  For each SDP block, the certificate
stores rational matrices $Q'_t$ and $R_t$ with
$M_t=R_tQ'_tR_t^\top$, where $R_t$ may be rectangular and $Q'_t$ is a smaller
symmetric matrix. Flagmatic's user guide~\cite{vaughan2012flagmaticguide}
explains the format: $Q'_t$ is
chosen to be positive definite, which implies the positive semidefiniteness of $M_t$.

Our verification takes these data through a different, equally elementary
witness.  Every PSD obligation is discharged through one library lemma, based
on an exact rational $LDL^\top$ factorization:
\begin{lstlisting}
theorem posSemidef_real_of_LDLt {n : (*@$\mathbb{N}$@*)}
    {M L : Matrix (Fin n) (Fin n) (*@$\mathbb{Q}$@*)} {d : Fin n -> (*@$\,\mathbb{Q}$@*)}
    (hd : forall (*@$\,$@*)i, 0 <= (*@$\,$@*)d i) (hM : M = L * Matrix.diagonal d * L(*@$^\top$@*))
    : (ratMatrixToReal M).PosSemidef := ...
\end{lstlisting}
The entrywise $\mathbb{Q}$-to-$\mathbb{R}$ cast, written
\lean{ratMatrixToReal}, happens inside the lemma, so the factorization itself
stays in exact rational arithmetic.  To use it, the tactic reconstructs $M_t$ from $Q'_t$ and $R_t$ in
exact rational arithmetic during elaboration, computes a factorization
$M_t=L_t\operatorname{diag}(d_t)L_t^\top$ with $L_t$ square and unit lower
triangular, and inlines $M_t$, $L_t$, and $d_t$ as literals in the
synthesized proof.  Lean treats the inlined factorization only as a candidate
witness: the \lean{psd\_real\_ldlt\_terms} tactic must prove the lemma's two
hypotheses.  Both are concrete decidable statements about literals, the
non-negativity of every entry of $d_t$ and an equality of finite matrices
over $\mathbb{Q}$, so \lean{decide +kernel} evaluates them and has the Lean
kernel certify the results; no eigenvalue or floating-point calculation
enters.  An $L_t$ or $d_t$ that does not factor $M_t$ is
therefore rejected.  What Lean does not check is that $M_t$ agrees with
the certificate's $Q'_t$ and $R_t$ fields: their reconstruction is done by the
tactic's unverified elaboration-time code, and \Cref{sec:flagmatic-trust}
takes up fidelity to the source certificate.

\boldparagraph{Expanding the target (Part~3).}
The certificate identity of \Cref{sec:flagmatic-background} is expressed at the
host size $N$, so Part~3 has nothing to do when the target graph $F$ already
has $N$ vertices.  When $F$ is smaller, the compiler proves an auxiliary
identity that expands $[F]$ over the $N$-vertex $\mathcal{H}$-free
$\emptyset$-flags,
\[
  [F]
  \;=\;
  \sum_{\substack{G \in \mathcal{F}^{\emptyset}_{N} \\ G\ \mathcal{H}\text{-free}}}
  \den{F}{G}\,[G]
  \qquad\text{modulo the \(\mathcal{H}\)-free condition,}
\]
with coefficients $\den{F}{G}$ given by the single-flag densities
generated in Part~1.

All certificates considered here have a singleton forbidden family
$\mathcal{H}=\{H\}$.  The generated proof invokes the
\lean{flag\_expand\_hfree N K} tactic, where $N$ is the host size and $K$ is the
reflection-layer \lean{Sym2Graph} encoding of $H$.  That the identity holds
only modulo the \(H\)-free condition matters, because our development works in
the ambient algebra.  In
the built-in \(H\)-free algebra (\Cref{rem:builtin-hfree}), the restricted sum
would simply be the expansion identity~\eqref{eq:expansion-identity}, all of
whose flags are \(H\)-free there.  In the ambient algebra, however, the flags containing
\(H\) are nonzero basis elements, so dropping their terms cannot give a
literal equality; the auxiliary identity instead says that every empty-type
positive homomorphism satisfying the \(H\)-freeness condition of
\Cref{def:ensemble-semantic-order} gives both sides the same value.
This identity is what lets the rest of the proof expand over only the flags
Part~1 emitted.

\boldparagraph{Assembling the main theorem (Part~4).}
Part~4 assembles the outputs of the preceding parts and realizes the
chain~\eqref{eq:cert-chain}.
The first step is to prove
$[F]\leq_H[F]+\sum_tQ_t$.  It uses the following lemma:
\begin{lstlisting}
theorem forbidLEWith_add_QuadraticForm {k n : (*@$\mathbb{N}$@*)} {(*@$\sigma$@*) : FlagType (Fin k)}
    {C : ForbidCondition} {f g : FlagAlgebra (*@$\emptyt$@*)}
    (M : Matrix (Fin n) (Fin n) (*@$\mathbb{R}$@*)) (hM : M.PosSemidef) (v : FlagAlgebraVec (*@$\sigma$@*) n)
    : forbidLEWith C f g -> (*@$\,$@*) forbidLEWith C f (g + (*@$\llbracket$@*)flagQuadraticForm M v(*@$\rrbracket_0$@*)) := ...
\end{lstlisting}
Here \lean{v} is the flag vector $v_t$ of \Cref{sec:flagmatic-background} and
\lean{flagQuadraticForm M v} is its quadratic form $v_t^\top M_t v_t$;
\lean{forbidLEWith} is the $\leq_H$ of \Cref{sec:forbidden}, instantiated at
\lean{forbiddenCondition H}, and \lean{hM} is the positive-semidefiniteness
proof generated in Part~2.
The proof starts from the reflexive inequality $[F]\leq_H[F]$ and applies this
lemma once for each SDP block.  After all blocks have been added, its conclusion is precisely
$[F]\leq_H[F]+\sum_tQ_t$.

It remains to prove $[F]+\sum_tQ_t\leq_H c\cdot\mathbf{1}$.  The compiler checks
this comparison by rewriting both sides in the basis of $H$-free flags at the
host size $N$.  When necessary, the Part~3 lemma first expands $[F]$ to that
basis.
Expanding each $Q_t=\bigl\llbracket v_t^\top M_t v_t\bigr\rrbracket_{\sigma_t}$
produces downward images of products of basis
elements; the \lean{reduce\_downward\_flagmul} tactic rewrites them using the
multiplication and downward identities generated in Part~1.  Then
\lean{expand\_one\_hfree\_at} expands $\mathbf{1}$ in the same basis, and
\lean{flagsum\_ac\_sort\_rhs\_pipeline} collects like basis terms.

After these rewrites, the difference between the right- and left-hand sides is
exactly the remainder $r$ of~\eqref{eq:cert-chain}.  The tactic
\lean{flag\_nonneg} checks that its coefficients are non-negative and then uses
the non-negativity of every basis element to conclude $0\leq_H r$.

\subsection{Trust Model}
\label{sec:flagmatic-trust}

There are two distinct questions of trust: is the theorem that Lean accepts
valid, and is it the statement the certificate was meant to yield?

For the first question, neither the certificate nor the compiler frontends
(the \lean{flag\_certificate} metaprogram and the scaffolding script) are
trusted for the logical validity of a theorem that Lean accepts.  No example
file introduces an axiom declaration or a \lean{sorry}, so every claim is carried by a proof term Lean checks: the
reconstructed matrices must be proved positive semidefinite, the certificate
calculation must normalize to non-negative coefficients, and the final
theorem must typecheck.  Subject to Lean's trusted base, erroneous data
cannot make an invalid proof term valid.

For the second question, the anchor is the theorem statement fixed in the
source file.  The tactic derives every ingredient of the proof from the
certificate, but it can only close the goal that the statement poses: a
mistranscribed target or forbidden graph, or a bound stated smaller than the
one the certificate proves, surfaces as an elaboration
failure.  What remains to audit is the statement itself: that
$[F]\leq_{\mathcal H}c\cdot\mathbf 1$, as written, expresses the intended
extremal claim.  In files emitted by \lean{gen-skeleton} the script writes that
statement, so reading it remains the auditability boundary.  In
the \lean{--materialize} form the certificate is not read at build time, so a
successful build does not say which certificate the file came from; Lean
still checks that the file's data proves the stated theorem.

The proof obligations also use two evaluation modes.  Five of the seven
generated files of \Cref{sec:results}, including the Erd\H{o}s pentagon,
discharge every finite computation by \lean{decide +kernel}, so no
compiled-evaluation axiom appears.  The $K_5$-free and $C_5$-free edge-density
files use \lean{native\_decide} instead and additionally trust Lean's native
compiler and runtime.  Finally, the custom
tactics are unverified metaprograms whose dependence on elaborated expression
shapes affects robustness and portability.  They construct proof terms that
Lean checks, so a tactic failure leaves the proof incomplete rather than
proving an invalid theorem.

\subsection{Evaluation}
\label{sec:results}\label{sec:flagmatic-scenarios}

We evaluated the compiler on seven certificates spanning host sizes
$N\in\{3,4,5\}$, targets at and below the host size, between one and four SDP
blocks, and four
forbidden graphs ($K_3$, $K_4$, $K_5$, and $C_5$).  All seven generated files
compile and contain neither \lean{sorry} nor generated axiom declarations.  No
certificate coefficient appears in the source at all: the data is consumed at
elaboration time.
Table~\ref{tab:flagmatic-scenarios} summarizes the seven cases;
\Cref{sec:compile-times} reports their compilation times and peak memory
under both \lean{decide +kernel} and \lean{native\_decide}.

\begin{table}
  \centering
\caption{Certificates carried through the compiler.  \(H\) is the
  forbidden graph, \(n_F\) the number of vertices of the target graph \(F\),
  \(N\) the host size, and \(T\) the number of SDP blocks.}
  \label{tab:flagmatic-scenarios}
  \begin{tabular}{lccccc}
    \toprule
    Target density & $H$ & $n_F$ & $N$ & $T$ & Bound \\
    \midrule
    Edge density (Mantel)               & $K_3$ & 2 & 3 & 1 & $1/2$    \\
    $P_3$ density                       & $K_3$ & 3 & 3 & 1 & $3/4$    \\
    $C_4$ density                       & $K_3$ & 4 & 4 & 2 & $3/8$    \\
    Edge density                        & $K_4$ & 2 & 4 & 2 & $2/3$    \\
    $C_5$ density (Erd\H{o}s pentagon)  & $K_3$ & 5 & 5 & 3 & $24/625$ \\
    Edge density                        & $K_5$ & 2 & 5 & 4 & $3/4$    \\
    Edge density                        & $C_5$ & 2 & 5 & 4 & $1/2$    \\
    \bottomrule
  \end{tabular}
\end{table}

The cases play complementary roles.  Mantel is a hand-checkable baseline: the
quadratic form of its single $2\times2$ SDP block is the explicit square from
\Cref{sec:flagmatic-background}.  The largest cases, the $K_5$- and $C_5$-free
edge densities, each use four blocks of $8\times8$
matrices.  These are Flagmatic's certificates as
produced, with no attempt to minimize the number of blocks or the matrix
dimensions; the unoptimized output already compiles end to end.

Each generated file ends by transferring its flag-algebra bound to the
specification level: an auto-generated theorem restates it as the generalized
Tur\'an-density bound $\tdensity{F}{\mathcal H}\le c$, through the bridge
lemma \lean{generalizedTuranDensity\_le\_of\_forbidLE}, with the target graph
spelled as an explicit edge list; the file's final statement therefore
mentions none of the generated \lean{FlagAlgebra\_n\_k\_m\_i} constants and can
be audited against the specification-layer definitions alone.  The certificate argument establishes only this
upper-bound direction.  For Mantel and the Erd\H{o}s pentagon,
the matching lower bounds are proved directly in Lean using the underlying
formalization (\Cref{sec:lowerbounds}), completing the
Tur\'an densities $\tfrac12$ and $\tfrac{24}{625}$ of Examples~\ref{ex:mantel}
and~\ref{ex:pentagon}; the Erd\H{o}s pentagon case is, to our knowledge, the
first machine-checked proof of that theorem.

\boldparagraph{Limitations.}
\label{sec:flagmatic-open}
We point out three limitations of the current compiler.
First, this evaluation establishes coverage of the seven cases above, not
completeness of the compiler.  Even a certificate within the compiler's
intended scope may expose a normalization shape not handled by the custom
tactics.  Such failures are safe: an unsupported or malformed input is
rejected or yields a skeleton containing \lean{sorry}, and neither case
produces an unchecked theorem.
Second, fully kernel-checked compilation does not yet scale.  For the
$K_5$-free and $C_5$-free edge-density certificates, we attempted to check the
generated files entirely by \lean{decide +kernel}; the runs did not finish in
an acceptable time and were abandoned, which is why those two files use
\lean{native\_decide} (\Cref{sec:flagmatic-trust}).
Third, the statement anchor narrows the fidelity concern of
\Cref{sec:flagmatic-trust} but does not eliminate it: the statement must still
be read, and in generated files the unverified script writes
it.  One smaller cost of the elaboration-time design is
also worth recording: the certificate becomes a build-time input that the
build system does not track, so editing a certificate does not by itself
trigger rechecking of the Lean file that reads it.

\section{Using the Formalization Beyond Certificate Compilation}
\label{sec:lowerbounds}

Underlying the compiler of \Cref{sec:flagmatic} is a reusable formalization of
flag-algebra expressions, their semantics, and the convergence theorems
(\Cref{thm:convergent-hom}) that tie that semantics to densities in finite
graphs.  These components can also be
used without a certificate at all, and this section illustrates three such
uses.  Mantel's theorem is proved in full against the formalization, without
the compiler.  For the Erd\H{o}s pentagon theorem, an explicit construction
and the convergence of normalized extremal numbers supply the lower bound
matching the compiled upper bound of \Cref{sec:results}, completing the exact
density.  Finally, two inequalities of
Goodman show that the same framework supports flag-algebra arguments outside
the Tur\'an-density problems considered in \Cref{sec:flagmatic}.

\boldparagraph{Mantel's theorem.}
\Cref{ex:mantel} states $\tdensity{K_2}{K_3}=\tfrac12$.  The Lean development
records it as follows:
\begin{lstlisting}
theorem Mantel_flag_bound : K2 (*@$\leq$@*) (1 / 2 : (*@$\mathbb{R}$@*)) (*@$\leansmul$@*) 1 + K3 := ...

theorem Mantel_flag_bound' : (*@$\forall$@*) ((*@$\phi$@*) : PositiveHom (*@$\emptyt$@*)), (*@$\phi$@*) K3 = 0 (*@$\to$@*) (*@$\phi$@*) K2 (*@$\leq$@*) 1 / 2 := ...

lemma extremal_density_K3_ge (n : (*@$\mathbb{N}$@*)) (hn2 : n (*@$\geq$@*) 2)
    : (extremalNumber n (completeGraph (Fin 3)) / n.choose 2 : (*@$\mathbb{R}$@*)) (*@$\geq$@*) 1 / 2 := ...

theorem Mantel_Turan : turanDensity (completeGraph (Fin 3)) = 1 / 2 := ...
\end{lstlisting}
Here \lean{K2} and \lean{K3} denote the basis elements $[K_2]$ and $[K_3]$.
The compiler of \Cref{sec:flagmatic} already produces the upper bound; this
development instead uses neither the compiler nor the
\lean{forbidLE} interface of \Cref{sec:forbidden}.

The direct upper-bound argument takes place in the ambient algebra and keeps
the triangle contribution explicitly:
\[
  K_2\leq\tfrac12\cdot\mathbf{1}+K_3.
\]
This inequality is formalized as \lean{Mantel\_flag\_bound}.  If a positive
homomorphism $\phi$ represents a triangle-free limit, then
$\phi([K_3])=0$, and evaluating the displayed inequality gives
$\phi([K_2])\leq\tfrac12$.  This pointwise statement is recorded as
\lean{Mantel\_flag\_bound'} and drives the remaining argument.

The matching lower bound comes from the balanced complete bipartite graph
$K_{\lfloor n/2\rfloor,\lceil n/2\rceil}$, whose
$\lfloor n^2/4\rfloor$ edges make up at least half of the $\binom{n}{2}$
vertex pairs; it is formalized as \lean{extremal\_density\_K3\_ge}.  Here
\lean{extremalNumber} denotes the maximum
number of edges in a $K_3$-free graph on $n$ vertices, which agrees with
$\mathrm{ex}(n,\,K_2;\,K_3)$ because the induced copies of $K_2$ are precisely
the edges.  Finally, \lean{Mantel\_Turan} combines this construction with
\lean{Mantel\_flag\_bound'} through \Cref{lem:semantic-bound}, without the
compiler or the ensemble semantics of \Cref{sec:forbidden}.

\boldparagraph{The Erd\H{o}s pentagon theorem.}
The full result of \Cref{ex:pentagon},
\begin{lstlisting}
theorem ErdosPentagon_Turan : generalizedTuranDensity K3 C5 = 24/625 := ...
\end{lstlisting}
is the conjunction of the upper bound of \Cref{sec:results} and the lower
bound
\begin{lstlisting}
theorem ErdosPentagon_Turan_lowerBound : generalizedTuranDensity K3 C5 (*@$\geq$@*) 24/625 := ...
\end{lstlisting}
proved by an explicit construction in Lean.  To our knowledge,
\lean{ErdosPentagon\_Turan} is the first machine-checked proof of the
Erd\H{o}s pentagon theorem; all known proofs of the theorem use flag
algebras~\cite{grzesik2012,hatami2012}, so a machine-checked proof
required a formalization of the method itself.

Since, by the convergence theorem
\lean{tendsto\_\allowbreak generalizedTuranDensity} of \Cref{sec:forbidden},
the normalized extremal numbers $\mathrm{ex}(n,\,C_5;\,K_3)/\binom{n}{5}$
converge to $\tdensity{C_5}{K_3}$, it suffices to exhibit, for infinitely many
graph sizes, one triangle-free graph whose normalized count of induced copies of
$C_5$ is at least $\tfrac{24}{625}$.  The
witness is the \emph{$n$-fold blow-up} of $C_5$: replace each vertex of
$C_5$ by a class of $n$ vertices, and join two vertices of distinct
classes exactly when the corresponding vertices of $C_5$ are adjacent;
vertices of a common class remain non-adjacent.  The Lean development is:
\begin{lstlisting}
def blowUp {V : Type} [Fintype V] (G : SimpleGraph V) (n : (*@$\mathbb{N}$@*))
    : SimpleGraph (V (*@$\times$@*) Fin n) :=
  { Adj v w := G.Adj v.1 w.1
    symm v w := by apply G.symm }

theorem blowUp_K3_free {m : (*@$\mathbb{N}$@*)} {G : SimpleGraph (Fin m)}
    (n : (*@$\mathbb{N}$@*)) (hfree : K3.Free G) : K3.Free (blowUp G n) := ...

lemma subgraphCount_blowUp_C5_ge (n : (*@$\mathbb{N}$@*))
    : subgraphCount C5 (blowUp C5 n) (*@$\geq$@*) n ^ 5 := ...

theorem generalizedExtremalNumber_K3_C5_ge (n : (*@$\mathbb{N}$@*))
    : generalizedExtremalNumber (5 * n) K3 C5 (*@$\geq$@*) n ^ 5 := ...
\end{lstlisting}
Here \lean{blowUp G n} realizes the classes as the product type
\lean{V $\times$ Fin n}, with adjacency read off the first coordinates alone.
The $n$-fold blow-up of $C_5$ is triangle-free, since adjacent vertices lie in
distinct classes and so the classes of a hypothetical triangle would form a
triangle in $C_5$ itself; \lean{blowUp\_K3\_free} states this for the blow-up
of any $K_3$-free base graph.  The blow-up also contains many copies of $C_5$:
picking one vertex from each of the five classes yields five vertices that
inherit exactly the adjacencies of $C_5$, so each of the $n^5$ choice
functions induces a copy.  \lean{subgraphCount\_blowUp\_C5\_ge} records that
count, where \lean{subgraphCount C5 G} counts induced subgraphs of \lean{G}
isomorphic to $C_5$.  So the blow-up is a triangle-free graph on $5n$ vertices
with at least $n^5$
induced copies of $C_5$, which
\lean{generalizedExtremalNumber\_K3\_C5\_ge} records as
$\mathrm{ex}(5n,\,C_5;\,K_3)\geq n^5$.  Normalizing gives, for every $n\geq1$,
\[
  \frac{\mathrm{ex}(5n,\,C_5;\,K_3)}{\binom{5n}{5}}
  \;\geq\;
  \frac{n^5}{\binom{5n}{5}}
  \;\geq\;
  \frac{n^5}{(5n)^5 / 5!}
  \;=\;
  \frac{24}{625}.
\]

There is one additional representation step in Lean.  The natural vertex set
of \lean{blowUp C5 n} is the product \lean{Fin 5 $\times$ Fin n}, while
\lean{generalizedExtremalNumber} takes its maximum over graphs on
\lean{Fin (5 * n)}.  The graph \lean{blowUp\_fin} simply reindexes the blow-up
on this latter vertex set; it is proved isomorphic to \lean{blowUp} and is used
as the witness.  Normalizing the finite bound above and passing to the limit
along the graph sizes $5n$ then gives
\lean{ErdosPentagon\_Turan\_lowerBound}.

\boldparagraph{Direct Goodman inequalities.}
The final examples use the core flag-algebra infrastructure directly for two
inequalities outside the Tur\'an-density problems considered in
\Cref{sec:flagmatic}.  Both are due to Goodman~\cite{goodman1959}.
The first is the algebraic lower bound
\[
  K_3 \;\geq\; K_2\cdot(2K_2-\mathbf{1}),
\]
which bounds the triangle density from below in terms of the edge density:
under a positive homomorphism with edge density $e$, the triangle density is
at least $e(2e-1)$, a quantity that becomes positive as soon as $e$ exceeds
the Mantel threshold $\tfrac12$.  The proof combines the size-three expansion
of $K_2$ with the Cauchy--Schwarz inequality
$\llbracket f\,g\rrbracket_\sigma^{\,2}\leq
\llbracket f^2\rrbracket_\sigma\cdot\llbracket g^2\rrbracket_\sigma$ for the
downward operator, applied to the edge flag labeled at one endpoint.  That
inequality is a separate theorem in Razborov's
paper~\cite{razborov2007flag}, proved in the development in this general form
as \lean{Cauchy\_\allowbreak Schwarz\_\allowbreak inequality}.  The resulting inequality is
formalized as
\lean{Goodman\_\allowbreak bound\_\allowbreak on\_\allowbreak triangle\_\allowbreak density}.

The second result is
\[
  \overline{K_3} + K_3 \;\geq\; \tfrac14\cdot\mathbf{1},
\]
where $\overline{K_3}$ is the empty graph on three vertices: asymptotically,
at least a quarter of all
vertex triples induce a triangle or an independent set.  Reading a graph and
its complement as the two color classes of a $2$-edge-colored complete graph
turns this into Goodman's classical bound on the number of monochromatic
triangles; the proof reuses the downward square certificate of Mantel's
theorem.  It is formalized as
\lean{Goodman\_theorem\_on\_Ramsey\_multiplicity}.

\section{Engineering Obstacles}
\label{sec:obstacles}

The specification layer of \Cref{sec:abstract} turns several steps that
ordinary combinatorial notation leaves implicit into explicit Lean proof
obligations.  These obligations do not change the underlying mathematics, but
discharging them repeatedly requires substantial proof-engineering effort.
This section examines three recurring sources of that effort.  First, because
flags are represented as isomorphism classes, every construction defined on
concrete graphs must be shown to be independent of the chosen representatives.
Second, some graph data, including flag sizes and the vertex types of tuple
entries, is encoded in Lean's types.  Reordering a tuple therefore changes not
only its entries but also the type-level order of their vertex types, and Lean
need not recognize the result as having the same type as a tuple constructed
directly in the new order.  Third, typeclass inference may construct different
terms representing the same finite enumeration, causing otherwise
straightforward rewrites to fail unless the relevant instances are controlled
explicitly.  The first three subsections examine these obstacles in turn.  The
final subsection describes a separate proof pattern that became common in the
specification layer: proving finite counting identities by constructing
explicit bijections.

\subsection{Computing with Graphs up to Isomorphism}

As described in \Cref{sec:flagtypes}, a flag is an isomorphism class of
concrete labeled graphs.  This representation ensures that later statements
do not depend on arbitrary vertex names, but it also imposes a recurring
well-definedness obligation.  A construction defined on concrete graphs gives
a construction on flags only after it has been shown to respect isomorphism.
Thus, for a rule $G\mapsto A(G)$, one must prove
that $G \simeq_f G'$ implies $A(G)=A(G')$.
In Lean, \lean{Quotient.lift} combines the rule with this compatibility proof
to produce a function on flags.  The quotient lift itself is brief; proving
the required compatibility is often the substantial part.

This obligation arises whenever a density or graph operation is first defined
on concrete labeled graphs and then lifted to flags.  It appears again at the
reflection boundary: \lean{Sym2Flag} is a quotient of computational graph
representations, so defining the decoding map
\lean{Sym2Flag.toFlag} of \Cref{sec:reflection-adequacy} requires proving that
equivalent concrete representations decode to the same abstract flag.  Once
this compatibility has been established, \lean{Quotient.lift} yields a
well-defined decoding map that the subsequent adequacy theorems can reuse
without exposing concrete representatives.

\subsection{Comparing Reordered Flag Tuples}
\label{sec:reordered-flag-tuples}

Part of a flag's data lives in its type, so objects that mathematics compares
freely can inhabit different types and fall outside ordinary equality.  What
the remedy costs depends on what the mismatch means.  A flag's size is such an
index: the expansion identity~\eqref{eq:expansion-identity} equates a size-$m$
flag with a combination of size-$n$ flags, and the dependent sum of
\Cref{sec:flagtypes} suffices, since it puts every size in
\lean{FinFlag}~$\sigma$ and the two sides are then distinct elements that the
quotient is there to identify, their difference being a generator of the zero
space.  A reordered tuple is the harder case.  A tuple's type records the
vertex types of its entries, so permuting the entries permutes the type too,
and Lean does not recognize the type so reached as the one it assigns to the
reordered tuple written out directly.  Here the two are meant to be equal
outright, so there is no quotient to appeal to: the types must be proved
equal first and the objects carried across that equality.

The multi-flag density is invariant under reordering its pattern flags.  One
instance is the symmetry
\(
  \den{F_0,F_1}{G}=\den{F_1,F_0}{G}
\),
which is what makes flag-algebra multiplication commutative.  On paper, the
invariance is proved once, in full generality: for any number of pattern
flags, reordering the tuple does not change its density in a host.  The
two-flag symmetry is then read off by taking the permutation that swaps the
two positions, an instantiation a mathematician would not count as a step of
the proof.  The formalization takes the same route, and the friction appears
exactly there: what costs nothing on paper is where the formal proof does its
work.

Recall from \Cref{sec:formal-subflag-densities} that a tuple of pattern
flags is a dependent function: an element of
\lean{FlagList}~$\sigma$~\lean{t Vl} assigns 
a $\sigma$-flag on the vertex type \lean{Vl i} to each index \lean{i : Fin t}, and \lean{flagListDensity}
computes the multi-flag density of such a tuple in a host flag.  Both the
entries and their types vary with the index, and the vertex-type family
\lean{Vl} is itself an argument of the tuple's type.  Reordering is defined
by precomposition:
\begin{lstlisting}
abbrev Perm (t : (*@$\mathbb{N}$@*)) := Fin t (*@$\simeq$@*) Fin t

def listTypePermute {t : (*@$\mathbb{N}$@*)} (Vl : Fin t (*@$\to$@*) Type) ((*@$\pi$@*) : Perm t) : Fin t (*@$\to$@*) Type :=
  fun i => Vl ((*@$\pi$@*) i)

def FlagList.permute {T : Type} {(*@$\sigma$@*) : FlagType T} {t : (*@$\mathbb{N}$@*)} {Vl : Fin t (*@$\to$@*) Type}
    (Fl : FlagList (*@$\sigma$@*) t Vl) ((*@$\pi$@*) : Perm t) : FlagList (*@$\sigma$@*) t (listTypePermute Vl (*@$\pi$@*)) :=
  fun i => Fl ((*@$\pi$@*) i)
\end{lstlisting}
The permuted tuple places \lean{Fl ($\pi$ i)} at position \lean{i}, so the
entry types are reordered too: the vertex-type family of the result is
\lean{listTypePermute Vl $\pi$}, which appears in the type of the result.
The invariance is proved once, for
every length and every permutation; the symmetry to be derived from it is the
pair case:
\begin{lstlisting}
variable {T W U(*@$_0$@*) U(*@$_1$@*) : Type} [Fintype T] [Fintype W] [Fintype U(*@$_0$@*)] [Fintype U(*@$_1$@*)]
         [DecidableEq W] [DecidableEq U(*@$_0$@*)] [DecidableEq U(*@$_1$@*)] {(*@$\sigma$@*) : FlagType T} 
         {t : (*@$\mathbb{N}$@*)} {Vl : Fin t (*@$\to$@*) Type} [FintypeList Vl] [DecidableEqList Vl]

theorem flagDensity_permute (Fl : FlagList (*@$\sigma$@*) t Vl) (G : Flag (*@$\sigma$@*) W) ((*@$\pi$@*) : Perm t)
    : flagListDensity Fl G = flagListDensity (Fl.permute (*@$\pi$@*)) G := ...

theorem flagPairDensity_comm (F(*@$_0$@*) : Flag (*@$\sigma$@*) U(*@$_0$@*)) (F(*@$_1$@*) : Flag (*@$\sigma$@*) U(*@$_1$@*)) (G : Flag (*@$\sigma$@*) W)
    : flagDensity(*@$_2$@*) F(*@$_0$@*) F(*@$_1$@*) G = flagDensity(*@$_2$@*) F(*@$_1$@*) F(*@$_0$@*) G := ...
\end{lstlisting}
To derive the second statement from the first, suppose the $\sigma$-flags
$F_0$ and $F_1$ have vertex types $U_0$ and $U_1$, and let \lean{Vl} be the
vertex-type family with \lean{Vl 0 = U}$_0$ and \lean{Vl 1 = U}$_1$, so that
\lean{[F}$_0$\lean{, F}$_1$\lean{]}$^{\mathtt{f}}$ is an element of
\lean{FlagList}~$\sigma$~\lean{2 Vl}.  Take $\pi$ to swap the two indices.
Then \lean{flagDensity\_permute} equates the density of
\lean{[F}$_0$\lean{, F}$_1$\lean{]}$^{\mathtt{f}}$ with that of its permuted
tuple, whose vertex-type family is \lean{listTypePermute Vl $\pi$}, whereas
\lean{flagPairDensity\_comm} speaks of the directly constructed pair
\lean{[F}$_1$\lean{, F}$_0$\lean{]}$^{\mathtt{f}}$, an element of
\lean{FlagList}~$\sigma$~\lean{2 Vl'} for the family \lean{Vl'} with
\lean{Vl' 0 = U}$_1$ and \lean{Vl' 1 = U}$_0$.  The two families are pointwise
equal, since both send $0$ to $U_1$ and $1$ to $U_0$, but Lean need not reduce
them to the same term, so the two tuples inhabit types that are not
definitionally the same.

The bridging lemma \lean{flagListDensity\_HEq\_eq} closes the gap:
\begin{lstlisting}
theorem flagListDensity_HEq_eq {Vl(*@$_0$@*) Vl(*@$_1$@*) : Fin t (*@$\to$@*) Type}
    [FintypeList Vl(*@$_0$@*)] [DecidableEqList Vl(*@$_0$@*)] [FintypeList Vl(*@$_1$@*)] [DecidableEqList Vl(*@$_1$@*)]
    {Fl(*@$_0$@*) : FlagList (*@$\sigma$@*) t Vl(*@$_0$@*)} {Fl(*@$_1$@*) : FlagList (*@$\sigma$@*) t Vl(*@$_1$@*)}
    (h_Vl_eq : Vl(*@$_1$@*) = Vl(*@$_0$@*)) (h_HEq : HEq Fl(*@$_0$@*) Fl(*@$_1$@*)) (G : Flag (*@$\sigma$@*) W)
    : flagListDensity Fl(*@$_0$@*) G = flagListDensity Fl(*@$_1$@*) G := ...
\end{lstlisting}
Its hypothesis \lean{h\_Vl\_eq} states \lean{Vl' = listTypePermute Vl $\pi$},
proved by checking the two indices of \lean{Fin 2}, while \lean{h\_HEq} uses
Lean's \emph{heterogeneous} equality, which relates terms whose types are not
yet known to be equal; here it says that both tuples place $F_1$ at position
$0$ and $F_0$ at position $1$.  Once the types are identified, that
heterogeneous equality can be used as an ordinary equality, and the lemma
concludes that the densities agree.  The
formal argument thus recovers what is immediate on paper: permuting
\lean{[F}$_0$\lean{, F}$_1$\lean{]}$^{\mathtt{f}}$ produces the same pair as
constructing \lean{[F}$_1$\lean{, F}$_0$\lean{]}$^{\mathtt{f}}$ directly.

\subsection{Controlling Finite Enumerations}
\label{sec:fintype}

Lean distinguishes knowing that a type is finite from having explicit finite
data listing all its elements.  An instance of \lean{Finite T} provides only
the former, whereas an instance of \lean{Fintype T} provides the latter:
\begin{lstlisting}
class Fintype (T : Type*) where
  elems : Finset T
  complete : (*@$\forall$@*) x : T, x (*@$\in$@*) elems
\end{lstlisting}
The field \lean{elems} contains every element of \lean{T}, as guaranteed by
\lean{complete}.  Our specification-layer definitions use this data when
taking cardinalities.  In Lean, the collections they count are typically
subtypes defined by predicates rather than explicit enumerations.  Examples
include tuples of pairwise-disjoint sets of unlabeled vertices in a host flag
that, together with the labeled vertices, induce the flags in a given flag
list, and type embeddings whose images produce a given flag.  Lean can
generally establish \lean{Finite} for these subtypes, and
\lean{Fintype.ofFinite} can then choose a \lean{Fintype} instance
noncomputably.  This is sufficient here because the proofs use the chosen
instance only to state and compare cardinalities, not to evaluate the
enumeration.

\lean{Fintype.ofFinite} is not Lean's only way to construct the required
instance.  At later use sites, operations such as \lean{Fintype.card} leave
their \lean{Fintype} arguments implicit, and typeclass inference may construct
them by another available route.  For example, Lean can construct a
\lean{Fintype} for a predicate-defined subtype by filtering the existing
\lean{Fintype} enumeration of its ambient type.  We call the resulting instance
\emph{structurally derived} because it is assembled from the subtype structure
and the existing instance on the ambient type.  Consequently, a count
definition may contain the instance chosen locally by
\lean{Fintype.ofFinite}, while a cardinality lemma receives this structurally
derived instance.  Both enumerate the same elements and are provably equal,
but if their terms are not definitionally equal, a direct rewrite can fail.

The formalization responds in two ways.  The first is preventive: the moment a
finite set is named, the intended instance is chosen and bound with
\lean{let}, so that every later occurrence elaborates against the same
binding.  The fragment below proves that two finite sets
related by a bijection have equal \lean{toFinset} counts:
\begin{lstlisting}
-- S0 and S1 are the sets that are finite. They are being counted.
-- h_iso_S0_S1 a bijection between them.
let hS0 : Fintype S0 := Fintype.ofFinite S0
let hS1 : Fintype S1 := Fintype.ofFinite S1
have card_eq : @Fintype.card S0 hS0 = @Fintype.card S1 hS1 :=
  @Fintype.card_congr S0 S1 hS0 hS1 h_iso_S0_S1
have count_eq : S0.toFinset.card = S1.toFinset.card := by
  rw [Set.toFinset_card, Set.toFinset_card]
  exact card_eq
\end{lstlisting}
The proof names the \lean{Fintype} instance for each set and uses the bijection
to derive an equality of their cardinalities.  The \lean{@} notation makes the
instance arguments to \lean{Fintype.card} and \lean{Fintype.card\_congr}
explicit.  By contrast, the corresponding arguments of the two
\lean{toFinset} expressions remain implicit in the displayed code; typeclass
inference supplies \lean{hS0} and \lean{hS1} from the preceding bindings.  The
final two lines confirm that these choices agree: \lean{Set.toFinset\_card}
rewrites each count using the instance carried by its \lean{toFinset}
expression, and \lean{exact card\_eq} succeeds because those inferred instances
are definitionally the same as the explicit instances in \lean{card\_eq}.
Without the bindings, separate occurrences may receive different instances,
causing the final step to fail.

The proof of \lean{flagDensity\allowbreak\_permute}, the invariance theorem of
\Cref{sec:reordered-flag-tuples}, uses this combination of an explicit
bijection and fixed \lean{Fintype} instances.  For each concrete representative
of the host flag, \lean{S0} and \lean{S1} consist of tuples of subflags obtained
from pairwise-disjoint choices of unlabeled host vertices.  The tuples in
\lean{S0} realize \lean{Fl}, while those in \lean{S1} realize
\lean{Fl.permute $\pi$}.  Reindexing the tuple entries gives a bijection between
\lean{S0} and \lean{S1}, and the fixed instances allow the resulting
cardinality equality to rewrite the numerator of the density.  A related
bijective count proves that \lean{isomorphismCount}, used in the downward
operator of \Cref{sec:formal-downward}, is invariant under replacing a labeled
graph by an isomorphic representative: the isomorphism transports label
placements and gives a bijection between the corresponding sets.

The second approach is corrective.  After the underlying mathematical objects
have been identified, Lean may still present an equality whose two sides look
identical but contain different hidden instance arguments.  The tactic
\lean{congr!} decomposes such an equality into congruence subgoals and proves
that the instances agree.  This is possible because \lean{Fintype T} is a
subsingleton, so any two \lean{Fintype} instances of the same type are provably
equal; the same holds for \lean{DecidableEq T}.  This is how
\lean{flagListDensity\allowbreak\_HEq\allowbreak\_eq} from
\Cref{sec:reordered-flag-tuples} concludes.  After the equality of the
vertex-type families and the heterogeneous equality of the flag tuples have
been used to identify the inputs, the remaining difference lies only in the
\lean{FintypeList} and \lean{DecidableEqList} instances supplied for their
component types, and \lean{congr!} resolves it.

Explicitly chosen \lean{Fintype} instances are therefore needed only when a
collection specified by a predicate is enumerated as a \lean{Finset}, which
stores an explicit finite enumeration, or when its cardinality is compared
with that of another collection.  Once the required counting identity has
been proved, the subsequent algebraic argument uses that identity without
referring to the chosen instances.

% The discipline above is what the \lean{Fintype} interface requires; Mathlib's
% own documentation says that one generally relies on congruence lemmas when
% rewriting expressions involving \lean{Fintype} instances.  The library also
% provides a cardinality that removes this class of mismatch at its root:
% \lean{Nat.card}, with \lean{Set.ncard} for sets, is a function of the type
% alone, so it carries no instance argument for elaboration paths to disagree
% about, and \lean{Nat.card\_congr} transports a count along a bijection with no
% instance management.  The documentation recommends it precisely when
% finiteness is propositional rather than structural, as it is throughout this
% layer.  In retrospect, stating the specification-layer counts with
% \lean{Nat.card} would have avoided the mismatches this subsection manages,
% leaving a single conversion to \lean{Finset} counts at the reflection
% boundary, where enumeration is genuinely structural.

\subsection{Proving Counting Identities by Bijections}
\label{sec:proof-style}

Explicit bijections became an unexpectedly prevalent proof pattern in the
specification layer.  We expected the counting equalities the layer requires
to follow by rewriting with existing counting lemmas or by symbolic
simplification, but case after case, unfolding the definitions exposed the
same structure: the two sides are cardinalities of finite sets of
configurations presented in different ways.  Our most reliable method is to
identify the two sets and construct an explicit bijection between them, which
Mathlib's cardinality lemmas then convert into the required equality.

Proving the chain rule (\Cref{lem:chain-rule}) made the pattern
clear.  Its two sampling procedures are visibly two descriptions of one
experiment on paper; in Lean, after unfolding the definitions, the proof came
down to a bijection between the finite sets of choices the two procedures
make, which yielded the equality of counts and hence the rule.

The same method proves permutation invariance, invariance under inserting an
empty flag, the downward-operator counting identities, and the adequacy of
the computational subgraph representation.  In each, the bijection carries the
mathematical argument, while the instance discipline of \Cref{sec:fintype}
handles the resulting cardinalities.  We do not claim bijections are the only
approach, but their effectiveness across these obligations made them our
default method.

\section{A Meta-Theory of Ensemble Semantics}
\label{sec:metatheory}

Defining \lean{forbidLE} as we did in \Cref{sec:forbidden}, by constraining
the interpretations rather than the algebra, exposed a mathematical problem in
its own right: when does this \emph{ensemble semantics} validate exactly the
same typed inequalities as a flag algebra in which the constraint is built in
from the beginning?  This section gives the definitions and main results
needed to state the answer.  The proofs are only sketched; a forthcoming paper
will present the full theory and its applications, and the results reported
here are machine-checked in the released development.

\boldparagraph{The constrained positive-homomorphism space.}
Let \(\mathcal G\) be the class of all finite simple graphs, let~\(\mathcal K\subseteq\mathcal G\) be a hereditary class, and fix a flag type
\(\sigma\).  Here \emph{hereditary} means closed under taking induced
subgraphs: if~\(G\in\mathcal K\) and \(S\subseteq V(G)\), then
\(G[S]\in\mathcal K\).  In particular, the \(H\)-free classes used earlier are
hereditary, because taking an induced subgraph cannot create a new
subgraph isomorphic to \(H\).

To make the underlying graph class explicit, we write
\(\FlagAlg{\sigma}[\mathcal G]\) for the ambient
\(\sigma\)-typed flag algebra and
\(\FlagAlg{\sigma}[\mathcal K]\) for the algebra constructed using only graphs
in \(\mathcal K\).  Only \(\sigma\)-flags whose underlying unlabeled graphs
belong to~\(\mathcal K\) occur in this constrained algebra.  Because
\(\mathcal K\) is hereditary, the remaining \(\sigma\)-flags, whose underlying
graphs lie outside \(\mathcal K\), span an ideal,\footnote{An ideal \(I\) here
is a linear subspace of the flag algebra that is also closed under
multiplication by arbitrary algebra elements.  That closure is what makes the
quotient by \(I\) again an algebra, one in which every element of \(I\) has
become zero.} and the usual constrained
algebra is canonically isomorphic to the quotient of the ambient algebra by
this ideal.  We therefore have a canonical quotient map
\[
  q_\sigma:
  \FlagAlg{\sigma}[\mathcal G]
  \twoheadrightarrow
  \FlagAlg{\sigma}[\mathcal K]
\]
that sends each \(\sigma\)-flag outside \(\mathcal K\) to zero and each
remaining flag to its constrained counterpart.

Let
\[
  X_\sigma
  :=
  \mathrm{Hom}^{+}\bigl(\FlagAlg{\sigma}[\mathcal G],\R\bigr)
\]
be the compact space of positive homomorphisms on the ambient
algebra, with the topology of pointwise convergence on \(\sigma\)-flags.
Pullback along \(q_\sigma\) embeds the positive homomorphisms of the constrained
algebra into \(X_\sigma\).  Its image is
\begin{equation}\label{eq:Q-sigma}
  Q_\sigma
  :=
  \left\{
    \psi\circ q_\sigma
    \;\middle|\;
    \psi\in
    \mathrm{Hom}^{+}\bigl(\FlagAlg{\sigma}[\mathcal K],\R\bigr)
  \right\}
  \subseteq X_\sigma.
\end{equation}
Thus, \(Q_\sigma\) is the space of all \(\sigma\)-typed positive homomorphisms
admitted by the constrained quotient.

There is also an intrinsic description:
\begin{equation}\label{eq:Q-sigma-intrinsic}
  Q_\sigma
  =
  \left\{
    \chi\in X_\sigma
    \;\middle|\;
    \chi([F])=0
    \text{ for every \(\sigma\)-flag \(F\) whose underlying graph lies outside
    \(\mathcal K\)}
  \right\}.
\end{equation}
The inclusion of \eqref{eq:Q-sigma} in \eqref{eq:Q-sigma-intrinsic} is
immediate, since \(q_\sigma\) already kills every \(\sigma\)-flag outside
\(\mathcal K\).  For the reverse inclusion, those flags span exactly the ideal
\(q_\sigma\) kills, so a homomorphism vanishing on them factors through
\(q_\sigma\).  The description also exhibits \(Q_\sigma\) as an intersection of
zero sets of continuous evaluation maps, hence as a closed subset of
\(X_\sigma\).  We write
\(Q_0:=Q_{\emptyset}\) for the corresponding space of constrained empty-type
positive homomorphisms.

\boldparagraph{Random extensions and the root-planting space.}
Assume that \(\sigma\) is \emph{non-degenerate} for \(\mathcal K\): that is,
there is some \(\phi_0\in Q_0\) for which
\(\phi_0(\langle\sigma\rangle_0)>0\).  By Razborov's random-extension result
recalled in \Cref{sec:background-downward}, every such \(\phi_0\) determines a
unique probability measure \(\Ext_\sigma(\phi_0)\) on \(X_\sigma\).  For a probability
measure~\(\mu\) on~\(X_\sigma\), its
support \(\operatorname{supp}(\mu)\) consists of those points
\(\chi\in X_\sigma\) for which every open neighborhood of \(\chi\) has positive
\(\mu\)-measure.
Define
\begin{equation}\label{eq:S-sigma}
  S_\sigma
  :=
  \overline{
    \bigcup_{\substack{\phi_0\in Q_0\\
                        \phi_0(\langle\sigma\rangle_0)>0}}
    \operatorname{supp}\bigl(\Ext_\sigma(\phi_0)\bigr)
  }
  \subseteq X_\sigma.
\end{equation}
Thus, to form \(S_\sigma\), we first collect the supports of
\(\Ext_\sigma(\phi_0)\) over all \(\phi_0\in Q_0\) for which \(\sigma\) has
positive density.  We then take the closure in \(X_\sigma\), adding every point
that can be approximated arbitrarily closely by points in those supports.

The support of each random-extension measure is contained in the constrained
space:
\[
  \operatorname{supp}\bigl(\Ext_\sigma(\phi_0)\bigr)
  \subseteq Q_\sigma
  \qquad
  \bigl(\phi_0\in Q_0,
        \phi_0(\langle\sigma\rangle_0)>0\bigr).
\]
Indeed, for a \(\sigma\)-flag \(F\) outside \(\mathcal K\), the downward image
of \([F]\) is a non-negative multiple of \([F|_\emptyset]\), which every
\(\phi_0\in Q_0\) evaluates to zero, so the random-extension identity
\eqref{eq:ensemble} forces \(\mathbb{E}[\phi([F])]=0\) and hence
\(\phi([F])=0\) almost surely.  Ranging over the countably many such flags
puts the sample in \(Q_\sigma\) with probability one, and \(Q_\sigma\) is
closed, so Equation~\eqref{eq:S-sigma} yields the
fundamental inclusion
\begin{equation}\label{eq:S-subset-Q}
  S_\sigma\subseteq Q_\sigma.
\end{equation}
We say that \((\mathcal G,\mathcal K,\sigma)\) is \emph{root-plantable} when
\(S_\sigma=Q_\sigma\).  The name records what the equality says: the labeled
vertices of \(\sigma\) are its \emph{roots}, and to \emph{plant} a
\(\sigma\)-flag's density profile in a larger graph of the class, without
choosing its labels in advance, is to make a fraction bounded away from zero
of all embeddings of \(\sigma\) reproduce that profile.  Root-plantability
asks that every constrained \(\sigma\)-typed positive homomorphism be
approximable by finite constructions of this kind.

\begin{definition}[Built-in and ensemble orders]\label{def:two-constrained-orders}
Let \(\mathcal K\subseteq\mathcal G\) be hereditary, let \(\sigma\) be a flag
type, and let \(f,g\in\FlagAlg{\sigma}[\mathcal G]\).  Define
\[
  f\leq^{\mathrm{quot}}_{\mathcal K,\sigma}g
  \quad\Longleftrightarrow\quad
  \chi(f)\leq\chi(g)
  \quad\text{for every }\chi\in Q_\sigma,
\]
and
\[
\begin{aligned}
  f\leq^{\mathrm{ens}}_{\mathcal K,\sigma}g
  \quad\Longleftrightarrow\quad
  \Ext_\sigma(\phi_0)
  \bigl(\{\chi\in X_\sigma\mid \chi(f)\leq\chi(g)\}\bigr)=1
  \quad\text{for every }\phi_0\in Q_0\text{ with }
  \phi_0(\langle\sigma\rangle_0)>0.
\end{aligned}
\]
\end{definition}

The first is the built-in order of \Cref{rem:builtin-hfree}, written
\(\leq^{\mathrm{quot}}\) because it is the ordinary semantic order in the
constrained quotient algebra: equivalently, every positive homomorphism
\(\psi:\FlagAlg{\sigma}[\mathcal K]\to\R\) satisfies
\(\psi(q_\sigma(f))\leq\psi(q_\sigma(g))\).  The second requires, for each
admissible \(\phi_0\), that the random positive homomorphism
\(\chi\sim\Ext_\sigma(\phi_0)\) satisfy \(\chi(f)\leq\chi(g)\) with
probability one.
For the \(H\)-free class,
\(\lean{forbidLE H f g}\) in our Lean formalization implements this ensemble order.  The corresponding
notions of non-negativity are \(0\leq^{\mathrm{quot}}_{\mathcal K,\sigma}f\) and
\(0\leq^{\mathrm{ens}}_{\mathcal K,\sigma}f\).

\begin{theorem}[Support-closure criterion]\label{thm:support-closure-summary}
Let \(\mathcal K\) be hereditary and let \(\sigma\) be non-degenerate for
\(\mathcal K\).  For every \(f,g\in\FlagAlg{\sigma}[\mathcal G]\),
\[
  f\leq^{\mathrm{quot}}_{\mathcal K,\sigma}g
  \quad\Longrightarrow\quad
  f\leq^{\mathrm{ens}}_{\mathcal K,\sigma}g.
\]
Moreover, the two orders agree for every pair \((f,g)\) if and only if the class
is root-plantable at \(\sigma\):
\[
  \left(
    \forall f,g\in\FlagAlg{\sigma}[\mathcal G],\quad
    f\leq^{\mathrm{quot}}_{\mathcal K,\sigma}g
    \ \Longleftrightarrow\
    f\leq^{\mathrm{ens}}_{\mathcal K,\sigma}g
  \right)
  \quad\Longleftrightarrow\quad
  S_\sigma=Q_\sigma.
\]
\end{theorem}

Read as soundness and completeness, the theorem says that built-in proofs are
always sound for the ensemble semantics, and that under root-plantability they
are complete for it as well.  Both halves come from one observation: each
order asks that the same continuous function \(\chi\mapsto\chi(g-f)\) be
non-negative, and they differ only in the set of \(\chi\) on which it is
tested.  The built-in order tests \(Q_\sigma\) by definition.  The ensemble
order asks for almost-sure non-negativity under each
\(\Ext_\sigma(\phi_0)\), which is non-negativity on that measure's support;
and because the function is continuous, testing it on those supports is the
same as testing it on their closure \(S_\sigma\).  Soundness is therefore
exactly the inclusion \(S_\sigma\subseteq Q_\sigma\) of
Equation~\eqref{eq:S-subset-Q}, and \(S_\sigma=Q_\sigma\) makes the two tests
identical.  When the inclusion is strict, a point of
\(Q_\sigma\setminus S_\sigma\) can be separated from the closed set
\(S_\sigma\), and a flag-algebra element realizing that separation is
non-negative in the ensemble order but not in the built-in one.

\boldparagraph{When root-plantability holds.}
The criterion in \Cref{thm:support-closure-summary} reduces semantic completeness to a structural problem about the
graph class.  Our main positive result gives a broad answer.  Given
\(G\in\mathcal K\), a vertex~\(v\in V(G)\), and a graph \(J\), write
\(G[v\mapsto J]\) for the graph obtained by replacing \(v\) by \(J\), keeping
the edges inside \(J\), and joining every vertex of \(J\) to every former
neighbor of \(v\).  Call \(\mathcal K\) \emph{blow-up-closed} if, for every
\(G\in\mathcal K\), \(v\in V(G)\), and \(N \in \mathbb{N}\), some graph \(J\) on exactly
\(N\) vertices satisfies \(G[v\mapsto J]\in\mathcal K\).  For a hereditary
class, this is equivalent to asking only for \(\lvert J\rvert\geq N\):
choose an induced \(N\)-vertex subgraph \(J'\) of \(J\), and observe that
\(G[v\mapsto J']\) is an induced subgraph of \(G[v\mapsto J]\).  The interior
graph \(J\) may depend on \(G\), \(v\), and \(N\); the definition requires the
existence of a suitable interior, not that every interior works.

\begin{theorem}[Blow-up closure implies root-plantability]
\label{thm:blowup-root-plantability-summary}
Every blow-up-closed hereditary graph class is root-plantable at every
nonempty flag type that is non-degenerate for the class.  The empty flag type is
root-plantable for every hereditary graph class, independently of blow-up
closure.  Consequently, the two orders agree at every non-degenerate flag type
for any blow-up-closed class.
\end{theorem}

We sketch the first claim.  Since \(S_\sigma\subseteq Q_\sigma\) always holds
by Equation~\eqref{eq:S-subset-Q}, root-plantability at \(\sigma\) amounts to
placing every \(\chi\in Q_\sigma\) in \(S_\sigma\), and the gap to close is
between a chosen labeling and a random one: \(\chi\) is a limit of
\(\sigma\)-flags in \(\mathcal K\), whose labelings may be chosen to suit it,
while \(S_\sigma\) asks that a \emph{uniformly chosen} labeling reproduce its
density profile with probability bounded away from zero.  Blow-ups close that gap by making the vertices that carry the profile numerous
enough for a random labeling to find them.  Blow-up closure supplies, inside
\(\mathcal K\), a large interior to replace the labeled vertex, so a uniformly
chosen embedding of \(\sigma\) lands on
one of them with probability bounded away from zero; at the same time, any
fixed number of randomly sampled vertices land in distinct interiors with
probability tending to one, so the blow-up leaves the densities they measure
unchanged.  Iterating this along
a sequence of flags representing \(\chi\) places \(\chi\) in the support of
\(\Ext_\sigma(\phi_0)\) for a limit \(\phi_0\) of the unlabeled blow-ups,
hence in \(S_\sigma\).

Independent blow-ups, in which every replacement graph \(J_v\) is an
independent set, cover all \(K_r\)-free classes
for \(r\geq3\), including triangle-free graphs.  Complete blow-ups, in which
every \(J_v\) is a clique, cover, for example, graphs whose components are
cliques.  Both are instances of
\Cref{thm:blowup-root-plantability-summary}.

\boldparagraph{Why the criterion is nontrivial.}
Heredity alone does not imply root-plantability.  Let \(\mathcal K\) be the
class of \(C_4\)-free graphs and let \(\sigma_1\) be the one-vertex type.
Every large \(C_4\)-free graph has \(o(n^2)\) edges, so a vertex chosen uniformly
at random has normalized degree tending to zero.  
Let \(d:=\fKtwobull\in\FlagAlg{\sigma_1}[\mathcal G]\), so that the density of
\(d\) in a finite \(\sigma_1\)-flag is the normalized degree of its labeled
vertex.  Every \(\chi\in S_{\sigma_1}\) therefore satisfies \(\chi(d)=0\), and
\(-d\) is non-negative in the ensemble semantics.  But arbitrarily large stars
are \(C_4\)-free: the star \(K_{1,t}\) consists of one vertex, its
\emph{center}, adjacent to \(t\) pairwise non-adjacent leaves.  Label the
center to obtain a \(\sigma_1\)-flag, and let \(\chi_*\in Q_{\sigma_1}\) be the
limit of that sequence.  Every leaf is
adjacent to the center, so \(\chi_*(d)=1\) and
\(\chi_*\notin S_{\sigma_1}\), hence \(S_{\sigma_1}\subsetneq Q_{\sigma_1}\).
Moreover, \(\chi_*(-d)=-1<0\), so \(-d\) fails to be non-negative in the
built-in semantics.
The example exhibits a genuine gap between the two typed orders: built-in
semantics permits a specially chosen vertex to carry the label, whereas
ensemble semantics labels a uniformly sampled vertex.

The gap does not invalidate the final unlabeled density bounds of this
paper.  The
certificate pipeline consumes its typed inequalities only through the
downward operator, which preserves non-negativity in the ensemble
semantics
(\lean{downward\_\allowbreak preserve\_\allowbreak semanticCone},
\Cref{sec:formal-downward}).  Its conclusions land at the empty flag type,
which is root-plantable for every hereditary class by
\Cref{thm:blowup-root-plantability-summary}.  The final conversion
(\lean{generalizedTuranDensity\_\allowbreak le\_\allowbreak of\_\allowbreak forbidLE},
\Cref{sec:forbidden}) then turns these conclusions into classical
statements about Tur\'an densities.  Root-plantability is therefore not a
hypothesis of any of these results.  It concerns only whether the intermediate
ensemble calculus proves exactly the constrained algebra's typed inequalities.

\boldparagraph{Formalization and scope.}
The Lean formalization of this section is included in the public release,
under \lean{LeanFlagAlgebras.MetaTheory}.  There
\(Q_\sigma\), \(S_\sigma\), and the condition \(S_\sigma=Q_\sigma\) are
\lean{Q$\sigma$}, \lean{S$\sigma$}, and \lean{RootPlantable}; the two orders of
\Cref{def:two-constrained-orders} appear as the non-negativity predicates
\lean{QuotientNonneg} and \lean{EnsembleNonneg}, with \(f\leq g\) expressed by
applying them to \(g-f\).  The theorems are
\lean{quotient\_implies\_ensemble} and \lean{support\_criterion} for the two
halves of \Cref{thm:support-closure-summary},
\lean{blowupClosed\_root\_plantable} and
\lean{heredClass\_emptyType\_rootPlantable} for the nonempty- and empty-type
claims of \Cref{thm:blowup-root-plantability-summary}, and
\lean{c4free\_not\_rootPlantable} with \lean{S$\sigma$\_subset\_Q$\sigma$} for
the strict inclusion \(S_{\sigma_1}\subsetneq Q_{\sigma_1}\).

These files were generated by Claude Code from an already-developed
mathematical account and its theorem targets, and Lean checks the resulting
proof terms.  The autoformalization was feasible because it extended rather
than rebuilt the manual development of the preceding sections, which had
already fixed the representations and semantic interfaces and supplied the
ambient algebra, the positive-homomorphism spaces, the random extensions, and
the downward operator; only the constrained layer, the forbidden ideal and its
quotient, had to be built.

This section isolates the core consequences of \lean{forbidLE}: built-in
inequalities always imply ensemble ones, root-plantability characterizes when
the converse holds, and blow-up closure is a broad sufficient condition for
it.  The forthcoming paper gives the complete account and goes further, with a
finite criterion implying root-plantability, further classes on either side of
the boundary, and a relative ensemble semantics that admits non-hereditary
constraints such as a fixed edge density.

\section{Open Design Questions}
\label{sec:unsure}

The formalization and certificate proofs presented above are fully
checked in Lean.  This section instead concerns design questions that the
present implementation cannot settle on its own.  For each question, we state
what our implementation establishes and what remains unknown because answering
it would require either building a competing implementation or running
substantially larger examples.

\boldparagraph{When should a forbidden-graph constraint enter the formalization?}
In our ensemble semantics (\Cref{sec:forbidden}), we defer the constraint until
elements are compared.  Consequently, the seven certificates of
\Cref{sec:results}, involving four different forbidden graphs, share a single
ambient algebra and a common library of algebraic lemmas.  This choice does not
compromise soundness: the ensemble semantics yields the asymptotic bounds
directly.  Comparing it with Razborov's standard, built-in semantics in
\Cref{sec:metatheory} is a separate mathematical question, not a
soundness obligation.

This does not determine which design is more economical overall.  Building the
constraint in from the outset may streamline arguments for a fixed extremal
problem, but we have not formalized comparable certificate proofs using that
design.  Even when root-plantability ensures that the two semantics validate
the same inequalities, their implementations may differ in engineering and
computational cost.

\boldparagraph{Should specification and computation use separate representations?}
Using separate representations for specification and computation kept the
abstract development close to standard mathematics while making finite
calculations executable.  In all seven case studies, the required density facts
were obtained by executable evaluation and transferred to the specification
layer by adequacy theorems; the multiplication facts were then derived from
these density facts.  Maintaining this bridge required substantial work.  An
alternative would be to formulate the theory directly using a concrete,
canonical representation.  This could reduce the bridging work, but choices
about representation and canonicalization would then appear in the statements
and proofs of the mathematical theory.  Determining which approach is less
costly would require a second development of comparable scope.  Nevertheless,
our separation follows the established methodology of Cohen et
al.~\cite{cohen2013refinements}.  It also appears independently in the
concurrent formalization of local flag algebras~\cite{davey2026local}, a
calculus normalized by maximum degree that takes no quotient and likewise
separates a noncomputable specification from a computable representation.

\boldparagraph{Can the computational layer scale beyond five vertices?}
Our seven certificates (\Cref{sec:results}) use host graphs on three,
four, or five vertices.  At these sizes, the optimized enumeration of
\Cref{sec:reflection-generation} keeps the required families of finite
identities manageable.  The Erd\H{o}s pentagon certificate, for example,
generates 2,832 pair-density theorems across its three flag types, all checked
by \lean{decide +kernel}.  Thus, the current reflection and automation layers
can handle certificates with five-vertex hosts.  Fully kernel-checked
compilation, however, is already near its practical limit: the Erd\H{o}s
pentagon certificate is the largest example completed entirely on the
kernel-checked path.

Going substantially beyond five vertices remains a major engineering
challenge.  As the host size grows, the number of flags and generated
pair-density and multiplication identities grows, as does the size of the
expressions that the final algebraic proof must normalize.  The other two
five-vertex certificates illustrate this growth: the $K_5$-free and
$C_5$-free edge-density cases require 10,332 and 8,172 pair-density theorems,
respectively, and therefore use \lean{native\_decide}.  Scaling further will
likely require more efficient executable graph representations and counting
algorithms in the reflection layer, as well as improvements to theorem
generation and algebraic normalization.  The same combinatorial growth also
constrains certificate generation: the Flagmatic user guide documents support
for bases of up to eight vertices for ordinary graphs and seven for
3-graphs~\cite{vaughan2012flagmaticguide}.  Our formal pipeline currently
reaches a practical limit earlier, and we have not yet determined which stage
is the dominant bottleneck or how much these optimizations could narrow the
gap.

\section{Related Work}
\label{sec:related}

\boldparagraph{Certificate checking and proof reconstruction.}
Checking the output of an untrusted solver without adding the solver to the
trusted base is a well-established theme in verification.  Examples include
clausal proofs for SAT~\cite{lammich2020unsat,tan2023cakelpr}, proof
reconstruction and verified checkers for
SMT~\cite{bohme2010z3,armand2011modular}, and sums-of-squares certificates for
polynomial
inequalities~\cite{harrison2007verifying,besson2006fast,morrison2026sos}.
These systems follow the certifying-algorithm paradigm: an untrusted procedure
performs the expensive search and returns evidence that a simpler checker
validates~\cite{mcconnell2011certifying}; such checkers can themselves be
formally verified~\cite{alkassar2014framework}.  Translation validation applies
the same per-run principle to compilers and other
translators~\cite{pnueli1998translation}.  Roux et al.\ apply the pattern to
numerical SDP-solver output for polynomial program
invariants~\cite{roux2018validating}.

Our approach is closest to Morrison's \lean{sos}
tactic~\cite{morrison2026sos}, which uses an external semidefinite solver to
find a floating-point Gram matrix, rounds it to rational data, and verifies the
resulting sums-of-squares decomposition through rational $LDL^\top$.  Our
certificate-checking tactic, \lean{flag\_certificate}, plays a similar role: it
treats the certificate as untrusted hint data and, for each sums-of-squares
block, constructs an exact rational factorization
$M=L\operatorname{diag}(d)L^\top$
and asks Lean to check both the matrix identity and the non-negativity of every
entry of $d$.  These checks establish positive semidefiniteness without
floating-point arithmetic.  The main difference is that \lean{sos} does not
have to construct and verify a separate family of combinatorial objects before
checking its certificate.  A flag-algebra certificate, by contrast, refers to
types, flags, and host graphs considered up to isomorphism.
Certificate-specific elaboration commands therefore generate canonical
representatives, prove the enumerations exhaustive, and establish the required
density and multiplication identities before \lean{flag\_certificate}
interprets and checks the certificate.  This verification also involves
normalization in the quotient algebra and forbidden-subgraph reasoning.

\boldparagraph{Proof by reflection and trusted evaluation.}
Proof by reflection has a long history in Rocq, from
Boutin~\cite{boutin1997reflection} to the extended treatment by
Chlipala~\cite{chlipala2013cpdt}.  Cohen et
al.~\cite{cohen2013refinements} develop a refinement methodology for relating
representations suited to mathematical reasoning to more efficient
representations suited to computation.  Our \lean{Sym2Graph}-based reflection
layer follows this general pattern, with adequacy theorems connecting its
computable representations to their specification-level counterparts.  The
trust boundary of reflected computation depends on how that computation is
evaluated.  Rocq's \lean{vm\_compute} uses a bytecode-based virtual machine
within its trusted reduction machinery~\cite{gregoire2002compiled}, whereas
\lean{native\_compute} compiles the computation to OCaml and therefore also
relies on the OCaml compiler and runtime~\cite{boespflug2011full}.  Lean offers
an analogous choice: \lean{decide +kernel} evaluates the decision procedure by
kernel reduction, whereas \lean{native\_decide} uses Lean's native compiler and
extends the trusted base accordingly.  Our development uses both mechanisms,
and \Cref{sec:flagmatic-trust} identifies which one is used for each certificate
obligation.

\boldparagraph{Proof engineering at scale.}
Sustaining large mechanized developments raises engineering questions that
are increasingly studied in their own right, from the survey of Ringer et
al.~\cite{ringer2019qed} to the retrospective on Project
Everest~\cite{ahman2026everest}.  That literature focuses largely on program
verification, whereas the obstacles we report (\Cref{sec:obstacles}) arise
from formalizing computational mathematics.  In our development, the main
sources of proof-engineering effort were proving constructions well defined on
graphs represented up to isomorphism, reconciling related objects whose
dependent types Lean does not recognize as identical, and keeping inferred
finite enumerations consistent.

\boldparagraph{Formalizations of combinatorics.}
Substantial combinatorial arguments have already been formalized in Lean.
Dahmen, H\"olzl, and Lewis formalized the Ellenberg--Gijswijt solution of
the cap set problem~\cite{dahmen2019capset}, Dillies and Mehta
Szemer\'edi's regularity lemma and the qualitative form of Roth's
theorem~\cite{dillies2022szemeredi}, and Mehta the Kruskal--Katona
theorem~\cite{mehta2022kruskal}.  Subercaseaux et
al.~\cite{subercaseaux2024hexagon} verified the empty-hexagon number against
a SAT certificate.  The collaborative PFR project~\cite{pfr2023} formalized
the Gowers--Green--Manners--Tao proof~\cite{gowers2025marton} of the
polynomial Freiman--Ruzsa conjecture in characteristic two.  Those
formalizations are each organized around a headline theorem, contributing
reusable infrastructure along the way. Ours is aimed at the method itself: it
formalizes the flag algebra, its semantics, and the finite machinery its
certificates need.
\Cref{sec:lowerbounds} shows the resulting library used directly, for
arguments the compiler does not cover.

\boldparagraph{Computer-assisted flag algebra arguments.}
Flagmatic~\cite{vaughan2013flagmatic} automates the computation of exact and
approximate flag-algebra bounds but does not produce proof-assistant proof
terms; Razborov's survey~\cite{razborov2013flag} organizes many early
applications of the method.  The seven certificates compiled here are
Flagmatic's output as produced, and our compiler consumes them as untrusted
input.

\boldparagraph{Concurrent flag-algebra work in Lean.}
Two Lean developments run concurrently with, and independently of, ours: one
formalizes a different flag-algebra calculus, the other aims to replace the
external solver rather than to validate its output.
Davey, Hurley, de Joannis de Verclos, Kang, and Volec formalized \emph{local
flag algebras}, the calculus introduced in their accompanying
paper~\cite{davey2026local} and applied in its companion~\cite{davey2026strong}.
It normalizes densities by the maximum degree $\Delta$
rather than by the number of vertices, so that the method still yields
nontrivial bounds for sparse graphs of small maximum degree, and under this
normalization no quotient is taken: their algebra is defined directly on
formal combinations of isomorphism classes.  Their formalization
machine-certifies the new bounds of both papers, and independently arrives at
the reflection pattern of \Cref{sec:reflection}: a
noncomputable specification mirrored by a computable representation (Boolean
adjacency in their case) with adequacy lemmas.  Our development formalizes the classical
theory instead (the quotient algebra, positive homomorphisms, and the ensemble
semantics) and is organized as a reusable library with a certificate-to-proof
compiler on top, so that certificates produced outside Lean become Lean
proofs.
The other is Spiegel's \emph{computational flag
algebras} in Lean~\cite{spiegel2026computational}, ongoing independent work
aimed at carrying out flag-algebra inequality proofs fully and automatically
inside Lean, without an external tool or solver.

\boldparagraph{Graph limits.}
Graphon theory provides an alternative analytic framework for the asymptotic
graph-density phenomena studied by flag algebras; see Lov\'asz and Szegedy and
Lov\'asz~\cite{lovasz2006limits,lovasz2012large}.  Freer has formalized a
substantial graphon theory in Lean~4, covering cut distance, weak regularity,
counting lemmas, and compactness~\cite{freer2026graphon}.  Our
development formalizes the flag-algebra side of this relationship rather than
graphons.  Formally connecting our positive-homomorphism semantics to Freer's
graphon formalization remains future work.

\section{Conclusion}
\label{sec:conclusion}

Flag-algebra certificates are produced by unverified numerical toolchains and
are far too large to audit by hand.  We have built a Lean~4
certificate-to-proof pipeline that consumes such certificates without trusting
them.  Certificate-specific elaboration commands first generate the required
density and multiplication identities; the certificate-checking tactic then
reads the certificate, checks its remaining claims, and assembles a complete
proof of a theorem stated in the source file.  Underlying the pipeline are a
specification layer that follows the standard mathematics and a reflection
layer that makes the finite parts of that mathematics executable, joined by
adequacy theorems.  We compiled seven Tur\'an-type upper bounds, and paired two
of them, Mantel's theorem and the Erd\H{o}s pentagon theorem, with lower-bound
constructions proved directly against the same formal interface.

One design choice shaped the development throughout.  In the standard
semantics, the forbidden-subgraph constraint is built into a separate algebra
for each forbidden family.  We instead kept one ambient algebra and imposed the
constraint only when its elements are compared.  This ensemble semantics keeps
the core definitions independent of the extremal problem being solved, and it
yields the intended asymptotic bounds on its own terms.  Comparing the two
semantics led to mathematics of its own: validity in the standard semantics
always implies validity in the ensemble semantics, and the converse holds for
all expressions exactly when the class is root-plantable.

We do not claim that the approach presented here is the right way to formalize
flag algebras.  Our aim
has been to make the design choices and their consequences explicit: the
engineering obligations they create, the layers needed to discharge them, and
the questions left open.  Some choices reflect current feasibility constraints
rather than principled preferences; most visibly, the two largest certificates
still rely on \lean{native\_decide}, so their proofs also trust Lean's native
compiler.

Several directions remain for future work.  One is to extend our Lean
development to additional parts of Razborov's flag-algebra calculus, especially
the differential structure developed in Section~4.3 of his original
paper~\cite{razborov2007flag}; this would exercise the quotient-level semantics
well beyond the certificate arguments considered here.  Another is to adapt the
formalization to richer combinatorial structures such as hypergraphs and
directed graphs.  A third is to build a formal bridge between our
positive-homomorphism semantics and graphon theory, connecting our development
to Lov\'asz's graph-limit theory~\cite{lovasz2012large} and the existing Lean
formalization of graphons~\cite{freer2026graphon}.

\subsection*{Use of AI Assistance}
The layers of this development used AI assistance to different degrees.  The
specification and reflection layers of \Cref{sec:abstract,sec:reflection},
including the quotient constructions, their interface theorems, and the
adequacy proofs, were formalized manually by the authors.  In the
certificate-to-proof compiler of \Cref{sec:flagmatic}, most of the required
metaprogramming, including the elaboration-time generation commands, the custom
tactics, and the certificate compiler, was developed with the help of
Claude Code.  Finally, the mathematical results of the meta-theory
summarized in \Cref{sec:metatheory} were obtained with the help of GPT~Pro,
and their Lean development was autoformalized entirely by Claude Code.  We
also used Codex and Claude Code to assist with writing and revising this paper.

\begin{acks}
We would like to thank Jineon Baek, Taeyoung Kim, Hyunwoo Lee, Joonkyung Lee, Christoph Spiegel, 
and Jan Volec for helpful discussions on Razborov's flag algebras and their Lean formalization.  
We are also grateful to
Ross Kang and Sidharth Hariharan for encouraging us to work harder to remove
\lean{native\_decide} in our Lean formalization.
This work was supported by the National Research Foundation of Korea (NRF)
grant funded by the Korean Government (MSIT) (No.~RS-2023-00279680) and by the
Institute for Basic Science (IBS-R029-C1).
\end{acks}

\bibliographystyle{ACM-Reference-Format}
\bibliography{refs}

\appendix
\section{Compilation Times}
\label{sec:compile-times}

Table~\ref{tab:compile-times} reports end-to-end compilation time and peak
memory for all seven compiler examples of \Cref{sec:results}, compiled with
Lean~4 (v4.27.0) on an
Intel Core i7-14700K (20 cores) with 64\,GB of RAM under Windows~11.  Each
measurement covers the complete source file: flag enumeration, pair-density and
multiplication lemma generation, PSD certificate checking, the objective
expansion, and the main theorem.  The generated
files differ only in how their bridging lemmas are discharged.  Under \lean{decide +kernel} (the
\lean{flagGen.kernelDecide} option), the file introduces no compiled-evaluation
axiom; under \lean{native\_decide}, it is faster but additionally depends on the
\lean{Lean.ofReduceBool} and \lean{Lean.trustCompiler} axioms.  We report both
modes for the five examples where kernel evaluation is tractable; for the two
largest, the $K_5$- and $C_5$-free edge-density examples, only
\lean{native\_decide} was measured.

Each figure is the mean over five successful runs of a single-file compilation
against the prebuilt library, with the sample standard deviation; memory is the
peak resident set, taken as the maximum over the five runs.  One run stalled on
system contention, showing an anomalous wall clock at an unchanged peak memory;
it was discarded and replaced by a further run.  Wall-clock times are highly
repeatable (standard deviation below $1\%$ of the mean in every configuration
except the $P_3$- and $C_4$-density cases under \lean{native\_decide}, at
$2.3\%$ and $1.8\%$).

\begin{table}[h]
  \centering
  \caption{End-to-end Lean compilation time (mean of five runs, $\pm$ sample
    standard deviation) and peak resident memory (maximum over five runs) for
    all seven examples, under \lean{decide +kernel} and \lean{native\_decide}.
    ``Lemmas'' is the total count of generated \lean{flagDensity}$_2$
    lemmas.  A dash marks a configuration not measured under kernel evaluation.}
  \label{tab:compile-times}
  \small
  \setlength{\tabcolsep}{4.5pt}
  \begin{tabular}{lccr rr rr}
    \toprule
    & & & & \multicolumn{2}{c}{\lean{decide +kernel}}
        & \multicolumn{2}{c}{\lean{native\_decide}} \\
    \cmidrule(lr){5-6}\cmidrule(lr){7-8}
    Case & Forbid & $N$ & Lemmas
      & Time (s) & RAM (GB) & Time (s) & RAM (GB) \\
    \midrule
    Mantel (edge density)              & $K_3$ & 3 &    15 & $15.6\pm0.1$   &  2.4 & $13.4\pm0.1$  & 2.0 \\
    $P_3$ density                      & $K_3$ & 3 &    15 & $14.9\pm0.0$   &  2.4 & $12.8\pm0.3$  & 1.9 \\
    $C_4$ density                      & $K_3$ & 4 &   210 & $76.6\pm0.7$   &  8.1 & $21.1\pm0.4$  & 2.1 \\
    Edge density                       & $K_4$ & 4 &   390 & $135.2\pm0.4$  & 11.2 & $29.0\pm0.0$  & 2.4 \\
    Erd\H{o}s pentagon ($C_5$ density) & $K_3$ & 5 &  2832 & $2713\pm13$    & 42.3 & $122.6\pm0.8$ & 4.1 \\
    Edge density                       & $K_5$ & 5 & 10332 & ---            & ---  & $325.0\pm0.6$ & 8.0 \\
    Edge density                       & $C_5$ & 5 &  8172 & ---            & ---  & $531.9\pm3.0$ & 8.5 \\
    \bottomrule
  \end{tabular}
\end{table}

\end{document}
